\documentclass[11pt]{article}
\usepackage[T1]{fontenc}
\usepackage[utf8]{inputenc}
\usepackage{amsmath,amssymb,amsthm,mathtools}
\usepackage{geometry,graphicx,booktabs,longtable,tabularx,array}
\usepackage{enumitem,needspace,xcolor,tikz,tcolorbox}
\usetikzlibrary{arrows.meta,positioning,calc,fit,backgrounds,decorations.pathreplacing,matrix,patterns}
\tcbuselibrary{breakable}
\usepackage{caption,hyperref,bookmark,cleveref,microtype}
\usepackage{komorebi-publication}
\renewcommand{\KomorebiTitle}{Komorebi Space}
\renewcommand{\KomorebiSubtitle}{Coherent realizations, obstructions, and descent}
\renewcommand{\KomorebiSeries}{Foundations of the Realization Calculus}
\renewcommand{\KomorebiKind}{Monograph}
\renewcommand{\KomorebiEdition}{Proof-expanded monograph edition 1.7.0 / 12 October 2026}
\renewcommand{\KomorebiShortTitle}{Komorebi Space}
\renewcommand{\KomorebiProfile}{monograph}
\KomorebiSetup
\setcounter{tocdepth}{1}
\setcounter{secnumdepth}{2}
\setlist{nosep,leftmargin=1.7em,topsep=4pt}
\hypersetup{pdfkeywords={Komorebi Space,realization calculus,descent,obstruction,Lean,holonomy,closed range}}
\newtheorem{theorem}{Theorem}[section]
\newtheorem{proposition}[theorem]{Proposition}
\newtheorem{lemma}[theorem]{Lemma}
\newtheorem{corollary}[theorem]{Corollary}
\theoremstyle{definition}
\newtheorem{definition}[theorem]{Definition}
\newtheorem{example}[theorem]{Example}
\theoremstyle{remark}
\newtheorem{remark}[theorem]{Remark}
\numberwithin{equation}{section}
\newcommand{\R}{\mathbb R}
\newcommand{\Z}{\mathbb Z}
\newcommand{\Q}{\mathbb Q}
\newcommand{\N}{\mathbb N}
\newcommand{\K}{\mathbb K}
\newcommand{\V}{\mathcal V}
\newcommand{\W}{\mathcal W}
\newcommand{\id}{\operatorname{id}}
\newcommand{\im}{\operatorname{im}}
\newcommand{\ran}{\operatorname{ran}}
\newcommand{\Hom}{\operatorname{Hom}}
\newcommand{\Ob}{\operatorname{Ob}}
\newcommand{\Ext}{\operatorname{Ext}}
\newcommand{\Match}{\operatorname{Match}}
\newcommand{\Exact}{\operatorname{Exact}}
\newcommand{\dist}{\operatorname{dist}}
\newcommand{\Aut}{\operatorname{Aut}}
\newcommand{\Stab}{\operatorname{Stab}}
\newcommand{\Set}{\mathbf{Set}}
\newcommand{\Cat}{\mathbf{Cat}}
\newcommand{\norm}[1]{\left\lVert#1\right\rVert}
\newcommand{\abs}[1]{\left\lvert#1\right\rvert}
\newcommand{\source}[1]{\par\nopagebreak[4]\phantomsection\label{src:#1}{\footnotesize\sffamily\color{KomorebiMuted}Source concordance: \hyperlink{claim:#1}{#1}.}\par\smallskip}
\newcommand{\code}[1]{\texttt{\detokenize{#1}}}
\newcommand{\chap}[1]{\par\Needspace{10\baselineskip}\section{#1}}
\newcommand{\partnote}[2]{\clearpage\part{#1}\noindent\begin{minipage}{.92\linewidth}\large #2\end{minipage}\par\vspace{12pt}}
\newtcolorbox{insight}{breakable,colback=KomorebiSoft!55,colframe=KomorebiPine!50,boxrule=.45pt,arc=1pt,left=9pt,right=9pt,top=7pt,bottom=7pt}
\newtcolorbox{scopebox}{breakable,colback=KomorebiWarm,colframe=KomorebiOchre!55,boxrule=.4pt,arc=1pt,left=9pt,right=9pt,top=7pt,bottom=7pt}
\usepackage{float}
\newcommand{\fixedplate}[3]{\begin{figure}[H]\centering\resizebox{.96\linewidth}{!}{\input{figures/#1}}\caption{#2}\label{#3}\end{figure}}
\newcommand{\plate}[3]{\begin{figure}[!htbp]\centering\resizebox{.96\linewidth}{!}{\input{figures/#1}}\caption{#2}\label{#3}\end{figure}}
\tikzset{box/.style={draw=KomorebiPine!70,rounded corners=2pt,fill=KomorebiSoft!50,align=center,inner sep=8pt,font=\small},arr/.style={-{Latex[length=2.2mm]},draw=KomorebiPine,thick},thinarr/.style={-{Latex[length=1.8mm]},draw=KomorebiMuted},lab/.style={font=\small,align=center},dot/.style={circle,fill=KomorebiPine,inner sep=2.3pt}}
\renewcommand{\arraystretch}{1.18}
\raggedbottom
\widowpenalty=7000\clubpenalty=7000
\setlength{\emergencystretch}{3em}

% Compact only the contents entries, without reducing body type size.
\usepackage{etoolbox}
\makeatletter
\patchcmd{\l@section}{1.0em}{0.6em}{}{}
\makeatother

% Generated by checks/edition_stats.py from the current executed reports.
\input{edition_stats}

% Keep readable half-page figures with their surrounding exposition.
\renewcommand{\topfraction}{.88}
\renewcommand{\bottomfraction}{.80}
\renewcommand{\textfraction}{.10}
\renewcommand{\floatpagefraction}{.78}
\setcounter{topnumber}{2}
\setcounter{bottomnumber}{2}
\setcounter{totalnumber}{3}

% WR evidence is generated from separately scoped accepted extension records.
\input{written_status}
\newcommand{\written}[1]{\par\nopagebreak[4]\smallskip{\footnotesize\sffamily\color{KomorebiMuted}Extension \hyperref[app:written]{#1}. \csname WRstatus#1\endcsname}\par\smallskip}

\begin{document}
\KomorebiCover
\pagenumbering{roman}
\section*{About this edition}
This monograph develops Komorebi Space as a theory of coherent realizations relative to a specified system of mathematical views. Its central constructions are ordinary limits, relative extension spaces, obstruction quotients, cocycle groupoids, an explicit stack on a composition site, and quantitative repair operators. The point of the synthesis is to keep these constructions connected without erasing the differences between their hypotheses.

The exposition follows the mathematics rather than the order in which the source library grew. Definitions precede existence questions. Existence is separated from uniqueness. Algebraic feasibility is separated from numerical stability. Changing a representation is separated from changing the data that the representation is required to preserve.

Publication version \textbf{1.7.0} deepens the existing proofs rather than enlarging the theorem program. Sixty-seven proof blocks have been rewritten to display the intermediate constructions, inverse equations, quantifier choices, and limiting arguments. The changes cover fifty KS results, all sixteen WR groups, and the weak-subsequence lemma. A new analytic appendix proves the projection, adjoint, Banach inverse, and Neumann-series tools used in those arguments.

The eighty source-linked KS entries retain the accepted version 1.2.0 canonical baseline, with \EditionFormalModules\ production modules and one root. The separate regularization extension retains its accepted version 1.6.0 focused replay. All files in both formal trees are unchanged. \emph{No new Lean compilation or axiom audit is claimed in this proof-exposition edition.} Fresh checks concern the source preservation, proof-edit inventory, exact finite suites, and publication build.

The WR catalogue still has sixteen groups. WR06, WR07, WR13, and WR14 have complete focused typed counterparts. WR09 and WR16 have formal pointwise estimates but written limiting consequences. The other ten groups remain written-only. Longer printed proofs do not change these evidence labels. The new background appendix is written mathematical support, not an additional formal release.

Lean's compiler, kernel, and pinned external dependency cache remain trusted. This is not an independent rebuild of mathlib or a proof that a Python program implements every formal construction. Source correspondence identifies the actual declarations and hypotheses. It does not turn explanatory prose into a kernel-checked language.

All mathematical figures are reproducible vector diagrams or Python plots of explicit formulas. Axes, support patterns, and worked values are checked against their stated mathematical models. The source package contains the manuscript, figure sources, exact checks, evidence archives, and the progress ledger. The separate visual guide offers a shorter route through the examples. The Lean guide explains the correspondence and its verification boundaries.

\begin{scopebox}
\textbf{Reading scope.} The word \emph{space} does not assert one universal topological object containing all of mathematics. The strict construction is a set of coherent families. Symmetry is retained in a groupoid, and effective descent is proved for a specified site. Analytic estimates apply only after a normed or Hilbert model has been declared. These are related layers, not interchangeable definitions.
\end{scopebox}

\clearpage
\section*{Orientation}
A description at one scale often omits information visible at another. A formula in one representation may fail to commute with a change of representation. Local pieces may each be valid and still fail to assemble. Even when assembly is possible, there may be many completions, and small discrepancies may or may not imply proximity to an exact answer.

Komorebi Space treats these questions in a common order. First specify the views and the maps between them. Then define the coherent realizations. Next identify the information that must remain fixed. Finally determine the obstruction to extension, the freedom in a successful extension, and the conditions under which repair is stable.

The resulting program has three central distinctions:
\[
\begin{gathered}
\text{a collection of appearances}\quad\ne\quad\text{a coherent realization},\\
\text{existence of an exact answer}\quad\ne\quad\text{a stable repair estimate},\\
\text{one isomorphism class}\quad\ne\quad\text{absence of symmetries}.
\end{gathered}
\]
The monograph makes these distinctions into explicit constructions and theorems.

\subsection*{The route through the book}
Part I constructs coherent sections and follows their obstructions as more data are prescribed. Part II develops resolution towers, finite observations, split repair coordinates, and multiscale reconstruction. Part III retains symmetries, proves effective descent, and distinguishes coherent phases from removable ones. Part IV develops analytic stability and an independent geometry of transformation costs. Part V follows approximate answers to their limits, distinguishes a coherent point from its intended Hilbert realization, and then derives the bounded-budget and regularization criteria that diagnose the difference. Part VI works through two finite loop models and returns to mathematical semantics. Appendix~\ref{app:workshop} contains fifteen solved problems. Appendix~\ref{app:hypotheses} gives a hypothesis atlas, and Appendix~\ref{app:written} records the exact evidence scope of the extensions. Appendix~\ref{app:proof-tools} proves the analytic background used in their arguments.

A short technical route is Sections~\ref{sec:views}, \ref{sec:obstruction}, \ref{sec:stack}, \ref{sec:closed}, and~\ref{sec:motion}. A reader interested in the role of packets and residuals should include Sections~\ref{sec:split} and~\ref{sec:cascade}. The graph models in Section~\ref{sec:graph} explain how a categorical realization problem acquires an actual numerical condition number. Sections~\ref{sec:observation} and~\ref{sec:repair-transport} give the representation-theoretic links. Sections~\ref{sec:signal-loop}, \ref{sec:twisted-loop}, and~\ref{sec:escape} can also be read first as concrete motivations. The last of these connects finite realization, closure, and failure of protected infinite existence in one calculation. Sections~\ref{sec:bounded-realization} through~\ref{sec:diagonal-regularization} give its positive converse, fixed-budget diagnostic, and regularized continuation. Sections~\ref{sec:inexact-repair} and~\ref{sec:stopping-rules} then explain how to certify an inexact calculation and choose its stopping tolerance.

\subsection*{What is being contributed}
The limit universal property, gauge action groupoids, standard stack definitions, orthogonal projection, and bounded inverse principles belong to established mathematics. They are used here as such. The contribution of this presentation is their explicit integration with the project's actual view systems, protected correction spaces, resolution differentials, graph transports, and checked examples. Historical novelty of every component is not inferred from formal verification. Conversely, the fact that the ingredients are classical does not remove the need to prove the comparison maps and to retain the right hypotheses.

\clearpage\begingroup\setlength{\parskip}{0pt}\tableofcontents\endgroup
\clearpage\pagenumbering{arabic}
\partnote{Realizations and their obstructions}{A global object is not added to local data by terminology. It is a family satisfying every prescribed transition equation. The first task is to construct the space of such families and to keep its possible emptiness visible.}
\section{Views and coherent realizations}\label{sec:views}
\subsection{The data of a view system}
Fix a small category \(\V\), or work in fixed universes so that the products below exist in the chosen ambient universe. A \emph{view} is an object of \(\V\). It can encode a resolution level, a coordinate presentation, or a context in which part of an object is visible. These interpretations are mathematical choices, not additional axioms.

\begin{definition}[View system]
A view system consists of sets \(F_v\) and maps
\[
 f^*:F_w\longrightarrow F_v\qquad(f:v\longrightarrow w)
\]
satisfying \((1_v)^*=\id\) and \((g\circ f)^*=f^*g^*\). Equivalently, it is a contravariant functor \(F:\V^{\mathrm{op}}\to\Set\).
\end{definition}
The direction matters. An arrow \(v\to w\) induces a restriction from the fiber at \(w\) to the fiber at \(v\). A restriction need not be invertible. A category with several arrows between two views can retain different routes between their descriptions. Replacing it by a preorder can lose those routes and the constraints they impose.

\begin{definition}[Strict Komorebi realization]
The ambient realization set is
\begin{equation}\label{eq:ambient}
 \mathfrak K(F)=\Gamma(F)=
 \left\{x\in\prod_{v\in\Ob\V}F_v:
 f^*(x_w)=x_v\text{ for every }f:v\to w\right\}.
\end{equation}
A point of \(\mathfrak K(F)\) is a whole coherent family. Its value in view \(v\) is the evaluation \(\operatorname{ev}_v(x)=x_v\).
\end{definition}
The raw product contains all families, including inconsistent ones. Equation~\eqref{eq:ambient} cuts out the exact families. No nonemptiness hypothesis is built into the definition.

\begin{theorem}[Joint determination and the universal property]\label{KS01}
The evaluations jointly determine a realization. If a set \(T\) carries maps \(c_v:T\to F_v\) satisfying \(f^*c_w=c_v\), there is a unique map \(c:T\to\Gamma(F)\) satisfying \(\operatorname{ev}_v c=c_v\) for every view.
\end{theorem}
\begin{proof}
We first verify that the evaluation maps themselves form a compatible family. If $x\in\Gamma(F)$ and $f:v\to w$, the defining equation for $x$ gives
\[
 f^*(\operatorname{ev}_w x)=f^*(x_w)=x_v=\operatorname{ev}_v x.
\]
If $x,y\in\Gamma(F)$ have equal evaluations, then $x_v=y_v$ for every $v$. Equality of functions into a dependent product gives equality of their underlying families. Membership in $\Gamma(F)$ is a property of a family, not a further choice of coordinate. Thus $x=y$.

Now fix the compatible maps $c_v:T\to F_v$. For each $t\in T$ define a family by $x(t)_v=c_v(t)$. For every arrow,
\[
 f^*x(t)_w=f^*c_w(t)=c_v(t)=x(t)_v.
\]
This calculation verifies membership in $\Gamma(F)$, so the rule $c(t)=x(t)$ defines an actual function $c:T\to\Gamma(F)$. Evaluating it gives $\operatorname{ev}_vc(t)=c_v(t)$, as required.

To prove uniqueness, let $b:T\to\Gamma(F)$ be another such function. At every $t$ and $v$,
\[
 (b(t))_v=\operatorname{ev}_v b(t)=c_v(t)=\operatorname{ev}_v c(t).
\]
Joint determination first yields $b(t)=c(t)$, and function extensionality then yields $b=c$. Nothing in this argument assumes that $T$ or $\Gamma(F)$ is inhabited. For example, when $T$ is empty there is the unique empty function. When $T$ is inhabited, existence of the compatible family $(c_v)$ is already substantive data.
\end{proof}
\source{KS01}
Thus the strict construction is precisely the ordinary limit of the view diagram, with the usual limit interpretation in the sequential case~\cite{stacks-limits}. This identifies its mathematical content and its universal property. It does not supply a universal view diagram for every possible mathematical question.

\plate{01_views}{The same realization is evaluated in several views. Compatibility is imposed along the arrows, not inferred from the fact that all displayed fibers exist. The ambient family belongs to the simultaneous solution set of the transition equations.}{fig:views}

\subsection{Maps of realizations}
Let \(\alpha:F\to G\) consist of maps \(\alpha_v:F_v\to G_v\) with
\[
 \alpha_v f_F^*=f_G^*\alpha_w.
\]
The induced map on realizations is \(\Gamma(\alpha)(x)_v=\alpha_v(x_v)\). The naturality square is exactly what is needed to show that the image family is coherent.

\begin{proposition}[Functorial transport]\label{KS02}
These induced maps preserve identities and composition:
\[
 \Gamma(\id_F)=\id_{\Gamma(F)},\qquad
 \Gamma(\beta\alpha)=\Gamma(\beta)\Gamma(\alpha).
\]
\end{proposition}
\begin{proof}
Evaluate each side at an arbitrary view and apply the corresponding identity for the coordinate maps. Theorem~\ref{KS01} then identifies the maps on the whole realization set.
\end{proof}
\source{KS02}
An arbitrary collection of fiber maps need not define a map on realizations. Its compatibility equations are part of the theorem, not a convention that can be omitted after a notation has been chosen.

\subsection{A global view and a useful warning}
\begin{proposition}[Terminal evaluation]\label{KS03}
If \(t\) is terminal in \(\V\), evaluation induces an equivalence \(\Gamma(F)\simeq F_t\).
\end{proposition}
\begin{proof}
For \(a\in F_t\), take its restriction along the unique arrow \(v\to t\) at each view. Uniqueness of the arrow to \(t\), together with the functor laws, proves coherence. Evaluation at \(t\) recovers \(a\). Conversely, any coherent family is the restriction of its coordinate at \(t\).
\end{proof}
\source{KS03}
This proposition is both useful and diagnostic. If the supposedly reconstructed global object has already been placed in a terminal fiber, the limit recovers that fiber. The proposition does not prove that arbitrary partial data have a preimage there.

The category of elements \(\int_\V F\) and \(\Gamma(F)\) also have different roles. The former organizes individual appearances and their transitions. The latter selects coherent families across all views. In the basic source layer the category of elements is explanatory notation, while the actual Lean record directly implements families and their coherence.

\chap{Partial data, gluing, and correspondence fibers}\label{sec:glue}
\subsection{Prescribed views}
For \(U\subseteq\Ob\V\), use the full subcategory on \(U\). A partial section \(s\in\Gamma_U(F)\) satisfies all transition equations whose endpoints lie in \(U\). Define
\begin{equation}\label{eq:extension}
 \Ext_U(s)=\{x\in\Gamma(F):x|_U=s\}.
\end{equation}
This fiber can be empty, a singleton, or much larger. A partial section is already coherent on its prescribed domain. A raw assignment that fails an internal transition equation is not a partial section of this kind.

Let \((U_i)_{i\in I}\) be full patches. A matching family consists of \(s_i\in\Gamma_{U_i}(F)\) with equal values on every overlap. The relevant covering condition is
\begin{equation}\label{eq:transitioncover}
 \text{for every }f:v\to w,\quad
 \text{some }U_i\text{ contains both }v\text{ and }w.
\end{equation}
Identities imply that this condition also covers every object.

\begin{theorem}[Strict gluing]\label{KS04}
Under \eqref{eq:transitioncover}, restriction induces an equivalence
\[
 \Gamma(F)\simeq\Match(F,(U_i)).
\]
Every fixed matching family has a unique global glue.
\end{theorem}
\begin{proof}
\emph{Step 1: every coordinate can be selected.} Apply the cover condition to the identity arrow of $v$. It gives an index containing $v$. Choose one such index $i(v)$ and put
\[
 x_v=(s_{i(v)})_v.
\]
The choice is an existence construction. It is not an algorithm selecting patches in a preferred order. If $v\in U_j$, matching on the overlap $U_{i(v)}\cap U_j$ gives
\[
 x_v=(s_j)_v.                                                   \tag{*}
\]
In particular the value is independent of the selected index.

\emph{Step 2: the selected coordinates satisfy every arrow equation.} For $f:v\to w$, choose a single index $j$ with $v,w\in U_j$. Formula (*) applies to both endpoints. Since the patch is full, it contains this actual arrow, even when other arrows have the same endpoints. Coherence of $s_j$ gives
\[
 f^*x_w=f^*(s_j)_w=(s_j)_v=x_v.
\]
Thus $x\in\Gamma(F)$. Object coverage alone would not supply the common patch used in this step.

\emph{Step 3: assembly and restriction are inverse.} Formula (*) proves $x|_{U_j}=s_j$ for every $j$. Conversely, start with $y\in\Gamma(F)$ and assemble its restrictions. The resulting coordinate at $v$ is $y_v$, regardless of which containing patch is chosen. Joint determination gives the original $y$. These two identities are exactly the inverse identities for the displayed equivalence.

For uniqueness with a fixed matching family, any other glue has coordinate $(s_{i(v)})_v=x_v$ at every $v$, hence equals $x$. If the view category is empty, its product has one empty family and the same conclusion follows without making any choices.
\end{proof}
\source{KS04}
The proof explains why object coverage alone is insufficient. Coherence is an equation attached to an arrow, so existence of a glue must see that equation somewhere in the local data.

\begin{example}[An object cover with no glue]\label{KS05}
Take two views \(0\to1\), Boolean fibers, and the identity restriction. The singleton patches cover all objects. A local value false at \(0\) and true at \(1\) agrees on the empty overlap, yet no global section realizes both values. A global family would have to assign the same Boolean to both views.
\end{example}
\source{KS05}
Object coverage does prove uniqueness of a possible glue. It does not prove its existence. This distinction persists when sets are replaced by objects with symmetries.

\subsection{Correspondences must also be coherent}
Suppose natural maps form a span \(F\xleftarrow{p}M\xrightarrow{q}G\), and fix \(a\in\Gamma(F)\), \(b\in\Gamma(G)\). At each view define the ordinary fiber
\[
 B_{a,b}(v)=\{m\in M_v:p_v(m)=a_v,\ q_v(m)=b_v\}.
\]
Naturality and coherence provide its restrictions. A global bridge is an element of
\[
 \mathfrak B(a,b)=\{z\in\Gamma(M):\Gamma(p)z=a,\ \Gamma(q)z=b\}.
\]

\begin{theorem}[Fiber and section compatibility]\label{KS06}
There is a coordinate preserving equivalence
\[
 \Gamma(B_{a,b})\simeq\mathfrak B(a,b).
\]
Under a transition cover, matching local bridge sections glue to a unique global bridge with both prescribed projections unchanged.
\end{theorem}
\begin{proof}
We must first check that the displayed fibers really form a view system. Let $m\in B_{a,b}(w)$ and $f:v\to w$. Naturality of $p$ and coherence of $a$ give
\[
 p_v(f_M^*m)=f_F^*p_w(m)=f_F^*a_w=a_v.
\]
The same calculation with $q,b$ gives $q_v(f_M^*m)=b_v$. Therefore $f_M^*m$ belongs to $B_{a,b}(v)$. Identity and composition of these restricted maps are inherited from $M$, because equality in a subset is equality of its underlying elements.

Given $u\in\Gamma(B_{a,b})$, forget each fiber-membership condition and call the underlying family $z_v\in M_v$. Coherence of $u$ is precisely coherence of $z$, so $z\in\Gamma(M)$. The equations $p_vz_v=a_v$ hold at every coordinate. Joint determination in $\Gamma(F)$ yields $\Gamma(p)z=a$, and the equations for $q$ similarly yield $\Gamma(q)z=b$. Thus $z$ is a global bridge.

In the opposite direction, let $z\in\mathfrak B(a,b)$. Evaluate its two global projection equations at $v$. They show that $z_v$ belongs to $B_{a,b}(v)$, and its original coherence shows that these elements form a section of the fiber system. Forgetting and reattaching membership conditions preserves every coordinate, so the two constructions are inverse.

Finally apply strict gluing to this specific fiber system. The glued section retains its membership in every $B_{a,b}(v)$, which is why both prescribed projections survive gluing. Uniqueness is for the fixed matching bridge family. Different matching families for the same $(a,b)$ can still yield different global bridges.
\end{proof}
\source{KS06}
These are ordinary fibers. They do not retain comparison isomorphisms or automorphisms. Section~\ref{sec:relative} constructs the appropriate groupoid version when this information matters.

\subsection{Locally inhabited does not mean globally inhabited}
\begin{example}[The loop with no coherent choice]\label{KS07}
Let the view category have one object and automorphism group \(C_2\). Let its nonidentity arrow act on a two element set by swapping the elements. The fiber is nonempty and every transport is bijective. Nevertheless a global section would be a fixed point of the swap, and there is none.
\end{example}
\source{KS07}
The problem is not missing local elements or noninvertible transition maps. It is the absence of an element satisfying the loop equation. Inserting a point in each fiber does not certify a coherent selection of those points.

\begin{insight}
The strict theory already distinguishes three levels: possible local appearances, compatible partial descriptions, and an actual global realization. A correspondence must pass through the same levels. Local nonemptiness of its fibers is not a globalization theorem.
\end{insight}

\chap{The additive obstruction and the space of extensions}\label{sec:obstruction}
The strict definition gives a system of equations. In an additive setting, that system has an actual differential, and the failure to extend prescribed data has a quotient valued obstruction.

\subsection{The compatibility differential}
Assume each fiber \(A_v\) is an abelian group and each restriction is additive. Set
\[
 E=\prod_v A_v,\qquad R=\prod_{f:v\to w}A_v,
 \qquad (\delta x)_f=f^*(x_w)-x_v.
\]
The group \(R\) retains an entry for every arrow. It is not silently identified with a cohomology group.

\begin{proposition}[Differential identities]\label{KS08}
The map \(\delta:E\to R\) is additive, \(\Gamma(A)=\ker\delta\), and
\[
 (\delta x)_{1_v}=0,\qquad
 (\delta x)_{g\circ f}=f^*((\delta x)_g)+(\delta x)_f.
\]
\end{proposition}
\begin{proof}
Additivity follows coordinate by coordinate. Vanishing is exactly the defining coherence equation. For composition, add and subtract \(f^*x_w\):
\[
 f^*g^*x_z-x_v
 =f^*(g^*x_z-x_w)+(f^*x_w-x_v).
\]
The identity equation uses \((1_v)^*=\id\).
\end{proof}
\source{KS08}
This exact formula is more informative than a scalar error score. It specifies how defects transform and combine. Norms can be added later, but they do not replace the typed residual.

\subsection{The admissible corrections}
Fix \(U\subseteq\Ob\V\). Write \(p_U:E\to\prod_{v\in U}A_v\) for raw restriction. A correction preserving the prescribed values belongs to \(\ker p_U\). Its possible defects form the subgroup
\[
 B_U=\delta(\ker p_U)\le R.
\]
For a partial section \(s\), choose any raw extension \(\widetilde s\), for example zero off \(U\), and define
\begin{equation}\label{eq:obstruction}
 Q_U=R/B_U,\qquad
 \mathfrak o_U(s)=[\delta\widetilde s]\in Q_U.
\end{equation}
If another raw extension is chosen, their difference vanishes on \(U\), so its differential belongs to \(B_U\). Thus the class is independent of the lift.

\begin{theorem}[Extension criterion]\label{KS09}
For every coherent partial section,
\[
 \mathfrak o_U(s)=0\quad\Longleftrightarrow\quad\Ext_U(s)\ne\varnothing.
\]
\end{theorem}
\begin{proof}
Let $\pi_U:R\to R/B_U$ be the quotient homomorphism. By the definition of a quotient by a subgroup,
\[
 \pi_U(r)=0\quad\Longleftrightarrow\quad r\in B_U.
\]
Here $B_U$ is the \emph{image} of $\ker p_U$ under $\delta$. Consequently membership of $\delta\widetilde s$ means the existence of a specific $h\in E$ satisfying both $p_Uh=0$ and $\delta h=\delta\widetilde s$. The subgroup need not be closed in any topology. No topology has been introduced.

If such $h$ exists, set $x=\widetilde s-h$. Additivity gives the two equations separately:
\[
 \delta x=\delta\widetilde s-\delta h=0,
 \qquad p_Ux=p_U\widetilde s-p_Uh=s.
\]
The first equation puts $x$ in $\Gamma(A)$, and the second puts that coherent family in $\Ext_U(s)$.

Conversely, take $x\in\Ext_U(s)$ and set $h=\widetilde s-x$. Since both raw families restrict to $s$, one has $p_Uh=0$. Since $\delta x=0$, one has $\delta h=\delta\widetilde s$. This witnesses $\delta\widetilde s\in B_U$, hence vanishing of its quotient class.

The proof is independent of the selected raw lift. Indeed, if $\widetilde s'$ is another lift, $k=\widetilde s'-\widetilde s$ lies in $\ker p_U$, and
\[
 \delta\widetilde s'-\delta\widetilde s=\delta k\in B_U.
\]
Thus the two quotients agree before either feasibility direction is invoked. A raw lift always exists in this additive model by using zero outside $U$. Constructing a raw lift does not solve the coherence equations.
\end{proof}
\source{KS09}
The quotient is constructed before the solvability test. The theorem therefore does not define obstruction circularly by saying that it vanishes when a solution exists. It also does not give a decision procedure for membership in an arbitrary subgroup of an infinite product.

\plate{02_extension}{The raw family is repaired only along directions invisible on the prescribed views. The obstruction records whether its defect lies in the image of those admissible directions. Once an extension exists, a different kernel measures the remaining freedom.}{fig:extension}

\subsection{Existence is not uniqueness}
Define the ambiguity group
\[
 H_U=\ker\delta\cap\ker p_U.
\]
It consists of coherent changes that are invisible on the prescribed domain.

\begin{theorem}[Classification of a nonempty extension fiber]\label{KS10}
If \(x_0\in\Ext_U(s)\), translation gives an equivalence
\[
 H_U\simeq\Ext_U(s),\qquad h\longmapsto x_0+h.
\]
Consequently the extension is unique exactly when \(\mathfrak o_U(s)=0\) and \(H_U=0\).
\end{theorem}
\begin{proof}
Define $\Theta:H_U\to\Ext_U(s)$ by $\Theta(h)=x_0+h$. It lands in the extension set because
\[
 \delta(x_0+h)=\delta x_0+\delta h=0,
 \qquad p_U(x_0+h)=s+0=s.
\]
Define $\Psi:\Ext_U(s)\to H_U$ by $\Psi(x)=x-x_0$. Its differential is zero, and its restriction is $s-s=0$, so the claimed target is correct. The formulas give
\[
 \Psi\Theta(h)=x_0+h-x_0=h,\qquad
 \Theta\Psi(x)=x_0+(x-x_0)=x.
\]
This proves the equivalence rather than only a comparison of cardinalities.

If $H_U=\{0\}$, the parameter of every extension is zero, so every extension equals $x_0$. Conversely, if $x_0$ is the only extension, then for any $h\in H_U$ the equality $x_0+h=x_0$ forces $h=0$ by cancellation. Combining this observation with the existence criterion proves the stated necessary and sufficient condition for a unique extension.

The choice of $x_0$ only supplies an origin for the parameters. If $x_1=x_0+h_1$ is used instead, the parameter of $x$ changes from $x-x_0$ to $x-x_1=(x-x_0)-h_1$. The affine set of answers itself has not changed.
\end{proof}
\source{KS10}
An extension fiber is therefore an affine copy, or torsor, of the invisible coherent group. Selecting one solution gives coordinates on the fiber. Without a selected solution, the affine structure does not designate an origin.

\begin{example}[Parity and an independent ambiguity]
For two views with \(f^*(m)=2m\) on \(\Z\), a prescribed value \(k\) at the coarse view extends exactly when \(k\) is even. The extension is \((k,k/2)\) and is unique. In particular \(1\) is obstructed and \(2\) extends nontrivially. The parity statement is an instance of Theorem~\ref{KS09} in the original source. If a separate, disconnected integer view is added and left unprescribed, every extendable seed has infinitely many completions. This last variant is a direct application of Theorem~\ref{KS10}, not a claim that the parity obstruction has changed.
\end{example}
\source{KS11}

\subsection{The exact sequence behind the criterion}
The additive structures make restriction a homomorphism
\(r_U:\Gamma(A)\to\Gamma_U(A)\). Using the zero extension, or lift independence, \(\mathfrak o_U\) is also a homomorphism.

\begin{proposition}[Exactness at partial sections]\label{KS12}
The sequence
\[
 \Gamma(A)\xrightarrow{r_U}\Gamma_U(A)
 \xrightarrow{\mathfrak o_U}Q_U
\]
is exact at \(\Gamma_U(A)\). Thus \(\ker\mathfrak o_U=\im r_U\), and restriction is onto precisely when \(\mathfrak o_U\) is the zero map.
\end{proposition}
\begin{proof}
First verify that $\mathfrak o_U$ is a homomorphism, since exactness is being asserted in an additive setting. Given $s,t\in\Gamma_U(A)$, choose raw lifts $\widetilde s,\widetilde t$. Their sum lifts $s+t$. Lift independence and additivity imply
\[
 \mathfrak o_U(s+t)
 =[\delta(\widetilde s+\widetilde t)]
 =[\delta\widetilde s]+[\delta\widetilde t]
 =\mathfrak o_U(s)+\mathfrak o_U(t).
\]
The zero lift gives $\mathfrak o_U(0)=0$. Restriction $r_U$ is a homomorphism because addition of global sections is coordinatewise.

If $s=r_Ux$ for a global section $x$, use $x$ itself as a raw lift. Its differential vanishes, so $\mathfrak o_U(s)=0$. This proves $\im r_U\subseteq\ker\mathfrak o_U$. If $s\in\ker\mathfrak o_U$, the extension criterion constructs a global $x$ restricting to $s$, proving the reverse inclusion.

Finally, $r_U$ is surjective exactly when its image is all of $\Gamma_U(A)$. By the equality just proved, this is equivalent to every partial section having zero obstruction, which says that $\mathfrak o_U$ is the zero homomorphism. Nothing here claims that its image is all of $Q_U$.
\end{proof}
\source{KS12}
No surjectivity onto the entire residual quotient is asserted. The residual group contains many entries that need not arise as the differential of a partial lift.

\subsection{Extension in stages}
For \(U\subseteq W\subseteq\Ob\V\), an extension first to \(W\) and then globally must retain the intermediate choice. There is an explicit equivalence
\begin{equation}\label{eq:stages}
 \Ext_U(s)\simeq
 \coprod_{\substack{t\in\Gamma_W(A)\\t|_U=s}}\Ext_W(t).
\end{equation}
It sends a global extension to its restriction to \(W\), together with itself as an extension of that restriction. The inverse forgets the intermediate coordinate. These operations are inverse because the required intermediate coordinate of a global section is its actual restriction. A coherent intermediate choice can still be obstructed, even when a different intermediate choice completes. Equation~\eqref{eq:stages} is why extension by stages is a dependent sum rather than a statement that every branch succeeds.
\source{KS13}

\chap{Protected correction as a general principle}\label{sec:protected}
The additive extension theorem is a particular instance of a useful general construction. Let \(E,R,S\) be abelian groups, with additive maps \(d:E\to R\) and \(p:E\to S\). The residual \(d\) measures compatibility, and \(p\) specifies the information that repair must preserve.

\begin{theorem}[Semantics preserving correction]\label{KS14}
For \(x\in E\), the following conditions are equivalent:
\[
\begin{aligned}
 &[dx]=0\text{ in }R/d(\ker p),\\
 &\exists h\in E:\ p(h)=0,\quad d(h)=d(x),\\
 &\exists y\in E:\ d(y)=0,\quad p(y)=p(x).
\end{aligned}
\]
The class of \(dx\) is unchanged when an admissible correction is subtracted from \(x\).
\end{theorem}
\begin{proof}
The image $d(\ker p)$ is a subgroup of $R$. The defining property of the quotient says
\[
 [dx]=0\quad\Longleftrightarrow\quad
 dx\in d(\ker p)
 \quad\Longleftrightarrow\quad
 \exists h\in E:\ ph=0,\ dh=dx.
\]
This is the first equivalence with its witness made explicit.

For the second, a permitted $h$ gives $y=x-h$. Then
\[
 dy=dx-dh=0,\qquad py=px-ph=px.
\]
Conversely an exact $y$ with $py=px$ gives $h=x-y$, which satisfies $ph=0$ and $dh=dx$. The two witness transformations are inverse: subtracting $h$ from $x$ and then subtracting that answer from $x$ recovers $h$, and conversely.

To verify invariance, fix any $h$ with $ph=0$. Then $d(x-h)=dx-dh$, and $dh$ lies in the subgroup divided out. Therefore $[d(x-h)]=[dx]$. This does not assert that the residual itself is unchanged. It asserts that permitted changes can alter the residual only within its already correctable coset.
\end{proof}
\source{KS14}
Here \emph{semantics} means the declared mathematical observable \(p\). It is not an assertion that one homomorphism captures every meaning of every object. In a tower, \(p\) might preserve a seed operator and every new output component. In a graph, it might preserve designated vertex values. Different choices give different problems.

\begin{proposition}[A certificate of impossibility]\label{KS15}
Let \(L:R\to D\) be additive. If \(L\) vanishes on \(d(\ker p)\) but \(L(dx)\ne0\), no exact \(y\) can preserve \(p(x)\).
\end{proposition}
\begin{proof}
Suppose, for contradiction, that $dy=0$ and $py=px$. Then $h=x-y$ satisfies $ph=0$ and $dh=dx$. By the annihilation hypothesis,
\[
 L(dx)=L(dh)=0,
\]
contrary to the assumed nonzero value. This already proves impossibility directly.

Equivalently, if $r-r'\in d(\ker p)$, then $L(r-r')=0$, so $L(r)=L(r')$. Hence $\overline L([r])=L(r)$ is a well-defined homomorphism on $R/d(\ker p)$. Its value on $[dx]$ is nonzero. This supplies the quotient interpretation of the same contradiction, without assuming that every nonzero quotient class can be detected by a map into one fixed preferred coefficient group.
\end{proof}
\source{KS15}
This is often the most economical negative proof. One need not classify every possible correction. It is enough to find an invariant annihilating all admissible corrections but not the observed defect.

\subsection{Why unrestricted cancellation is empty information}
If every correction is permitted, taking \(h=x\) leaves the exact candidate zero. Thus \(\exists h:d(x-h)=0\) is automatic. It says nothing about recovering an object of interest. A repair theorem becomes informative only when it preserves something, is quantitatively controlled, or satisfies a comparison property that excludes arbitrary replacement.

\begin{example}[A rigid obstruction]
Take \(E=R=S=\Z\), \(d=p=\id\), and \(x=1\). The only admissible correction is zero, while the residual is nonzero. Hence there is no exact answer preserving the observed value. This is not a failure of the formalism. It is the correct certificate for an incompatible demand.
\end{example}

\subsection{What an obstruction does not yet provide}
An algebraic quotient does not automatically carry a useful norm. A zero infimum of distances to its boundary subgroup does not imply membership unless closure is controlled. Nor does vanishing supply a preferred correction when \(\ker d\cap\ker p\) is nontrivial. These two issues lead to different parts of the book: ambiguity is treated algebraically and through groupoids, while closed range and condition numbers enter the analytic theory.

\begin{insight}
A repair problem is not determined by its defect alone. It is determined by the defect \emph{and} the admissible changes. Keeping the pair \((d,p)\) explicit prevents a successful cancellation from being mistaken for a successful reconstruction.
\end{insight}

\input{chapters/02a_prescription}
\partnote{Resolution, repair, and retained detail}{Resolution gives a concrete family of views. Its compatibility differential connects exact maps, protected information, and residual coordinates. Recovering every finite coordinate is then separated from convergence in a chosen norm.}
\section{Resolution towers and exact transport}\label{sec:towers}
\subsection{The sequential model}
A tower in a category \(\mathcal C\) is a sequence
\[
 \cdots\xrightarrow{\rho_2}X_2\xrightarrow{\rho_1}X_1
 \xrightarrow{\rho_0}X_0.
\]
For sets, a point is a sequence \(x_n\in X_n\) with \(\rho_n(x_{n+1})=x_n\). Composite restrictions turn the tower into a view system on \(\N\), with an arrow \(i\to j\) when \(i\le j\). Its section space is the same coherent point type, with the same coordinates.

A candidate transport from \(X\) to \(Y\) consists of maps \(f_n:X_n\to Y_n\). It is exact when
\begin{equation}\label{eq:exacttower}
 \rho^Y_n f_{n+1}=f_n\rho^X_n\quad\text{for all }n.
\end{equation}

\begin{proposition}[The exact tower category]\label{KS16}
Identity maps and composites of exact transports are exact. An exact transport induces a map of coherent point spaces by \(x\mapsto(f_nx_n)_n\), preserving identities and composition.
\end{proposition}
\begin{proof}
For two exact transports, substitute \(\rho^Z_ng_{n+1}=g_n\rho^Y_n\) and \(\rho^Y_nf_{n+1}=f_n\rho^X_n\) into the composite square. For a coherent point,
\[
 \rho^Y_n f_{n+1}x_{n+1}
 =f_n\rho^X_nx_{n+1}=f_nx_n.
\]
Identity, associativity, and composition on points are checked coordinatewise.
\end{proof}
\source{KS16}
This is a category on a fixed schedule. A full procategory allows additional reindexing and equivalence relations on maps. Those constructions must be supplied rather than identified with the fixed schedule by terminology.

\subsection{The actual residual of a transport}
Assume \(\mathcal C\) is preadditive, so differences of maps make sense. The tower differential is
\[
 (\partial f)_n=\rho^Y_nf_{n+1}-f_n\rho^X_n.
\]
Its target has components \(\Hom(X_{n+1},Y_n)\), not \(\Hom(X_n,Y_n)\).

\begin{proposition}[The transport Leibniz identity]\label{KS17}
For composable candidates,
\[
 \partial(gf)_n=(\partial g)_nf_{n+1}+g_n(\partial f)_n.
\]
In particular exactness is equivalent to vanishing of the differential.
\end{proposition}
\begin{proof}
Expand the right side. The terms \(-g_n\rho^Y_nf_{n+1}\) and \(+g_n\rho^Y_nf_{n+1}\) cancel, leaving \(\rho^Z_ng_{n+1}f_{n+1}-g_nf_n\rho^X_n\). This is the differential of the composite. The exactness statement is \eqref{eq:exacttower}.
\end{proof}
\source{KS17}
There is no additional mixed defect term in this identity. Such a term would require a different definition of the defect or a different expansion. Typed cancellation is a useful safeguard against importing an attractive but incorrect error formula.

\subsection{Cofinal views preserve points, not every measurement}
Let \(\phi:\N\to\N\) be monotone and unbounded. The sampled tower uses levels \(X_{\phi(n)}\) and their composite restrictions.

\begin{theorem}[Cofinal reconstruction of points]\label{KS18}
Sampling induces an equivalence between the coherent points of \(X\) and of its sampled tower.
\end{theorem}
\begin{proof}
Write $\rho_{i,j}:X_j\to X_i$ for the composite restriction when $i\le j$, with $\rho_{i,i}=I$. These maps satisfy $\rho_{i,j}\rho_{j,k}=\rho_{i,k}$. A coherent sampled point obeys
\[
 \rho_{\phi(n),\phi(m)}y_m=y_n\qquad(n\le m),
\]
by iterating its adjacent coherence equations.

For each $k$, unboundedness gives an index $n$ with $k\le\phi(n)$. Set $x_k=\rho_{k,\phi(n)}y_n$. If $m$ is another permissible index, choose $\ell\ge n,m$. Monotonicity permits both composite restrictions, and sampled coherence gives
\[
 \rho_{k,\phi(n)}y_n
 =\rho_{k,\phi(n)}\rho_{\phi(n),\phi(\ell)}y_\ell
 =\rho_{k,\phi(\ell)}y_\ell
 =\rho_{k,\phi(m)}y_m.
\]
Thus $x_k$ is independent of the choice, including when the schedule repeats a level.

For the adjacent equation at $k$, choose one $n$ with $k+1\le\phi(n)$. Then
\[
 \rho_{k,k+1}x_{k+1}
 =\rho_{k,k+1}\rho_{k+1,\phi(n)}y_n
 =\rho_{k,\phi(n)}y_n=x_k.
\]
The recovered family is therefore a point of the original tower.

Sampling a full point and recovering its $k$th coordinate gives $\rho_{k,\phi(n)}x_{\phi(n)}=x_k$ by its original coherence. Recovering a sampled point and evaluating at $k=\phi(n)$ allows the index $n$ itself, giving $\rho_{\phi(n),\phi(n)}y_n=y_n$. These are both inverse identities. No surjectivity of the bonding maps or of the point evaluations was used.
\end{proof}
\source{KS18}
This theorem is about exact coherent points. Raw residual size is a different quantity. On identity scalar towers, \(f_n=(-1)^n\) has nonzero fine differences, but its even subsequence is constant. The coarse residual vanishes because adjacent fine defects cancel. A cofinal equivalence of limits does not imply equality of residual norms or an absence of aliasing in uncorrected data.

\plate{03_tower}{Exact transport commutes with every adjacent restriction. A cofinal subsequence retains coherent points, but a raw defect can cancel during coarsening. The retained coarse value alone need not record the omitted fine discrepancy.}{fig:tower}

\chap{Finite observations and realized transformations}\label{sec:observation}
The point-space construction has a converse question. Given a function between coherent realization spaces, can it be represented by compatible maps between finite views? The answer is not automatically yes. This section constructs the relevant category of reindexed maps and identifies its realized functions under an explicit observability hypothesis.

\subsection{An output level may require a finer input level}
Let \(X\) and \(Y\) be sequential set-valued resolutions. Write \(\rho^X_{i,j}:X_j\to X_i\) for composite restriction. A reindexed transport consists of a monotone schedule \(\phi:\N\to\N\) and maps
\[
 f_n:X_{\phi(n)}\longrightarrow Y_n
\]
satisfying, for \(i\le j\),
\begin{equation}\label{eq:scheduled-naturality}
 \rho^Y_{i,j}f_j=f_i\rho^X_{\phi(i),\phi(j)}.
\end{equation}
The schedule records how much input is used for each output level. Monotonicity means that asking for more output never forces the representation to discard input already available. It does not require strict growth, a bounded look-ahead, or an efficiency estimate.

The realized map is
\[
 |f|:\operatorname{Point}(X)\to\operatorname{Point}(Y),\qquad
 |f|(x)_n=f_n(x_{\phi(n)}).
\]
Equation~\eqref{eq:scheduled-naturality} proves that this is a coherent output. A second transport \(g_n:Y_{\psi(n)}\to Z_n\) composes with \(f\) using schedule \(\phi\circ\psi\) and coordinate map \(g_n f_{\psi(n)}\). The order of the schedules is fixed by the types, not by the order in which their letters are written in prose.

\subsection{Equality is witnessed at finite stages}
Two different schedules can represent the same transformation. To compare \((\phi,f)\) and \((\psi,g)\), require that for every output level \(n\) there be a finite level \(k\ge\phi(n),\psi(n)\) with
\begin{equation}\label{eq:eventual-equality}
 f_n\rho^X_{\phi(n),k}=g_n\rho^X_{\psi(n),k}
 \quad\text{as functions on all of }X_k.
\end{equation}
Call this eventual equality. The witnessing \(k\) may depend on \(n\). Equality is required on the whole finite-level fiber, not only on states that happen to extend globally.

\begin{theorem}[The eventual transport category]\label{KS61}
Eventual equality is an equivalence relation and is preserved by composition. Its quotient therefore gives a category of resolutions and eventual transports. Realization descends to a functor into sets, preserving identities and composition.
\end{theorem}
\begin{proof}
A finite-stage equality remains true after moving to a later stage. Indeed, if it holds on $X_k$ and $\ell\ge k$, precompose both sides with $\rho^X_{k,\ell}$ and use the composite-restriction laws. Reflexivity uses $k=\phi(n)$, symmetry reverses the equality, and transitivity moves two witnesses to their maximum. The shared middle presentation then cancels as a literal function on that common fiber. This proves an equivalence relation without referring to global points.

We check compatibility with composition with all indices displayed. Write $f,f'$ with schedules $\phi,\phi'$ and $g,g'$ with schedules $\psi,\psi'$. Fix an output level $n$. Choose $j\ge\psi(n),\psi'(n)$ with
\[
 g_n\rho^Y_{\psi(n),j}=g'_n\rho^Y_{\psi'(n),j}.
\]
Next choose $k\ge\phi(j),\phi'(j)$ with
\[
 f_j\rho^X_{\phi(j),k}=f'_j\rho^X_{\phi'(j),k}.
\]
Monotonicity implies that $k$ is also at least both composite input levels. Applying the naturality equations and then the two witness equations gives
\[
\begin{aligned}
 g_nf_{\psi(n)}\rho^X_{\phi(\psi(n)),k}
 &=g_n\rho^Y_{\psi(n),j}f_j\rho^X_{\phi(j),k}\\
 &=g'_n\rho^Y_{\psi'(n),j}f_j\rho^X_{\phi(j),k}\\
 &=g'_n\rho^Y_{\psi'(n),j}f'_j\rho^X_{\phi'(j),k}\\
 &=g'_nf'_{\psi'(n)}\rho^X_{\phi'(\psi'(n)),k}.
\end{aligned}
\]
This is the required equality on every element of $X_k$. Thus composition descends to classes. The representative composition is associative because function composition and schedule composition are associative. The schedule $n\mapsto n$ with identity coordinate maps gives both identity laws. Hence the quotient has the claimed category structure.

Finally, evaluate a witness equation at the $k$th coordinate of a coherent point $x$. Its two restricted arguments become $x_{\phi(n)}$ and $x_{\psi(n)}$, proving equality of the realized outputs. Thus realization is well defined on the quotient. For a composite, direct evaluation gives
\[
 |g\circ f|(x)_n=g_n\bigl(f_{\psi(n)}(x_{\phi(\psi(n))})\bigr)
 =(|g|(|f|(x)))_n.
\]
The identity law is equally coordinatewise. Joint determination of output points proves the functor laws.
\end{proof}
\source{KS61}
This is an actual quotient category in the formal source. No identification with the full procategory of an arbitrary base category is needed or claimed. The objects here are sequential set-valued resolutions, and the equivalence relation is the finite-stage relation just defined.

\subsection{When global behavior determines the finite presentation}
Say that \(X\) has surjective point projections when
\[
 \forall k\quad \operatorname{ev}_k:\operatorname{Point}(X)\to X_k
 \quad\text{is onto}.
\]
This means that every finite input state really occurs in some coherent global input. It is stronger than nonemptiness of the global space.

\begin{theorem}[Faithfulness on an observable source]\label{KS62}
If \(X\) has surjective point projections, two reindexed maps out of \(X\) are eventually equal exactly when they induce the same realized function. Surjective adjacent bonding maps are a sufficient condition for surjective point projections.
\end{theorem}
\begin{proof}
Realized equality follows from eventual equality by the preceding theorem, so only its converse remains. Suppose \(|f|=|g|\). Fix the output level \(n\), and let \(k=\max\{\phi(n),\psi(n)\}\). To prove equality of two maps defined on \(X_k\), fix an arbitrary element \(a\in X_k\). The assumed surjectivity of \(\operatorname{ev}_k\) supplies a coherent point \(x\) with \(x_k=a\). Equality of the two realized maps gives equality of their output coordinates:
\[
 f_n(x_{\phi(n)})=g_n(x_{\psi(n)}).
\]
Coherence of \(x\) rewrites this as
\[
 f_n\bigl(\rho^X_{\phi(n),k}(a)\bigr)
 =g_n\bigl(\rho^X_{\psi(n),k}(a)\bigr).
\]
The index \(k\) was selected before \(a\), so this is one common finite-stage equality for \emph{all} elements of \(X_k\). It is exactly the witness required by eventual equality at output \(n\). Repeating for each \(n\) proves faithfulness.

For the sufficient condition on bonding maps, fix \(k\) and \(a\in X_k\). Define coordinates below \(k\) by restriction of \(a\). At each later stage use surjectivity of the adjacent bonding map to choose a preimage of the coordinate already chosen. Recursion with classical choice constructs the upper tail. Adjacent equations hold by construction, and induction on the interval length gives every composite equation. This produces a point evaluating to \(a\). If a fiber is empty, the statement that each of its elements lifts is vacuous, as required by surjectivity onto that fiber. Without the projection-surjectivity assumption, equality on realizable inputs alone need not identify finite maps on values that never occur in a global point.
\end{proof}
\source{KS62}
The hypothesis is sufficient, not presented as a characterization of every source on which faithfulness happens to hold. Its role in the proof is precise: it converts a statement about globally realized inputs into a statement about all finite inputs.

Without it, global behavior can lose information. For example, use \(X_n=\{m\in\N:m\ge n\}\) with inclusion bonds and a constant Boolean output tower. The source point space is empty. The two transports that always output false and always output true induce the same function from the empty set, but disagree at every finite stage. Every \(X_k\) is nonempty, so no eventual equality witness exists. This calculation uses the finite-prefix boundary of KS54 and does not identify empty global behavior with equality of finite presentations.

\subsection{The exact image of realization}
For a function \(F:\operatorname{Point}(X)\to\operatorname{Point}(Y)\), define finite observation by the quantifiers
\begin{equation}\label{eq:finite-observation}
 \forall n\ \exists k\ \exists a:X_k\to Y_n\ \forall x:\quad
 F(x)_n=a(x_k).
\end{equation}
The input level is uniform over all global inputs for each fixed output level. It may vary with the output level. This is not the weaker condition that each individual input can be inspected to some input-dependent finite depth.

\begin{theorem}[Characterization of realized transformations]\label{KS63}
If \(X\) has surjective point projections, realization induces a bijection
\[
 \operatorname{EventualHom}(X,Y)
 \simeq
 \{F:\operatorname{Point}(X)\to\operatorname{Point}(Y):F
 \text{ is finitely observed}\}.
\]
\end{theorem}
\begin{proof}
For a reindexed map $(\phi,f)$, the formula $|f|(x)_n=f_n(x_{\phi(n)})$ immediately supplies a finite-observation witness for each $n$. Eventual equality does not affect the realized function by KS61, so this defines a map from the quotient to the stated set of functions.

For surjectivity, let $F$ be finitely observed. Choose one pair $(k_n,a_n)$ at each output level, satisfying $F(x)_n=a_n(x_{k_n})$ for every coherent input. Define
\[
 \phi(n)=\max_{0\le i\le n}k_i,
 \qquad f_n=a_n\rho^X_{k_n,\phi(n)}.
\]
The maximum is of a nonempty finite set. It gives a monotone schedule with $k_n\le\phi(n)$. For every point $x$, coherence in $X$ gives
\[
 f_n(x_{\phi(n)})
 =a_n\bigl(\rho^X_{k_n,\phi(n)}x_{\phi(n)}\bigr)
 =a_n(x_{k_n})=F(x)_n.                                      \tag{*}
\]

The remaining obligation is naturality on \emph{all} finite states. Fix $i\le j$ and $u\in X_{\phi(j)}$. By surjectivity of that point projection, choose $x$ with $x_{\phi(j)}=u$. Equation (*) and output coherence yield
\[
\begin{aligned}
 \rho^Y_{i,j}f_j(u)&=\rho^Y_{i,j}F(x)_j=F(x)_i,\\
 f_i\rho^X_{\phi(i),\phi(j)}(u)&=f_i(x_{\phi(i)})=F(x)_i.
\end{aligned}
\]
The state $u$ was arbitrary. Hence the two finite-level functions are equal, which proves scheduled naturality. The resulting transport realizes $F$ by (*).

Different witness choices can give different schedules and coordinate maps, but the realized function is still $F$. By KS62 these transports are eventually equal, so their quotient class is independent of all choices. The same faithfulness result says that two classes with equal realized functions are identical. We have proved both surjectivity and injectivity, including the well-definedness of the inverse choice construction. No bound on the growth of $\phi$ follows from these quantifiers.
\end{proof}
\source{KS63}
The finite-level compatibility is derived in this proof. It is not inserted as an additional field in the hypothesis that \(F\) is finitely observed. This is one of the useful bridges between a global description of a transformation and an actual presentation in the resolution calculus.

\plate{20_schedule}{Independent finite-observation witnesses need not use a monotone input schedule. Taking the running maximum produces one without changing the realized function. The map at the enlarged level first restricts to the originally sufficient level. This construction gives no efficiency bound on how large those levels are.}{fig:schedule}

\subsection{Binary prefixes: a complete interpretation}
Let \(X_n\) consist of binary words of length \(n\), with prefix restriction. A coherent point is exactly an infinite binary sequence: its length-\(n\) view is its first \(n\) bits. Every prefix extends, for example by appending zeros. Hence the source satisfies the hypothesis of the characterization theorem.

Bitwise complementation is represented at level \(n\) by complementing the first \(n\) bits. The left shift \(F(x)_i=x_{i+1}\) uses the schedule \(\phi(n)=n+1\). A nearest-neighbor difference over \(\mathbb F_2\),
\[
 G(x)_i=x_i+x_{i+1},
\]
also uses \(n+1\) input bits for \(n\) output bits. Applying \(G\) twice uses \(n+2\) bits, and the middle terms cancel:
\[
 G(G(x))_i=x_i+x_{i+2}.
\]
The compatible finite maps make each statement meaningful on the entire coherent sequence, not merely on one chosen truncation. These are worked instances of the scheduled formulas. The finite regression suite checks their prefix compatibility independently.

Not every set function on these points is finitely observed. Consider the Boolean output that is one exactly when every input bit is zero, regarded as a point of the constant Boolean tower. Given any proposed finite level \(k\), compare the all-zero sequence with a sequence whose first one occurs at position \(k\). Their first \(k\) bits agree, but the outputs differ. Therefore no finite witness \(a:X_k\to\{0,1\}\) can compute even that one output coordinate. This rules out a reindexed presentation by the forward direction of the theorem.

The obstruction in this example is one of information, not a runtime benchmark. No amount of computation on the specified finite prefix distinguishes those two inputs. The theorem does not prohibit a different mathematical representation that supplies more information as part of each view. Changing the available view is changing the input model.

\subsection{Cofinal changes also preserve the morphism calculus}
A monotone unbounded schedule \(\theta\) gives the sampled resolution \(X_\theta\). There is a transport from \(X\) to \(X_\theta\) that reads the selected level. For the reverse direction choose, for each \(n\), a level \(j(n)\) with \(n\le\theta(j(n))\), and restrict the available value down to \(X_n\). Taking a least such index gives a monotone schedule.

\begin{theorem}[Cofinal invariance of eventual transports]\label{KS64}
Sampling and recovery are inverse in the eventual transport category. Consequently, for monotone unbounded schedules \(\theta\) and \(\eta\), there is a bijection
\[
 \operatorname{EventualHom}(X,Y)
 \simeq
 \operatorname{EventualHom}(X_\theta,Y_\eta).
\]
On point spaces it agrees with the previously constructed cofinal realization equivalences.
\end{theorem}
\begin{proof}
Let \(s:X\to X_\theta\) be sampling, with schedule \(\theta\) and identity level maps. For recovery choose the least \(j(n)\) with \(n\le\theta(j(n))\). Unboundedness makes this set nonempty. Increasing \(n\) only strengthens the inequality, so \(j\) is monotone. The level map of \(r:X_\theta\to X\) is \(\rho^X_{n,\theta(j(n))}\). Composite restriction proves its naturality.

The composite \(rs:X\to X\) reads level \(\theta(j(n))\) and restricts to \(n\). At that very stage its map agrees with the identity transport after restriction from \(\theta(j(n))\) to \(n\). Thus \(rs\) is eventually equal to the identity.

For \(sr:X_\theta\to X_\theta\), the schedule at output \(n\) is \(j(\theta(n))\), and the level map is \(\rho^X_{\theta(n),\theta(j(\theta(n)))}\). Compare it with the identity at a sampled index
\[
 k\ge\max\{n,j(\theta(n))\}.
\]
Restricting either route from level \(k\) of the sampled system gives the same original restriction \(\rho^X_{\theta(n),\theta(k)}\). This proves eventual equality even when \(\theta\) has repeated values. Hence the two classes are inverse isomorphisms, without requiring them to be literal inverse schedules.

Write \(s_X,r_X\) and \(s_Y,r_Y\) for the sampling and recovery isomorphisms of source and target. The asserted Hom map is
\[
 [f]\longmapsto s_Y[f]r_X,
\]
with inverse \([g]\mapsto r_Y[g]s_X\). Associativity and the two inverse equations cancel the inserted pairs, proving both round trips. Applying the realization functor transports these equations to point spaces, where sampling and recovery are the cofinal point equivalences already constructed. No norm invariance or surjectivity of source point projections is needed for this categorical conjugation argument.
\end{proof}
\source{KS64}
This strengthens cofinal invariance from points to a specified category of transformations. It still does not make raw residual norms invariant, identify every global set function with a finite presentation, or preserve the numerical conditioning of a repair. Those are separate questions, addressed by the defect maps in the next sections.

\begin{insight}
For a projection-surjective source, eventual finite presentations correspond exactly to uniformly finite observations at each output level. This is a precise machine-oriented representation theorem. It is not a claim that every represented map is computationally efficient, or that every arbitrary function of an infinite input has such a presentation.
\end{insight}


\chap{Banach repair and constructed split coordinates}\label{sec:split}
\subsection{An actual complete candidate space}
Let \(X_n,Y_n\) be Banach spaces over a complete nontrivially normed field, with continuous linear bonding maps and uniform bounds \(\norm{\rho^X_n}\le B_X\), \(\norm{\rho^Y_n}\le B_Y\). Use the spaces of bounded operator families
\[
 E=\ell^\infty_n\mathcal L(X_n,Y_n),\qquad
 R=\ell^\infty_n\mathcal L(X_{n+1},Y_n).
\]
The source formalization derives the required completeness from the operator spaces and their bounded product. The same differential satisfies
\[
 \norm{\partial f}_\infty\le(B_X+B_Y)\norm f_\infty.
\]
Thus it is a continuous linear map, and evaluation at a fixed level is continuous. These facts connect algebraic compatibility to an analytic limit argument.

\subsection{From a correction operator to a projection}
The next statement is general. Let \(E\) be Banach, let \(d:E\to R\), \(p:E\to S\), and \(A:R\to E\) be continuous linear maps. Suppose \(pA=0\) and, for \(0\le q<1\),
\begin{equation}\label{eq:contraction}
 \norm{df-dA(df)}\le q\norm{df}\qquad(f\in E).
\end{equation}
The hypothesis concerns residuals of actual candidates. It need not hold for every independent vector in \(R\).

\begin{theorem}[Constructed geometric repair]\label{KS19}
Starting from \(f^{(0)}=f\), define
\[
 f^{(m+1)}=f^{(m)}-A(df^{(m)}).
\]
There is a bounded linear projection \(P_A:E\to\ker d\) with
\[
 f^{(m)}\longrightarrow P_Af,\qquad p(P_Af)=p(f),
\]
and
\[
 \norm{f^{(m)}-P_Af}
 \le\frac{\norm A}{1-q}\norm{df^{(m)}}.
\]
It fixes every exact candidate and satisfies
\[
 \norm{P_A}\le1+\frac{\norm A\norm d}{1-q}.
\]
\end{theorem}
\begin{proof}
Put $r_m=df^{(m)}$. Additivity of $d$ gives
\[
 r_{m+1}=r_m-dAr_m.
\]
Each $r_m$ is the residual of an actual candidate, so the stated hypothesis applies at every iteration. Induction yields, for $j\ge0$,
\[
 \norm{r_{m+j}}\le q^j\norm{r_m},\qquad
 \norm{f^{(m+j+1)}-f^{(m+j)}}
 \le\norm A\,q^j\norm{r_m}.
\]
For $n>m$, the finite triangle inequality therefore gives
\[
 \norm{f^{(n)}-f^{(m)}}
 \le\norm A\,\norm{r_m}\sum_{j=0}^{n-m-1}q^j
 \le\frac{\norm A}{1-q}\norm{r_m}.
\]
Since $\norm{r_m}\le q^m\norm{df}$ tends to zero, this is a Cauchy estimate. Completeness of $E$ supplies a limit, denoted $P_Af$. Passing $n\to\infty$ in the same norm inequality proves the required tail estimate.

Continuity of $d$ implies $d(P_Af)=\lim_m r_m=0$. Protection holds before taking a limit: $pA=0$ gives $pf^{(m+1)}=pf^{(m)}$, hence $pf^{(m)}=pf$ for every $m$. Continuity of $p$ gives $p(P_Af)=pf$. This explains why continuity of the observable belongs to this analytic statement.

For linearity, write $L=I-Ad$, so $f^{(m)}=L^mf$. For fixed $f,g$ and scalars $\lambda,\mu$,
\[
 L^m(\lambda f+\mu g)=\lambda L^mf+\mu L^mg.
\]
Continuity of addition and scalar multiplication allows passage to the norm limit and proves linearity of $P_A$. The estimate with $m=0$ gives
\[
 \norm{P_Af}\le\norm f+\frac{\norm A}{1-q}\norm{df}
 \le\left(1+\frac{\norm A\norm d}{1-q}\right)\norm f,
\]
so the linear map is bounded. If $df=0$, every iterate equals $f$, hence $P_Af=f$. Its image is contained in $\ker d$ and it fixes that kernel, proving surjectivity onto its stated image and $P_A^2=P_A$. For $q=0$ the residual vanishes after the first update and the same inequalities remain valid.
\end{proof}
\source{KS19}
This theorem constructs the dynamics from \(A\). It does not construct a suitable \(A\) for every residual problem. That missing input can, however, be derived for a useful class of actual towers.

\subsection{A right inverse from target splittings}
Suppose each target bonding map has a continuous linear section
\[
 s_n:Y_n\to Y_{n+1},\qquad \rho^Y_ns_n=I,\qquad \norm{s_n}\le S.
\]
Assume the propagation gap
\begin{equation}\label{eq:splitgap}
 SB_X<1,\qquad K=\frac{S}{1-SB_X}.
\end{equation}
For \(r\in R\), recursively set
\begin{equation}\label{eq:greenrecurrence}
 (Ar)_0=0,\qquad
 (Ar)_{n+1}=s_n\bigl(r_n+(Ar)_n\rho^X_n\bigr).
\end{equation}

\begin{theorem}[Split tower correction]\label{KS20}
The recurrence defines a continuous linear map \(A:R\to E\) with
\[
 \partial A=I_R,\qquad \norm A\le K.
\]
The operator \(P=I-A\partial\) is an exact projection and
\[
 (Pf)_0=f_0,\qquad
 \norm{f-Pf}_\infty\le K\norm{\partial f}_\infty.
\]
\end{theorem}
\begin{proof}
Let $\theta=SB_X<1$. The number $K=S/(1-\theta)$ is nonnegative and satisfies
\[
 S+SB_XK=K.
\]
Fix $r\in R$ and define its coordinates recursively as in \eqref{eq:greenrecurrence}. The initial coordinate has norm zero. If $\norm{(Ar)_n}\le K\norm r_\infty$, submultiplicativity of operator norms gives
\[
 \begin{aligned}
 \norm{(Ar)_{n+1}}
 &\le\norm{s_n}\bigl(\norm{r_n}+\norm{(Ar)_n}\norm{\rho^X_n}\bigr)\\
 &\le S(1+KB_X)\norm r_\infty=K\norm r_\infty.
 \end{aligned}
\]
Induction proves the estimate at every coordinate, so the recursively defined family actually belongs to the bounded product $E$.

The coordinate recurrence preserves addition and scalar multiplication, starting from the zero initial coordinate. Induction therefore proves these laws for the family map $A$. Its uniform estimate proves $\norm A\le K$, giving a continuous linear map with the claimed domain and codomain.

For every $n$, apply the target restriction to the recurrence. The equation $\rho^Y_ns_n=I$ gives
\[
 \rho^Y_n(Ar)_{n+1}=r_n+(Ar)_n\rho^X_n,
 \qquad (\partial Ar)_n=r_n.
\]
Coordinate equality yields the operator identity $\partial A=I_R$. In particular this identity holds for every bounded residual, not just residuals already known to lie in the image.

Set $P=I-A\partial$. Then
\[
 \partial P=\partial-\partial A\partial=0,\qquad
 P^2=I-2A\partial+A(\partial A)\partial=I-A\partial=P.
\]
The expression $2A\partial$ abbreviates addition of the same operator twice and uses no division by two. Every exact $f$ is fixed because $\partial f=0$. Since $(Ar)_0=0$, the seed of $Pf$ is $f_0$. Finally $f-Pf=A\partial f$, and the bound on $A$ gives the asserted residual estimate. The denominators above are justified exactly by the strict propagation gap.
\end{proof}
\source{KS20}
The surjectivity is onto all bounded residual families. It is stronger than the assertion that some exact output can be assigned to each candidate. The gap and the chosen splittings are genuine inputs.

\subsection{Protected information beyond the seed}
When the target grows, its new kernel directions are independent information. Define
\[
 N_n(f)=(I-s_n\rho^Y_n)f_{n+1},\qquad
 (p_Lf)_0=f_0,\quad(p_Lf)_{n+1}=N_n(f).
\]
Each \(N_n(f)\) lands in \(\ker\rho^Y_n\). The observable \(p_L\) retains both the seed and every new output component, relative to the chosen sections.

\begin{theorem}[Split coordinate theorem]\label{KS21}
The observable is a bounded projection with \(p_LA=0\) and \(p_LP=p_L\). It separates exact candidates. If
\[
 \mathcal S_L=\{s\in E:p_Ls=s\},
\]
then
\[
 \Phi:E\simeq_L\mathcal S_L\times R,
 \qquad \Phi(f)=(p_Lf,\partial f),
 \qquad \Phi^{-1}(s,r)=Ps+Ar.
\]
Repair becomes \((s,r)\mapsto(s,0)\). It is the unique exact repair preserving \(p_Lf\).
\end{theorem}
\begin{proof}
Put \(B_n=I-s_n\rho^Y_n\). Since \(\rho^Y_ns_n=I\), direct multiplication gives
\[
 B_n^2=B_n,\qquad \rho^Y_nB_n=0,\qquad B_ns_n=0.
\]
Moreover, \(B_ny=y\) whenever \(\rho^Y_ny=0\). Thus \(B_n\) really projects onto the target kernel. Its norm is at most \(1+SB_Y\). Coordinate zero of \(p_L\) is the identity, and every later coordinate applies \(B_n\), so \(p_L^2=p_L\) and \(\norm{p_L}\le1+SB_Y\). In particular the formula defines a map on the bounded product, not just on arbitrary coordinate families.

The recurrence for \(Ar\) gives zero seed and places \((Ar)_{n+1}\) in the image of \(s_n\). The identities just proved make all its innovations zero. Therefore \(p_LA=0\), and
\[
 p_LP=p_L(I-A\partial)=p_L.
\]
We next prove separation on the exact space. If \(f\) is exact, its restriction equation rewrites as
\[
 f_{n+1}=s_n\rho^Y_nf_{n+1}+B_nf_{n+1}
          =s_nf_n\rho^X_n+N_n(f).
\]
Two exact families with the same seed and innovations agree first at level zero and then at every later level by this recurrence. More generally, if two arbitrary families have the same observable and differential, their difference is exact with zero observable, so separation on the exact space makes their difference zero.

The proposed chart has the claimed target because \(p_L(p_Lf)=p_Lf\). Given \(s\in\mathcal S_L\) and \(r\in R\), set \(f=Ps+Ar\). The preceding identities and \(\partial A=I\) give
\[
 p_Lf=p_Ls=s,\qquad \partial f=r.
\]
This proves one round trip. For the other, \(f\) and \(P(p_Lf)+A\partial f\) have the same observable and residual, so the joint separation just proved identifies them. All maps are linear and bounded on their indicated subspaces and finite products, establishing a continuous linear equivalence.

Finally, \(Pf\) has coordinates \((p_Lf,0)\). Any exact answer with the same observable has those same two coordinates, hence equals \(Pf\). This uniqueness is relative to the complete observable of seed and innovations. It is not obtained by assuming that the seed alone determines an exact growing-target family.
\end{proof}
\source{KS21}
This is an actual analytic coordinate equivalence, not generally an isometry. The protected coordinate depends on the declared splitting. It should not be called an intrinsic semantic coordinate without a further comparison theorem.

\plate{04_split}{In the split coordinate chart, a candidate separates into retained information and compatibility residual. The exact subspace is the zero residual slice. The chart and its protected coordinate depend on the chosen target sections.}{fig:split}

\subsection{A growing example and the gap boundary}
Use \(X_n=Y_n=\R^{n+1}\) with maximum norms. Let \(\pi_n\) delete the last coordinate, take \(\rho^X_n=\tfrac12\pi_n\), \(\rho^Y_n=\pi_n\), and let \(s_n\) append zero. Then \(S=B_Y=1\), \(B_X=1/2\), and \(K=2\).

\begin{example}[Repair of the identity family]\label{KS22}
The bounded candidate \(f_n=I\) is not exact. Its repair is
\[
 (Pf)_n=\operatorname{diag}(2^{-n},2^{-(n-1)},\ldots,2^{-1},1).
\]
The new last output row is retained at every stage while inherited rows are propagated by the factor \(1/2\).
\end{example}
\begin{proof}
At level zero the value is one. Exactness forces every inherited row at the next level to be one half of its previous value. Preservation of the innovation forces the new row to remain the corresponding identity row. Induction gives the displayed diagonal.
\end{proof}
\source{KS22}
For constant identity towers, by contrast, \((\partial f)_n=f_{n+1}-f_n\). The constant residual \(r_n=1\) has no bounded primitive, since any primitive satisfies \(f_n=f_0+n\). Here \(SB_X=1\). This counterexample excludes discarding the gap indiscriminately. It does not prove that the gap is necessary for every individual tower. Exact repair preserving the initial coordinate still exists in this constant example, by replacing a candidate with its initial constant family. Surjectivity onto all bounded residuals does not follow from that repair.
\source{KS23}

\subsection{What is canonical}
For fixed splitting and protected observable, the repair above is unique. Changing the splitting can change both the observable and the repaired answer. In the growing example, replace \(s_n(y)=(y,0)\) by \(s_n'(y)=(y,y_0)\). At the first positive level, repair of the identity becomes
\[
 \begin{pmatrix}1/2&0\\-1/2&1\end{pmatrix}
 \quad\text{instead of}\quad
 \begin{pmatrix}1/2&0\\0&1\end{pmatrix}.
\]
Both are exact continuations for their respective innovation profiles and have the same seed. Thus a unique answer for a declared preservation problem is not an answer independent of all preservation choices.
\source{KS24}

\chap{Repair under a change of representation}\label{sec:repair-transport}
An exact transformation should carry coherent inputs to coherent outputs. It need not carry the selected repair of an incoherent input to the selected repair of its transformed image. The distinction has an exact formula. In split coordinates, a compatible change of representation is generally block triangular rather than block diagonal. Its off-diagonal term records precisely the discrepancy between the two repair routes.

\subsection{A section is a choice of residual coordinates}
Let \(d:E\to R\) be continuous linear between normed spaces over a nontrivially normed field. Suppose it has a continuous linear right inverse \(A:R\to E\), so \(dA=I_R\). Set
\[
 K=\ker d,\qquad P_A=I_E-Ad.
\]
Then \(dP_A=0\), \(P_A\) fixes \(K\), and \(P_A^2=P_A\). The right inverse chooses where each residual is represented inside the candidate space. Different right inverses can choose different representatives of the same residual.

\begin{theorem}[The split defect chart]\label{KS65}
The maps
\[
 \chi_A:E\simeq_L K\times R,\qquad
 \chi_A(x)=(P_Ax,dx),\qquad
 \chi_A^{-1}(u,r)=u+Ar
\]
are inverse continuous linear equivalences. In these coordinates, \(P_A\) sends \((u,r)\) to \((u,0)\). If \(B\) is another right inverse, then \(A-B\) takes values in \(K\), and
\[
 P_Bx-P_Ax=(A-B)dx.
\]
The projections commute if and only if they are equal, and they are equal if and only if \(A=B\).
\end{theorem}
\begin{proof}
The identity \(dA=I\) gives \(dP_A=d-dAd=0\), so the proposed first chart coordinate actually belongs to \(K\). If \(u\in K\), then \(du=0\) and \(P_Au=u\). Therefore, for \((u,r)\in K\times R\),
\[
 d(u+Ar)=r,\qquad P_A(u+Ar)=u+Ar-A(r)=u.
\]
In the other direction, \(P_Ax+A(dx)=x-Adx+Adx=x\). These are both inverse equations, including the membership condition for the kernel coordinate. All component maps are bounded linear, so they define a continuous linear equivalence on the standard finite product. Applied to \(P_Ax\), the chart gives \((P_Ax,0)\), proving the repair formula.

For a second right inverse \(B\), the equation \(d(A-B)=I-I=0\) places its difference from \(A\) in \(K\). Direct subtraction gives
\[
 P_B-P_A=(I-Bd)-(I-Ad)=(A-B)d.
\]
Both projections have range \(K\) and fix every point of \(K\). Thus applying \(P_A\) after \(P_B\) gives \(P_B\), while applying \(P_B\) after \(P_A\) gives \(P_A\). These two compositions are equal if and only if the projections are equal.

If \(P_A=P_B\), evaluate their difference at \(Ar\). Since \(dAr=r\), the difference formula yields \((A-B)r=0\) for every \(r\), hence \(A=B\). Equality of sections conversely makes their projection formulas identical. The argument uses the right inverse to test every residual. It does not infer equality of sections merely from agreement on some restricted set of observed residuals.
\end{proof}
\source{KS65}
This is a linear splitting, not necessarily an orthogonal one. The exact coordinate \(u\) depends on the chosen section. The residual coordinate \(r=dx\) does not. A specified Hilbert geometry can make one section distinguished through orthogonality, but the algebraic split alone does not do so.

The tower construction is stronger in a different direction. Its recurrence builds a right inverse from the target splittings, and its observable identifies what that repair preserves. The abstract chart here isolates the mechanism shared by those concrete constructions.

\subsection{A change of representation has two components}
Let \(D:F\to T\) be a second continuous linear defect map. A compatible change consists of bounded maps
\[
 S:E\to F,\qquad C:R\to T,\qquad DS=Cd.
\]
The equation is the commuting defect square. It says that transforming a candidate and then measuring its defect agrees with transforming its actual defect. No sections have yet been used.

In particular, \(S\) carries \(\ker d\) into \(\ker D\). It may still have a kernel, lose information, or fail to be onto. Calling it a change of representation does not assert an equivalence unless inverse maps have also been supplied.

Choose sections \(A\) of \(d\) and \(B\) of \(D\). Define
\begin{equation}\label{eq:repair-mismatch}
 M=SA-BC:R\longrightarrow F.
\end{equation}
The commuting square and section identities give \(DM=0\). Thus this mismatch takes residual information into an exact direction of the target.

\begin{theorem}[Triangular transport and repair mismatch]\label{KS66}
In split defect coordinates, the change of representation is
\begin{equation}\label{eq:triangular-chart}
 \chi_BS\chi_A^{-1}(u,r)=(Su+Mr,Cr).
\end{equation}
The two repair routes differ by
\begin{equation}\label{eq:repair-change}
 P_BSx-SP_Ax=M(dx).
\end{equation}
Consequently they agree on exact inputs, and
\[
 \norm{P_BSx-SP_Ax}\le\norm M\norm{dx},
 \qquad \norm M\le\norm S\norm A+\norm B\norm C.
\]
They agree on all inputs if and only if \(SA=BC\).

For two consecutive compatible changes \((S_1,C_1)\), \((S_2,C_2)\), with compatible choices of sections at the three ends, their mismatches satisfy
\[
 M_{21}=S_2M_1+M_2C_1.
\]
If \(dx_j\to0\) in norm for a fixed change and fixed sections, the repair discrepancy tends to zero.
\end{theorem}
\begin{proof}
First \(DS=Cd\) sends \(\ker d\) into \(\ker D\), and
\[
 DM=DSA-DBC=CdA-C=C-C=0.
\]
Thus both terms in the proposed exact coordinate have the correct codomain. Given \(u\in\ker d\), compute the residual and exact coordinates of \(S(u+Ar)\) separately:
\[
 DS(u+Ar)=Cr,\qquad
 S(u+Ar)-BCr=Su+(SA-BC)r.
\]
This proves the triangular formula as an equality in \(\ker D\times T\).

For a general input, use the same commuting square to obtain
\[
 \begin{aligned}
 P_BSx-SP_Ax
 &=Sx-BDSx-Sx+SAdx\\
 &=(SA-BC)dx=M(dx).
 \end{aligned}
\]
Its norm is at most \(\norm M\norm{dx}\). The triangle inequality and submultiplicativity give \(\norm M\le\norm S\norm A+\norm B\norm C\). Exact inputs have \(dx=0\), so the two routes agree there regardless of the section choices. If they agree on \emph{every} input, substitute \(x=Ar\). Then \(dAr=r\) makes \(Mr=0\) for each residual, hence \(SA=BC\). That equality conversely makes the discrepancy vanish everywhere.

For composition, denote the three sections by \(A,B,A_2\). The composite residual map is \(C_2C_1\), so its mismatch is \(S_2S_1A-A_2C_2C_1\). Insert and subtract \(S_2BC_1\):
\[
 \begin{aligned}
 S_2S_1A-A_2C_2C_1
 &=S_2(S_1A-BC_1)+(S_2B-A_2C_2)C_1\\
 &=S_2M_1+M_2C_1.
 \end{aligned}
\]
All terms map the first residual space into the last candidate space, in the written order. Finally, for fixed bounded \(M\), the estimate \(\norm{M(dx_j)}\le\norm M\norm{dx_j}\) proves convergence of the discrepancy when \(dx_j\to0\). Allowing the section or representation norms to grow with \(j\) would require a separate uniform bound.
\end{proof}
\source{KS66}
The off-diagonal term in \eqref{eq:triangular-chart} has a specific meaning. A raw residual can become part of the target's chosen exact coordinate after a change of presentation. Deleting the residual before changing presentation removes that contribution. Deleting it afterwards does not. This is not a failure of exactness preservation. Both routes end in the target exact space, but they can end at different points in it.

\subsection{A two-coordinate model in which the routes differ}
Use \(E=F=\R^2\), write candidates as \((u,r)\), and take
\[
 d(u,r)=D(u,r)=r,\qquad A(r)=B(r)=(0,r).
\]
The exact space is the horizontal axis and both repairs are \(P(u,r)=(u,0)\). For a real parameter \(c\), set
\[
 S_c(u,r)=(u+cr,r),\qquad C=I.
\]
The defect square commutes because the second coordinate is unchanged. The mismatch is \(M_c(r)=(cr,0)\). Hence
\[
 PS_c(u,r)=(u+cr,0),\qquad S_cP(u,r)=(u,0).
\]
With \(c=2\) and input \((0,1)\), one route gives \((2,0)\) and the other gives \((0,0)\). Both have zero residual. A report containing only their final residuals would miss the entire difference.

\plate{21_mismatch}{The shear preserves the residual coordinate but mixes it into the chosen exact coordinate. Repair after the shear reaches \((2,0)\). Repair before the shear reaches \((0,0)\). Their difference is the exact-valued mismatch applied to the original residual. The example uses the ordinary Euclidean drawing, while the algebraic identity does not depend on that drawing.}{fig:mismatch}

Two shears compose by adding their parameters. In the mismatch law this appears as \(M_{c_2+c_1}=S_{c_2}M_{c_1}+M_{c_2}\), and \(S_{c_2}\) fixes the exact horizontal axis. The identity is therefore a simple instance of the same composition rule, not an unrelated coordinate trick.

\subsection{Actual tower resampling gives a defect map}
Return to Banach towers with contractive source and target bonds. Let \(\phi\) be a monotone schedule. Candidate sampling is
\[
 (S_\phi f)_n=f_{\phi(n)}.
\]
The sampled bonding map is the actual composite restriction between the selected levels. The residual map \(C_\phi\) is not pointwise sampling of a scalar error score. It accumulates transported fine defects across each omitted block.

For a block from \(a\) to \(b\), its value on a residual family \(r\) is
\begin{equation}\label{eq:block-defect-sum}
 (C_{a,b}r)=\sum_{k=a}^{b-1}
 \rho^Y_{a,k}\,r_k\,\rho^X_{k+1,b}.
\end{equation}
Every summand maps \(X_b\) into \(Y_a\). These surrounding restrictions are necessary when the spaces change with the level. For an actual differential, the sum telescopes to
\[
 C_{a,b}(\partial f)=\rho^Y_{a,b}f_b-f_a\rho^X_{a,b}.
\]
The dyadic coarsener in Section~\ref{sec:cascade} is the case \(b=a+2\).

\begin{theorem}[Bounded-gap resampling and selected repair]\label{KS67}
Suppose \(\phi(n+1)-\phi(n)\le m\) for a fixed natural number \(m\). Then the actual sampling maps are bounded and satisfy
\[
 \norm{S_\phi}\le1,\qquad \norm{C_\phi}\le m,
 \qquad \partial_\phi S_\phi=C_\phi\partial.
\]
Exact candidates remain exact under every monotone common resampling, even without a gap bound. If \(A\) and \(B\) are bounded right inverses for the fine and sampled differentials, then
\[
 \norm{P_BS_\phi f-S_\phi P_Af}
 \le(\norm A+m\norm B)\norm{\partial f}.
\]
\end{theorem}
\begin{proof}
For a bounded candidate family, \(\sup_n\norm{f_{\phi(n)}}\le\sup_k\norm{f_k}\), proving the sampling bound. For a bounded residual family, each summand in \eqref{eq:block-defect-sum} is a map from \(X_b\) to \(Y_a\) and has norm at most \(\norm r_\infty\), because every composite bond is contractive. A selected block has length at most \(m\), so the triangle inequality bounds its sum by \(m\norm r_\infty\). Taking the supremum over blocks proves that \(C_\phi\) maps into the bounded residual product and has the asserted operator norm. If a block has length zero its sum is zero, in agreement with the identity restriction.

To verify the commuting square, insert
\(r_k=\rho^Y_kf_{k+1}-f_k\rho^X_k\). The \(k\)-th term becomes
\[
 \rho^Y_{a,k+1}f_{k+1}\rho^X_{k+1,b}
 -\rho^Y_{a,k}f_k\rho^X_{k,b}.
\]
Successive internal terms cancel in the finite sum, leaving \(\rho^Y_{a,b}f_b-f_a\rho^X_{a,b}\). This is exactly the sampled differential. No infinite series interchange is involved.

If the fine differential vanishes, the finite telescope is zero on every selected interval. That conclusion remains true for unbounded gap lengths, because each individual interval is finite. What may fail without a gap bound is boundedness of the residual transformation on \emph{arbitrary} bounded residual families, not preservation of exact inputs.

When the two bounded sections exist, Theorem~\ref{KS66} applies to the commuting square just established. Substitute \(\norm{S_\phi}\le1\) and \(\norm{C_\phi}\le m\) in its estimate to get \((\norm A+m\norm B)\norm{\partial f}\). Existence of those sections is an additional hypothesis, not a consequence of the sampling estimate.
\end{proof}
\source{KS67}
The existence of the two right inverses is not inferred merely from contractivity of the bonds. The split-tower theorem provides them under its additional section and propagation-gap conditions. Keeping these hypotheses separate prevents a resampling theorem from being mistaken for an unconditional reconstruction algorithm.

\subsection{Two distinct losses under coarsening}
On scalar identity towers, a fine residual block \((1,-1)\) has zero two-step sum. Thus \(C_\phi r=0\) need not imply \(r=0\), even for a bounded-gap schedule. This is cancellation. Retaining an independent detail coordinate remedies that information loss in a split coarsener, as the next section shows.

Unbounded gaps cause a different issue. With \(\phi(n)=n^2\), the bounded residual family \(r_k=1\) has block sums \(2n+1\). These are unbounded. Therefore block accumulation cannot be a map from every bounded residual family to a bounded sampled residual family in this example. The constant residual is not claimed to be the differential of a bounded candidate on the identity tower. It is a witness about the proposed operator on the entire residual space.

Neither example contradicts cofinal invariance of coherent points or eventual transports. Those theorems concern exact compatibility. A norm on all raw defects, the retained fine detail, and the choice of a repair section are additional structures.

\begin{insight}
A compatible change of representation transports equations correctly. It transports the selected repair correctly only when the chosen residual sections also commute. The mismatch is an exact-valued operator, so it can remain invisible to a residual-only check.
\end{insight}


\chap{Multiscale residuals and the limits of reconstruction}\label{sec:cascade}
\subsection{A split coarsener retains its missing detail}
Let \(C:R\to\overline R\) and \(J:\overline R\to R\) be linear maps with \(CJ=I\). Put \(H=I-JC\). Then
\[
 r=JCr+Hr,\qquad CH=0,\qquad H^2=H.
\]
The two coordinates \((Cr,Hr)\) retain all information. The coarse coordinate alone generally does not.

The actual dyadic tower coarsener has the typed formula
\begin{equation}\label{eq:dyadic}
 (Cr)_n=\rho^Y_{2n}r_{2n+1}+r_{2n}\rho^X_{2n+1}.
\end{equation}
It is obtained by telescoping two adjacent compatibility squares. With target sections, a lift is
\[
 (Jc)_{2n}=0,\qquad (Jc)_{2n+1}=s_{2n}c_n.
\]
The section identity gives \(CJ=I\), so the missing detail is explicitly \(Hr\), not an unspecified reconstruction error. The original normed construction imposes the bounds needed to make these maps bounded on its operator family spaces.
\source{KS25}

\subsection{A whole cascade}
At scale \(j\), let \(C_j:R_j\to R_{j+1}\) and \(J_j:R_{j+1}\to R_j\) be bounded linear maps with \(C_jJ_j=I\). For \(x\in R_0\), define
\[
 x_0=x,\quad x_{j+1}=C_jx_j,\quad
 L_0=I,\quad L_N=J_0\cdots J_{N-1},
\]
\[
 h_j=x_j-J_jx_{j+1},\qquad Q_Nx=L_Nx_N,
 \qquad S_N=I-Q_N.
\]
The map \(Q_N\) restores the surviving coarse remainder to the original space. It is this restored remainder, not just \(x_N\), that controls reconstruction.

\begin{theorem}[Finite reconstruction and nested remainders]\label{KS26}
For every \(N\),
\[
 x=Q_Nx+\sum_{j<N}L_jh_j,
 \qquad S_Nx=\sum_{j<N}L_jh_j.
\]
Each \(Q_N\) is a projection, and for \(N\le M\),
\[
 Q_NQ_M=Q_MQ_N=Q_M.
\]
The vectors with all details zero are exactly the persistent subspace
\[
 \mathcal P_\infty=\{x:Q_Nx=x\text{ for every }N\}.
\]
\end{theorem}
\begin{proof}
Let $D_0=I$ and $D_N=C_{N-1}\cdots C_0:R_0\to R_N$. Thus $x_N=D_Nx$, $Q_N=L_ND_N$, and repeated use of $C_jJ_j=I$ gives
\[
 D_NL_N=I_{R_N}.
\]
This cancellation takes place between adjacent maps with matching domains. No identification of different residual spaces is needed.

Since $L_{j+1}=L_jJ_j$ and $x_{j+1}=C_jx_j$,
\[
 L_jh_j=L_j(x_j-J_jx_{j+1})=Q_jx-Q_{j+1}x.
\]
Summing from $j=0$ to $N-1$ leaves $x-Q_Nx$ because $Q_0=I$. This proves both finite reconstruction identities, including $N=0$, when the detail sum is empty.

Idempotence follows from $D_NL_N=I$:
\[
 Q_N^2=L_N(D_NL_N)D_N=Q_N.
\]
For $M\ge N$, put $C_{N,M}=C_{M-1}\cdots C_N$ and $J_{N,M}=J_N\cdots J_{M-1}$, using identities for $M=N$. Then
\[
 D_M=C_{N,M}D_N,\quad L_M=L_NJ_{N,M},\quad C_{N,M}J_{N,M}=I.
\]
The identities $D_NL_M=J_{N,M}$ and $D_ML_N=C_{N,M}$ now give
\[
 Q_NQ_M=L_NJ_{N,M}D_M=Q_M,\qquad
 Q_MQ_N=L_MC_{N,M}D_N=Q_M.
\]
Both orders have been checked, rather than inferred from one range inclusion.

If all $h_j$ vanish, the finite telescope gives $Q_Nx=x$ for every $N$. Conversely, if every $Q_Nx=x$, then $L_jh_j=Q_jx-Q_{j+1}x=0$. Apply the left inverse $D_j$ of $L_j$ to obtain $h_j=0$. The injectivity of the restoration composite is the step that recovers the original detail from its transported value.
\end{proof}
\source{KS26}

\plate{05_cascade}{A split cascade retains every discarded component as a detail. The finite sum is always exact after the restored coarse remainder is included. An infinite detail expansion requires that remainder to vanish in the topology being used.}{fig:cascade}

\subsection{Convergence has its own theorem}
Call \(x\) norm recoverable when \(Q_Nx\to0\). This is equivalent to \(S_Nx\to x\), but it does not follow merely from vanishing persistent subspace. Let
\[
 \mathcal F=\bigcup_N\ker Q_N
\]
be the finite core.

\begin{theorem}[Density criterion for norm recovery]\label{KS27}
Suppose \(\sup_N\norm{Q_N}\le K<\infty\). Then the norm recoverable vectors are exactly \(\overline{\mathcal F}\). In particular all vectors are norm recoverable if and only if the finite core is dense.
\end{theorem}
\begin{proof}
The nested identities first imply that the kernels increase: if $Q_My=0$ and $N\ge M$, then $Q_Ny=Q_NQ_My=0$. Thus the finite core really is an increasing union of linear subspaces.

Suppose $Q_Nx\to0$. Put $y_N=S_Nx=x-Q_Nx$. Idempotence gives $Q_Ny_N=Q_Nx-Q_N^2x=0$, so $y_N\in\mathcal F$. Since $x-y_N=Q_Nx\to0$, these are finite-core approximants converging to $x$. Hence $x\in\overline{\mathcal F}$.

Conversely, fix $x\in\overline{\mathcal F}$ and $\epsilon>0$. Choose $y\in\mathcal F$ satisfying $\norm{x-y}<\epsilon/(K+1)$. Some $M$ has $Q_My=0$. For every $N\ge M$, nesting and the common operator bound give
\[
 \norm{Q_Nx}=\norm{Q_N(x-y)}\le K\norm{x-y}
 <\epsilon\frac{K}{K+1}\le\epsilon.
\]
If $K=0$, the norm on the left is already zero, so no division by $K$ is required. The approximation $y$ is chosen before the depth $N$, giving precisely the quantifier order in norm convergence. This proves equality of the two sets. Applying it to every $x\in R_0$ gives the density criterion.
\end{proof}
\source{KS27}
There are other sufficient conditions. If the transported details are absolutely summable in the Banach space \(R_0\), their sum exists and the remaining boundary lies in \(\mathcal P_\infty\). If that persistent subspace is zero, the sum equals \(x\). Alternatively,
\[
 \norm{x-S_Nx}\le\left(\prod_{j<N}\norm{J_j}\right)\norm{x_N}
\]
gives a direct sufficient decay test. A decreasing coarse norm by itself is not enough when restoration gains grow.

\subsection{The actual dyadic result}
For the source's dyadic residual hierarchy, the persistent residual is zero. Zero detail at a scale reveals every even coordinate, while an odd coordinate is passed to a smaller index at the next scale. Strong induction on the coordinate index therefore proves that a residual with all details zero is zero.

\begin{theorem}[Dyadic determination and finite coordinate recovery]\label{KS28}
In the actual dyadic system, all detail coordinates jointly determine the residual. Moreover,
\[
 k+1<2^N\quad\Longrightarrow\quad
 (Q_Nr)_k=0,\qquad(S_Nr)_k=r_k.
\]
Thus every fixed residual coordinate is recovered exactly at a sufficiently large finite depth. The statement does not assert convergence in the supremum norm over all coordinates.
\end{theorem}
\begin{proof}
Write \(R_j\) for the residual space at scale \(j\). Its lift has the actual coordinate formula
\[
 (J_jz)_{2m}=0,\qquad (J_jz)_{2m+1}=s^{(j)}_{2m}z_m,
\]
where \(s^{(j)}_{2m}\) is the corresponding target section at that scale. The expression on the right is composition of the section with the residual operator. Its source and target are those of the odd coordinate at scale \(j\).

Set \(L_{j,0}=I_{R_j}\) and \(L_{j,N+1}=J_jL_{j+1,N}\). We prove, simultaneously for all starting scales \(j\), that
\[
 k+1<2^N\quad\Longrightarrow\quad (L_{j,N}z)_k=0
 \qquad(z\in R_{j+N}).
\]
For \(N=0\), the inequality \(k+1<1\) has no solution. Suppose the assertion holds at depth \(N\). At depth \(N+1\), an even coordinate vanishes by the displayed lift formula. For an odd coordinate \(k=2m+1\), the hypothesis gives
\[
 2(m+1)<2^{N+1},\qquad m+1<2^N.
\]
The induction hypothesis at starting scale \(j+1\) makes \((L_{j+1,N}z)_m\) zero. Composing that zero operator with \(s^{(j)}_{2m}\) gives zero at coordinate \(k\). This completes the induction without identifying residual spaces at different scales.

Now \(Q_Nr=L_{0,N}x_N\), so the support assertion applies to the restored remainder. Since \(S_Nr=r-Q_Nr\), it gives the reconstruction equation. Finally, if two residuals have identical details at every scale, the finite telescope makes their \(S_N\) reconstructions equal for every \(N\). At any fixed coordinate choose \(N\) with \(k+1<2^N\). The reconstruction equation then makes their original coordinates equal. Hence the residuals are equal.
\end{proof}
\source{KS28}
The strict inequality matters for recovery of a whole initial block. At its final coordinate the remainder can still be nonzero. There is also a sharper, nonconsecutive description of the support.

\begin{proposition}[The dyadic support pattern]\label{KS55}
For the same lifts,
\[
 2^N\nmid(k+1)\quad\Longrightarrow\quad(Q_Nr)_k=0.
\]
Thus coordinate \(k\) is certainly recovered once \(N>\nu_2(k+1)\), where \(\nu_2(m)\) is the largest power-of-two exponent dividing the positive integer \(m\). The only possible unrecovered coordinates at depth \(N\) are
\[
 2^N-1,\quad 2\cdot2^N-1,\quad3\cdot2^N-1,\quad\ldots.
\]
\end{proposition}
\begin{proof}
Prove the assertion for \(L_{j,N}\) by induction, uniformly in \(j\). At depth zero every positive integer is divisible by one, so there is nothing to prove. At depth \(N+1\), even coordinates vanish. If \(k=2m+1\), then
\[
 2^{N+1}\mid(k+1)\quad\Longleftrightarrow\quad2^N\mid(m+1).
\]
When divisibility fails, the inner coordinate in \(J_jL_{j+1,N}\) vanishes by induction. The formula for the odd lift finishes the proof. Apply the result to \(Q_Nr=L_{0,N}x_N\).
\end{proof}
\source{KS55}
This refinement is formalized by a scale-uniform induction on the actual varying-fiber residual system in \code{DyadicSupport.lean}. It asserts where a remainder \emph{may} survive, not that every allowed coordinate is nonzero. In the scalar identity-bonding model with \(C(a,b)=a+b\), the residual constantly equal to one has value \(2^N\) at every surviving position. That example proves sharpness of the uniform prefix cutoff.

\plate{16_dyadic_support}{The possible support of the restored dyadic remainder is sparse and periodic. At depth \(N\), only positions whose index plus one is divisible by \(2^N\) can remain. Open cells are guaranteed recovered. Marked cells are allowed to survive, not asserted nonzero for every residual.}{fig:dyadicsupport}

\subsection{Two boundary examples}
On \(\ell^\infty(\N,\R)\), let \(C\) remove the first coordinate and let \(J\) prepend zero. Then \(Q_Nx\) is the tail of \(x\), placed at its original coordinates with zeros before \(N\). Every fixed coordinate is eventually zero in the remainder, and \(\norm{Q_N}\le1\). Nevertheless the constant sequence \(\mathbf1\) satisfies \(\norm{Q_N\mathbf1}=1\) for every \(N\).

\begin{proposition}[The recovery domain is \(c_0\)]\label{KS29}
For this shift cascade,
\[
 S_Nx\to x\text{ in }\ell^\infty
 \quad\Longleftrightarrow\quad x_n\to0.
\]
For \(x_n=q^n\), \(0\le q<1\), the error is exactly \(q^N\).
\end{proposition}
\begin{proof}
The error norm is \(\sup_{n\ge N}|x_n|\). It tends to zero precisely when the sequence tends to zero. For a geometric sequence its largest tail value is the first one, \(q^N\).
\end{proof}
\source{KS29}
As a second example, take scalar spaces with \(C=\tfrac12I\), \(J=2I\). The coarse value \(x_N=2^{-N}x\) tends to zero, but restoration gives \(Q_Nx=x\). Every detail is zero and the entire input is persistent. The example belongs to this scalar cascade, not to the actual dyadic hierarchy, whose persistent subspace was proved to vanish.

\begin{insight}
Retained packets make omission explicit. They do not guarantee convergence. Coordinate recovery, bounded reconstruction, summability, and attainment remain different mathematical properties.
\end{insight}

\input{chapters/04_symmetry}
\input{chapters/04a_boundary_lab}
\input{chapters/04b_phase}
\addtocontents{toc}{\protect\newpage}
\partnote{The geometry of exact realization}{A categorical space of realizations does not come with a numerical distance by default. In a declared linear model, the compatibility equations become a bounded operator. Closed range, holonomy, and conditioning then determine what small residuals actually guarantee.}
\section{Graph views, holonomy, and reconstruction}\label{sec:graph}
\subsection{Actual transports on a finite graph}
Let a finite connected simple graph have vertex set \(V\), root \(o\), and a chosen route \(p_i:o\to i\) for every vertex, with the route at \(o\) empty. Use a common real normed fiber \(H\). An inverse consistent connection assigns continuous linear equivalences \(\tau_{ij}\), with \(\tau_{ji}=\tau_{ij}^{-1}\). Our convention is that \(\tau_{ij}\) carries the value at \(j\) into the coordinates at \(i\).

The candidate and residual spaces are finite products with maximum norms:
\[
 E=\prod_{i\in V}H,\qquad
 R=\prod_{(i,j)\in\vec E}H,\qquad
 (dx)_{ij}=\tau_{ij}x_j-x_i.
\]
Both orientations of an edge are retained. If \(M\) bounds the norms of all measured edge maps, then \(\norm d\le M+1\). The empty edge case has the usual zero residual.

Distinct walks are retained as arrows of the view category. The product along a walk \(q=(i_0,\ldots,i_m)\) is
\[
 T_q=\tau_{i_0i_1}\cdots\tau_{i_{m-1}i_m}.
\]
Concatenation gives the restriction composition law. An exact vertex family defines a coherent walk-view section, and a coherent walk-view section satisfies every edge equation. Thus this graph construction returns to the original ambient realization type rather than creating a disconnected numerical model.

\subsection{Transported residuals and path gains}
\begin{proposition}[The path telescope]\label{KS40}
For a walk \(q=(i_0,\ldots,i_m)\),
\[
 T_qx_{i_m}-x_{i_0}
 =\sum_{k=0}^{m-1}
 (\tau_{i_0i_1}\cdots\tau_{i_{k-1}i_k})(dx)_{i_ki_{k+1}}.
\]
The empty prefix is the identity. If \(\gamma(q)\) is the sum of the products of preceding edge norms, then
\[
 \norm{T_qx_{i_m}-x_{i_0}}\le\gamma(q)\norm{dx}_\infty.
\]
\end{proposition}
\begin{proof}
Write \(U_0=I\) and \(U_k=\tau_{i_0i_1}\cdots\tau_{i_{k-1}i_k}\) for \(1\le k\le m\). The transport convention makes every \(U_k\) a map from the fiber at \(i_k\) to that at \(i_0\). Since \(U_{k+1}=U_k\tau_{i_ki_{k+1}}\), the \(k\)-th transported defect is
\[
 U_k(dx)_{i_ki_{k+1}}
 =U_{k+1}x_{i_{k+1}}-U_kx_{i_k}.
\]
Summing over \(k=0,\ldots,m-1\) cancels each interior value once with each sign. The remaining terms are \(U_mx_{i_m}-U_0x_{i_0}=T_qx_{i_m}-x_{i_0}\), proving the identity with its actual ordered products. For an empty walk both sides are zero.

The triangle inequality and the operator norm bound now give
\[
 \norm{T_qx_{i_m}-x_{i_0}}
 \le\sum_{k=0}^{m-1}\norm{U_k}\norm{(dx)_{i_ki_{k+1}}}
 \le\left(\sum_{k=0}^{m-1}\prod_{j=0}^{k-1}\norm{\tau_{i_ji_{j+1}}}\right)\norm{dx}_\infty.
\]
The bracket is \(\gamma(q)\), with the empty product equal to one. Repeated vertices or edges do not affect the cancellation. Reordering the edge maps would, in general, invalidate it.
\end{proof}
\source{KS40}
For isometric maps, \(\gamma(q)\le|q|\). For arbitrary bounded equivalences, route length cannot replace the amplification gain. Reversing a route can have a different gain, even though its transport is the inverse map.

\subsection{Holonomy is the remaining exactness condition}
Write \(T_i=T_{p_i}\). Reconstruction from a proposed root value is
\[
 (R_pz)_i=T_i^{-1}z.
\]
For each directed edge define the root loop operator
\[
 K_{ij}=T_i\tau_{ij}T_j^{-1},\qquad
 K=\bigcap_{(i,j)\in\vec E}\ker(K_{ij}-I).
\]
Each \(K_{ij}\) is the transport around an actual rooted loop. No flatness assumption has been made.

\begin{theorem}[Holonomy classification]\label{KS41}
Evaluation at the root and reconstruction give an equivalence
\[
 K\simeq_L\ker d.
\]
A root value extends to an exact vertex family precisely when it lies in \(K\), and that extension is unique. The space \(K\) is independent of the chosen routes. Reconstruction is route independent on \(K\).
\end{theorem}
\begin{proof}
Every \(T_i\) is a continuous linear equivalence, being a finite ordered composition of such maps. Reconstruction is therefore well defined. For an edge \((i,j)\), apply \(T_i\) to its defect:
\[
 T_i\bigl(\tau_{ij}T_j^{-1}z-T_i^{-1}z\bigr)
 =T_i\tau_{ij}T_j^{-1}z-z=K_{ij}z-z.
\]
Because \(T_i\) is injective, that edge defect is zero if and only if the corresponding root loop fixes \(z\). Thus \(dR_pz=0\) holds exactly for \(z\in K\). The root route is empty, so \((R_pz)_o=z\).

Conversely, let \(x\in\ker d\). The path telescope along \(p_i\) has all zero summands, hence \(T_ix_i=x_o\) for every vertex. Applying \(T_i^{-1}\) gives \(x_i=T_i^{-1}x_o\), so \(x=R_p(x_o)\). The first part then places \(x_o\) in \(K\). We have proved both inverse equations and uniqueness of an exact extension of a root value.

The maps are linear. On the finite vertex product,
\[
 \norm{R_pz}_\infty\le\left(\max_i\norm{T_i^{-1}}\right)\norm z,
 \qquad \norm{x_o}\le\norm x_\infty.
\]
These bounds also hold on the subspaces \(K\) and \(\ker d\), proving continuity of both directions. No completeness is needed for these finite bounds.

For route independence, the proved equivalence identifies \(K\) with the image of root evaluation from the intrinsically defined space \(\ker d\). That image does not refer to any chosen routes, so every route system yields the same \(K\). For \(z\in K\), reconstruction using either system is an exact family with root \(z\), and uniqueness makes the families equal. For \(z\notin K\), this uniqueness argument does not assert equality of the two generally inconsistent reconstructed families.
\end{proof}
\source{KS41}
With isometric transports the equivalence is an isometry for the vertex maximum norm, because every reconstructed coordinate has the root norm. With general bounded equivalences, it need not be an isometry.

\plate{09_holonomy}{Tree routes reconstruct from an anchor. The remaining edges close loops and impose fixed-vector equations at that anchor. A nontrivial loop can obstruct an arbitrary seed without preventing every nonzero seed from extending. The fixed space, rather than the identity of the whole loop operator, is the relevant realization condition.}{fig:holonomy}

\subsection{Prescribed observations remain visible}
If a prescribed vertex set \(U\) contains \(o\), a candidate \(x\) has an exact repair preserving \(x|_U\) precisely when
\[
 x_o\in K\quad\text{and}\quad (R_px_o)|_U=x|_U.
\]
Necessity follows from the classification. Sufficiency is the reconstructed family itself. The source identifies this criterion with the earlier obstruction class in \(R/d(\ker p_U)\). A graph seed test is therefore an instance of the same protected correction problem, not a replacement for it.

\chap{Closed range is the stability boundary}\label{sec:closed}
\subsection{Projection exists before a residual estimate}
Let \(E\) be a real Hilbert space and \(R\) a real Banach space. Let \(D:E\to R\) be continuous linear. The kernel \(K=\ker D\) is closed, so its orthogonal projection \(q=P_K\) exists~\cite{mathlib-projection}.

\begin{proposition}[Exact orthogonal repair]\label{KS42}
For every \(x\in E\),
\[
 D(qx)=0,\quad q^2=q,\quad\norm{qx}\le\norm x,
\]
and \(qx\) is the unique closest point to \(x\) in \(K\).
\end{proposition}
\begin{proof}
Boundedness of \(D\) makes \(K=\ker D\) closed. Lemma~\ref{lem:projection-tool} constructs \(q=P_K\) and gives the orthogonal decomposition \(x=qx+(x-qx)\). The first term belongs to \(K\), so \(Dqx=0\). Projection fixes every point of \(K\), giving \(q(qx)=qx\).

Pythagoras applied to that decomposition gives \(\norm x^2=\norm{qx}^2+\norm{x-qx}^2\), hence \(\norm{qx}\le\norm x\). For any \(y\in K\), the vector \(qx-y\) lies in \(K\), whereas \(x-qx\) lies in its orthogonal complement. Consequently
\[
 \norm{x-y}^2=\norm{x-qx}^2+\norm{qx-y}^2\ge\norm{x-qx}^2.
\]
Equality holds only when \(\norm{qx-y}=0\), or \(y=qx\). This proves nearest-point uniqueness without any closedness assumption on \(\ran D\). Closedness of the kernel and of the range serve different purposes.
\end{proof}
\source{KS42}
No closed range assumption on \(D\) is needed for this projection. Such an assumption enters when the size of the residual is asked to control the distance to the exact space.

\begin{theorem}[Closed range and quantitative exactness]\label{KS43}
The following are equivalent:
\begin{enumerate}
\item \(\ran D\) is closed.
\item There is \(C\ge0\) with \(\norm{x-qx}\le C\norm{Dx}\) for all \(x\).
\item There is \(C\ge0\) such that every \(x\) has some exact \(y\) with \(\norm{x-y}\le C\norm{Dx}\).
\end{enumerate}
When the range is closed, the least admissible constant is the operator norm of the inverse of \(D:K^\perp\to\ran D\). In the zero operator case that inverse has zero dimensional domain and its norm is zero.
\end{theorem}
\begin{proof}
Let $K=\ker D$ and $K^\perp$ carry the induced Hilbert norm. Define $D_\perp u=Du$, but regard its target as $\ran D$ with the norm inherited from $R$. It is injective: $Du=0$ puts $u$ in both $K$ and $K^\perp$, so $u=0$. It is surjective onto that target: if $r=Dx$, then $u=x-qx\in K^\perp$ and $Du=Dx=r$. Thus it is a bounded linear bijection before any continuity of its inverse is claimed.

Suppose first that the range is closed. Because $R$ is Banach, $\ran D$ is then Banach. The domain $K^\perp$ is also complete. The bounded inverse theorem, proved in Appendix~\ref{app:proof-tools}, gives a bounded inverse $A=D_\perp^{-1}$. Applying it to $Dx=D(x-qx)$ yields
\[
 x-qx=A(Dx),\qquad \norm{x-qx}\le\norm A\norm{Dx}.
\]
Hence condition 2 holds. If any $C\ge0$ satisfies condition 2, then for every $r\in\ran D$ apply it to $u=Ar\in K^\perp$. Since $qu=0$, it gives $\norm{Ar}\le C\norm r$. Taking the operator norm proves $\norm A\le C$. This establishes the least constant, not just a sufficient one. If $D=0$, both $K^\perp$ and the range are zero spaces and $\norm A=0$.

Now assume condition 2. To prove closedness, take an arbitrary sequence $r_n\in\ran D$ with $r_n\to r$ in $R$. Choose $x_n$ with $Dx_n=r_n$ and put $u_n=x_n-qx_n$. Then $u_n\in K^\perp$ and $Du_n=r_n$. Applying the estimate to $u_n-u_m$ gives
\[
 \norm{u_n-u_m}\le C\norm{r_n-r_m}.
\]
Thus $(u_n)$ is Cauchy in the complete space $K^\perp$. Let its limit be $u$. Continuity gives $Du=\lim_n r_n=r$, so $r\in\ran D$. In a normed space, a set containing the limits of all its convergent sequences is closed. This proves condition 1 without presupposing that the initially unbounded inverse is continuous.

Condition 2 implies condition 3 by choosing $y=qx$. Conversely, if condition 3 supplies an exact $y$ for a given $x$, the nearest-point property of $qx$ gives
\[
 \norm{x-qx}\le\norm{x-y}\le C\norm{Dx}.
\]
The same $C$ works for all $x$, proving condition 2. This last implication does not require the selected answers $y$ to depend linearly or continuously on $x$.
\end{proof}
\source{KS43}
The last equivalence is important. A nonlinear choice of exact answers cannot bypass the failure of a uniform residual bound. Failure of closed range is not merely a defect of one selected projection algorithm.

\subsection{A small residual far from exactness}
On \(\ell^2(\N)\), index the standard basis from one and define \(De_n=n^{-1}e_n\). The map is bounded and injective, so its exact space is zero. Yet
\[
 \norm{e_n}=1,\qquad\norm{De_n}=1/n\longrightarrow0.
\]
No finite \(C\) can bound \(\dist(e_n,\ker D)\) by \(C\norm{De_n}\). The range is not closed. For example finite truncations of \((1/n)_{n\ge1}\) belong to the range and converge in \(\ell^2\), but a preimage of the limit would be the constant sequence, which does not belong to \(\ell^2\).

This printed diagonal model illustrates the criterion. The general unit-witness obstruction is present in the recovered Hilbert source. The particular infinite diagonal construction is not represented as a newly compiled Lean theorem in this publication.

\plate{10_closed_range}{Residual size controls distance to the exact space only when transverse directions have a positive lower bound. The diagonal example sends unit vectors to residuals tending to zero, so no uniform inverse gain exists. Exact projection still exists.}{fig:closed}

\subsection{Norm zero versus algebraic zero}
If an additive correction image is not closed, distance to that image can vanish without actual membership. Thus an infimum of residual error can be zero while the exact protected correction is unattained. The algebraic obstruction quotient and its possible seminorm answer different questions. One must identify a closed image, compactness condition, or other attainment principle before replacing one question by the other.

\subsection{Stable repair with protected data}
The preceding orthogonal projection does not generally preserve a separately declared observable. To join the analytic theory to Section~\ref{sec:protected}, let \(p:E\to S\) be continuous linear into a normed space. Define
\[
 M=\ker p,\qquad B=D|_M:M\to R,\qquad H=M\cap\ker D.
\]
The closed subspace \(M\) inherits its Hilbert structure from \(E\). A candidate \(x\) is \emph{feasible} for protected repair when an exact \(y\) with \(p(y)=p(x)\) exists. The algebraic criterion is still \(Dx\in\ran B\).

\begin{theorem}[Nearest feasible protected repair]\label{KS56}
Every feasible \(x\) has a unique closest exact answer \(y_*\) preserving \(p(x)\). Its correction \(h_*=x-y_*\) is the unique vector in the orthogonal complement of \(\ker B\) inside \(M\) satisfying \(Bh_*=Dx\). For every other feasible answer \(y\),
\[
 \norm{x-y}^2=\norm{x-y_*}^2+\norm{y_*-y}^2.
\]
Moreover, the following conditions are equivalent:
\[
 \ran B\text{ is closed}
 \quad\Longleftrightarrow\quad
 \exists C\ge0\ \forall\text{ feasible }x:
 \norm{x-y_*}\le C\norm{Dx}.
\]
When these conditions hold, the least nonnegative constant is the transverse inverse norm of \(B\), including the zero-range case.
\end{theorem}
\begin{proof}
Feasibility gives some $h_0=x-y_0\in M$ with $Bh_0=Dx$. Conversely, every $h\in M$ satisfying $Bh=Dx$ produces an exact protected answer $y=x-h$. Thus the answer problem and the correction problem are equivalent, with correction set
\[
 \{h\in M:Bh=Dx\}=h_0+\ker B.
\]
This is a nonempty closed affine subset of the Hilbert space $M$. Closedness of $\ker B$ follows from continuity of $B$, independently of whether its image is closed.

Let $k_0=P^M_{\ker B}h_0$ and $h_*=h_0-k_0$. Then $Bh_*=Bh_0=Dx$ and $h_*\perp\ker B$ in $M$. If a different initial correction is chosen, it differs from $h_0$ by an element of $\ker B$. Its orthogonal-complement component is therefore again $h_*$. This proves that the construction is independent of the feasibility witness.

Every other correction can be written uniquely as $h_*+k$ with $k\in\ker B$. Orthogonality gives
\[
 \norm{h_*+k}^2=\norm{h_*}^2+\norm k^2.
\]
Writing $y_*=x-h_*$ and $y=x-h_*-k$, we have $x-y=h_*+k$ and $y_*-y=k$. Substitution proves the stated Pythagorean identity. Its last term vanishes only when $y=y_*$, establishing the unique closest answer.

If $\ran B$ is closed, KS43 applied on $M$ gives a bounded inverse of $B$ on $(\ker B)^{\perp_M}$. The vector $h_*$ is exactly that inverse applied to $Dx$, so its norm obeys the required bound.

For the converse, suppose the bound holds for every feasible raw candidate. Every $x\in M$ is feasible, because $y=0$ has $Dy=0$ and $py=px=0$. Its closest protected exact answer is $P^M_{\ker B}x$. Thus the assumed estimate reduces on $M$ to
\[
 \norm{x-P^M_{\ker B}x}\le C\norm{Bx}\qquad(x\in M).
\]
KS43 forces $\ran B$ to be closed. The same restriction also proves that no smaller nonnegative constant than the transverse inverse norm can work. If the range is zero, feasible candidates are already exact under permitted correction and the least correction is zero, so the optimal gain is zero.
\end{proof}
\source{KS56}
The relevant stability boundary is therefore \(D(\ker p)\), not automatically \(D(E)\). Restricting permitted changes can turn a well-conditioned unconstrained problem into an obstructed or poorly conditioned protected one. The nearest protected answer is generally not \(P_{\ker D}x\).

For \(D(x_1,x_2,x_3)=x_1-x_2\), \(p(x)=(x_1,x_3)\), and \(x=(2,0,5)\), the ordinary orthogonal answer is \((1,1,5)\). The protected answer is \((2,2,5)\). The respective squared errors are two and four. Both have zero residual, but only the latter preserves the declared data. The theorem is formalized in \code{ProtectedHilbertRepair.lean}, using the actual operator on \(\ker p\). The displayed three-coordinate values are an illustrative calculation, not a separate universal algorithm.

\plate{15_protected_geometry}{The same compatibility equation permits different repair problems. Orthogonal repair chooses the nearest point on the diagonal. Preserving the first coordinate forces a vertical correction to a different point. The third coordinate is fixed at five throughout.}{fig:protectedgeometry}

\chap{Holonomy repair and transport amplification}\label{sec:holonomyrepair}
\subsection{An orthogonal anchor, not a global least squares answer}
Return to the graph model with a real Hilbert fiber. The root residual is
\[
 \mathcal H z=(K_{ij}z-z)_{ij},\qquad K=\ker\mathcal H.
\]
Let \(q=P_K\), and define
\[
 Px=R_p(qx_o).
\]
This is always exact, fixes exact families, and is idempotent. It is the unique exact family preserving the feasible component \(qx_o\) of the anchor. Its orthogonality is on the anchor fiber. The vertex and edge norms are still maximum norms.

For example on a two-vertex identity connection, \((0,2)\) with root value zero is repaired to \((0,0)\). The constant family \((1,1)\) is closer in the maximum norm, but it does not preserve the anchor. The operator \(P\) is not being claimed to be a global nearest-point projection in that norm.

\subsection{Every factor in the bound has a meaning}
Define the route gain
\[
 A_p=\max_i\gamma(\overline{p_i}),
\]
using reverse routes so that the error is measured in each destination view. The path telescope gives
\[
 \norm{x-R_p(x_o)}_\infty\le A_p\norm{dx}_\infty.
\]
Let \(N_p\) be the residual normalization \((N_pr)_{ij}=T_ir_{ij}\). The exact identity \(\mathcal H=N_pdR_p\) implies
\[
 \norm{\mathcal H(x_o)}_\infty
 \le\norm{N_p}(\norm d A_p+1)\norm{dx}_\infty.
\]

\begin{theorem}[Amplified holonomy repair]\label{KS44}
If \(\ran\mathcal H\) is closed and \(\kappa\) is its transverse inverse gain, then
\begin{equation}\label{eq:graphgain}
 \norm{x-Px}_\infty\le
 \left[A_p+\norm{R_p}\kappa\norm{N_p}(\norm d A_p+1)\right]
 \norm{dx}_\infty.
\end{equation}
The range condition is necessary for the existence of any uniform bound of distance to an exact family by \(\norm{dx}_\infty\). It is automatic for a fixed finite dimensional fiber, but not uniformly quantitative across all connections.
\end{theorem}
\begin{proof}
First quantify the residual seen at the root. The identity \(\mathcal H=N_pdR_p\), applied to \(x_o\), gives
\[
 \begin{aligned}
 \norm{\mathcal H(x_o)}_\infty
 &=\norm{N_p\bigl(d(R_px_o-x)+dx\bigr)}_\infty\\
 &\le\norm{N_p}\bigl(\norm d\norm{R_px_o-x}_\infty+\norm{dx}_\infty\bigr)\\
 &\le\norm{N_p}(\norm d A_p+1)\norm{dx}_\infty.
 \end{aligned}
\]
Closed range of \(\mathcal H\) lets Theorem~\ref{KS43} control the anchor correction by
\(\norm{x_o-qx_o}\le\kappa\norm{\mathcal H(x_o)}_\infty\). Insert the uncorrected reconstruction between \(x\) and \(Px=R_pqx_o\):
\[
 \norm{x-Px}_\infty
 \le\norm{x-R_px_o}_\infty+\norm{R_p}\norm{x_o-qx_o}.
\]
Substituting the two preceding bounds yields \eqref{eq:graphgain}. The finite maximum-norm residual product is a Banach space because its fiber is Hilbert, so the target completeness needed in Theorem~\ref{KS43} is available.

For necessity, suppose a constant \(C\) satisfies
\(\dist(w,\ker d)\le C\norm{dw}_\infty\) for every vertex family \(w\). Fix \(z\in K^\perp\) and take \(w=R_pz\). Every exact family \(e\) has \(e_o\in K\), and root evaluation has norm at most one. Therefore
\[
 \norm z=\dist(z,K)\le\norm{z-e_o}\le\norm{R_pz-e}_\infty.
\]
Taking the infimum over \(e\in\ker d\) is legitimate without assuming that this global nearest point is attained. It gives
\[
 \norm z\le C\norm{dR_pz}_\infty
 \le C\norm{N_p^{-1}}\norm{\mathcal Hz}_\infty.
\]
Here \(N_p\) is a continuous linear equivalence on the finite edge product because every component \(T_i\) is one. For a general root vector, remove its projection onto \(K\); this leaves the same \(\mathcal H\)-image and puts it in \(K^\perp\). Thus the displayed estimate is a uniform residual bound of the kind in Theorem~\ref{KS43}, which forces \(\ran\mathcal H\) to be closed. A finite-dimensional fiber makes that image finite dimensional and hence closed. It does not give a uniform numerical value of the inverse gain across different connections.
\end{proof}
\source{KS44}
For isometric transports, one can use \(\norm{R_p}\le1\), \(\norm{N_p}\le1\), \(\norm d\le2\), and \(A_p\le r\), where \(r\) is the route radius. The bound reduces to
\[
 \norm{x-Px}_\infty\le[r+\kappa(2r+1)]\norm{dx}_\infty.
\]
The exact repair is route independent for a fixed root and inner product. A chosen upper bound involving route gains need not have the same numerical value for different routes.

\subsection{Flatness does not remove amplification}
Consider a scalar chain with \(n+1\) vertices and coordinates \(e_i(z)=2^iz\). Set \(\tau_{ij}=e_ie_j^{-1}\). This connection is flat, but reverse transport along one edge multiplies norms by two. Let
\[
 x_i=8\cdot2^i-1,\qquad y_i=7\cdot2^i.
\]
The family \(y\) is exact and preserves \(x_o=7\). On a forward edge the residual of \(x\) is \(1/2\), and on its reverse it is \(-1\). Thus
\[
 \norm{dx}_\infty=1,\qquad
 \norm{x-Px}_\infty=2^n-1.
\]
A route length bound of \(n\) fails. This is amplification with trivial holonomy, distinct from poor conditioning of the fixed space of a nearly trivial loop.
\source{KS45}
\plate{17_amplification}{A flat connection can still amplify a fixed residual. Both orientations of every edge are measured, giving residual norm one. Root-preserving repair has error \(2^n-1\), which exceeds the route-length estimate once \(n\ge2\). The vertical scale is logarithmic.}{fig:amplification}

\subsection{Small holonomy can still be poorly conditioned}
In \(\R^2\), the rational rotation
\[
 U_t=\frac1{1+t^2}
 \begin{pmatrix}1-t^2&-2t\\2t&1-t^2\end{pmatrix}
\]
satisfies
\[
 \norm{(U_t-I)z}^2=\frac{4t^2}{1+t^2}\norm z^2.
\]
For \(t\ne0\) its fixed space is zero and its residual gain is \(\sqrt{1+t^2}/(2|t|)\), which diverges as \(t\to0\). At \(t=0\), the entire fiber is fixed. Therefore small loop discrepancy does not supply a uniform stable realization estimate. The fixed directions and their conditioning must be analyzed together.

\chap{Quantitative descent in a common Hilbert model}\label{sec:quantitative}
\subsection{Local exactness and matching are separate operators}
Let \(E\) be a real Hilbert space of raw local candidates. Let \(D:E\to R\) measure local equation defects and \(C:E\to Y\) measure overlap mismatch, with Banach targets. Put
\[
 L=\ker D,\qquad B=C|_L,
 \qquad M=\ker D\cap\ker C.
\]
The exact matched families form \(M\). First project to \(L\), then project inside \(L\) to \(\ker B\). Denote the resulting map by \(P\).

\begin{theorem}[The two-stage repair is the joint projection]\label{KS46}
The operator \(P\) is the orthogonal projection onto \(M\), equivalently onto the kernel of \(\Delta=(D,C)\). If \(u=x-P_Lx\) and \(v=P_Lx-Px\), then
\[
 u\perp v,\qquad\norm{x-Px}^2=\norm u^2+\norm v^2.
\]
\end{theorem}
\begin{proof}
The subspace \(L=\ker D\) is closed, hence Hilbert with the inherited inner product. The restriction \(B=C|_L\) is continuous, so its kernel is a closed subspace of \(L\). Under the inclusion into \(E\), this kernel is exactly \(M=L\cap\ker C\). Thus both projections used to define \(P\) exist.

For an input \(x\), write \(\ell=P_Lx\) and let \(Px\) be the projection of \(\ell\) onto \(\ker B\) inside \(L\). Then \(Px\in M\), \(u=x-\ell\in L^\perp\), and \(v=\ell-Px\in L\) is orthogonal to \(M\). Since \(M\subseteq L\), the vector \(u\) is also orthogonal to \(M\). Their sum \(x-Px=u+v\) is therefore orthogonal to \(M\). A vector in \(M\) with that error is its unique orthogonal projection, so \(P=P_M\).

The two errors themselves are orthogonal because \(u\in L^\perp\) and \(v\in L\). Expanding the squared norm of their sum gives \(\norm{x-Px}^2=\norm u^2+\norm v^2\). Finally \(\ker(D,C)=\ker D\cap\ker C=M\), so this is also the projection onto the kernel of the joint defect. The proof uses a second projection \emph{inside} \(L\). Replacing it by the ambient projection onto \(\ker C\) would invalidate the membership step and need not preserve local exactness.
\end{proof}
\source{KS46}
This identity does not require the local and overlap operators to commute. The second projection is taken in the already exact subspace, not in the whole raw space without restriction.

For a concrete failure of the unrestricted procedure, take \(L=\{(a,0)\}\) and \(\ker C=\{(a,a)\}\) in \(\R^2\). Their intersection is zero. Starting from \((1,0)\), projection to \(L\) does nothing, while subsequent projection to \(\ker C\) gives \((1/2,1/2)\), which is no longer in \(L\). The restricted second stage instead gives the correct joint answer zero.

\plate{11_descent}{Local correction and overlap correction are orthogonal errors when the second problem is restricted to the locally exact subspace. Their composition is the projection onto the joint exact space. The overlap range inside that subspace is the additional stability obligation.}{fig:quantdescent}

\subsection{A stable repair requires a stable matching problem}
Suppose \(\ran D\) and \(\ran B\) are closed, with residual gains \(a,b\). Then
\[
 \norm u\le a\norm{Dx},\qquad
 \norm v\le b\norm{CP_Lx}
 \le b(\norm{Cx}+\norm C\,a\norm{Dx}).
\]
Therefore
\begin{equation}\label{eq:jointgain}
 \norm{x-Px}\le
 a(1+b\norm C)\norm{Dx}+b\norm{Cx}.
\end{equation}
The Pythagorean identity gives the sharper estimate
\[
 \norm{x-Px}^2\le
 a^2\norm{Dx}^2+b^2\bigl(\norm{Cx}+a\norm C\norm{Dx}\bigr)^2.
\]
For the maximum norm on \(R\times Y\), two sufficient gains are therefore
\[
 G_2=\sqrt{a^2+b^2(1+a\norm C)^2}
 \le G_1=a(1+b\norm C)+b.
\]
The inequality follows from \(\sqrt{u^2+v^2}\le u+v\) for nonnegative \(u,v\). The archived interface supplies the sum-type gain \(G_1\), while its squared-error theorem also yields \(G_2\). Neither is claimed to be the least joint gain. In either bound, \(\norm C\) may be replaced by any certified bound on \(C\) restricted to \((\ker D)^\perp\).

\begin{theorem}[Closed range for the joint defect]\label{KS47}
Assume \(\ran D\) is closed. Then
\[
 \ran(D,C)\text{ is closed}\quad\Longleftrightarrow\quad
 \ran(C|_{\ker D})\text{ is closed}.
\]
Under either condition, a finite joint residual-to-repair bound exists.
\end{theorem}
\begin{proof}
Suppose the range of the joint map is closed. To test closedness of $\ran B$, take a sequence $Bk_n\to y$ with $k_n\in\ker D$. The joint images are
\[
 (Dk_n,Ck_n)=(0,Bk_n)\longrightarrow(0,y).
\]
Closedness supplies $x\in E$ with $(Dx,Cx)=(0,y)$. The first coordinate gives $x\in\ker D$, and the second gives $Bx=y$. Hence the matching range is closed. This direction did not require closedness of $\ran D$.

For the reverse direction, assume both $\ran D$ and $\ran B$ are closed, and take joint images converging to $(r,y)$. Thus $Dx_n\to r$ and $Cx_n\to y$. Decompose each input orthogonally as
\[
 x_n=u_n+k_n,\quad
 u_n\in(\ker D)^\perp,\quad k_n\in\ker D.
\]
Because $r\in\ran D$ and the transverse inverse $A$ of $D$ is bounded by KS43, the first components determine
\[
 u_n=A(Dx_n)\longrightarrow A(r)=u,
 \qquad Du=r.
\]
Continuity of $C$ now gives
\[
 Bk_n=Ck_n=Cx_n-Cu_n\longrightarrow y-Cu.
\]
Closedness of $\ran B$ supplies $k\in\ker D$ with $Bk=y-Cu$. There is no need for the sequence $k_n$ itself to converge. The single vector $u+k$ satisfies
\[
 D(u+k)=r,
 \qquad C(u+k)=Cu+Bk=y.
\]
Thus the joint limit lies in the joint range, proving closedness. The target product is Banach under its maximum norm, so KS43 applies to the joint map and gives a finite repair gain. The earlier two-stage estimate supplies an explicit, though not necessarily least, value.
\end{proof}
\source{KS47}


This Hilbert model is a quantitative counterpart of gluing, but it is not a norm on an arbitrary categorical stack. It starts from a fixed raw candidate space, fixed norms, and actual local and overlap operators. Changing those choices requires its own comparison.

\chap{Controlled motion of exact spaces}\label{sec:motion}
\subsection{A two-sided condition number theorem}
Let \(D,D':E\to R\) be continuous linear maps from a real Hilbert space to the same normed target. Write \(q=P_{\ker D}\), \(q'=P_{\ker D'}\). Suppose
\[
 \norm{(I-q)x}\le C\norm{Dx},\qquad
 \norm{(I-q')x}\le C'\norm{D'x},\qquad C,C'\ge0,
\]
and \(\norm{D-D'}\le\varepsilon\).

\begin{theorem}[Motion of the exact projector]\label{KS48}
Without either kernel inclusion as a hypothesis,
\begin{equation}\label{eq:motion}
 \norm{q-q'}\le\max(C,C')\varepsilon.
\end{equation}
In particular a uniform residual gain \(M\) makes \(D\mapsto P_{\ker D}\) \(M\)-Lipschitz on the class of operators carrying that gain.
\end{theorem}
\begin{proof}
For $z\in\ker D'$, the old residual bound gives
\[
 \norm{(I-q)z}\le C\norm{Dz}
 =C\norm{(D-D')z}\le C\varepsilon\norm z.
\]
Interchanging the roles gives
\[
 \norm{(I-q')w}\le C'\varepsilon\norm w
 \qquad(w\in\ker D).                                      \tag{*}
\]
We need a dual form of (*). For $v\in(\ker D')^\perp$ and any $y\in E$, self-adjointness of $q$ and $v=(I-q')v$ give
\[
 |\langle qv,y\rangle|
 =|\langle v,qy\rangle|
 =|\langle v,(I-q')qy\rangle|
 \le C'\varepsilon\norm v\norm{qy}
 \le C'\varepsilon\norm v\norm y.
\]
The Hilbert norm equals the supremum of these pairings on the unit ball, so
$\norm{qv}\le C'\varepsilon\norm v$. This proves the precise restriction estimate needed on $(I-q')x$, not just a coarser bound in terms of $\norm x$.

Now write
\[
 (q-q')x=q(I-q')x-(I-q)q'x.
\]
The two displayed terms lie in $\ker D$ and $(\ker D)^\perp$, respectively, so they are orthogonal. The input components $(I-q')x$ and $q'x$ are also orthogonal. Applying the dual estimate to the first and the original cross estimate to the second gives
\[
\begin{aligned}
 \norm{(q-q')x}^2
 &=\norm{q(I-q')x}^2+\norm{(I-q)q'x}^2\\
 &\le (C'\varepsilon)^2\norm{(I-q')x}^2
      +(C\varepsilon)^2\norm{q'x}^2\\
 &\le\varepsilon^2\max(C,C')^2
       \bigl(\norm{(I-q')x}^2+\norm{q'x}^2\bigr)\\
 &=\varepsilon^2\max(C,C')^2\norm x^2.
\end{aligned}
\]
The bound $\norm{D-D'}\le\varepsilon$ implies $\varepsilon\ge0$, so taking square roots yields the operator estimate. If the common gain is $M$, substitute $C=C'=M$ and $\varepsilon=\norm{D-D'}$. This is exactly the Lipschitz inequality on the class carrying that gain. The use of two orthogonal decompositions explains the maximum coefficient instead of a sum.
\end{proof}
\source{KS48}
There is no smallness threshold when both certificates are already supplied. The estimate can be uninformative when their gains are large. Knowing only the old certificate is a different situation.

For changing input as well as model, the same theorem gives
\[
 \norm{qx-q'y}\le\norm{x-y}+\max(C,C')\varepsilon\norm y.
\]
It also bounds the change in distance to the exact space by \(\max(C,C')\varepsilon\norm x\).

\subsection{Constructing the new certificate}
Assume the target \(R\) is Banach and \(A:R\to E\) is a continuous linear right inverse of \(D\), with \(DA=I\), \(\norm A\le a\), and \(a\varepsilon<1\). Set
\[
 T=(D-D')A,\qquad Z=(I-T)^{-1},\qquad A'=AZ.
\]
The inverse is supplied by the convergent Neumann series because \(\norm T\le a\varepsilon<1\).

\begin{theorem}[Stable surjectivity and derived repair]\label{KS49}
The perturbed map is surjective, with
\[
 D'A'=I,\qquad \norm{A'}\le\frac{a}{1-a\varepsilon}.
\]
Its exact projection satisfies
\[
 \norm{x-q'x}\le\frac{a}{1-a\varepsilon}\norm{D'x},\qquad
 \norm{q-q'}\le\frac{a\varepsilon}{1-a\varepsilon}.
\]
\end{theorem}
\begin{proof}
The map $T=(D-D')A$ acts on the Banach target $R$ and satisfies $\norm T\le a\varepsilon<1$. The operator series
\[
 Z=\sum_{n=0}^\infty T^n
\]
converges in operator norm and obeys $\norm Z\le(1-a\varepsilon)^{-1}$. Multiplying its finite partial sums on either side by $I-T$ gives $I-T^{N+1}$. Passing to the operator-norm limit proves $(I-T)Z=Z(I-T)=I$. Appendix~\ref{app:proof-tools} supplies the full series argument, including completeness of the operator space.

Since $DA=I$, one has $D'A=DA-(D-D')A=I-T$. Therefore $A'=AZ$ satisfies
\[
 D'A'=(I-T)Z=I,\qquad
 \norm{A'}\le\norm A\norm Z\le\frac a{1-a\varepsilon}.
\]
The first equation explicitly supplies a preimage of every target vector, proving surjectivity.

For any $x$, put $y=x-A'D'x$. Then $D'y=D'x-D'A'D'x=0$, so $y$ is a valid exact competitor for orthogonal projection. Its error gives
\[
 \norm{x-q'x}\le\norm{x-y}=\norm{A'D'x}
 \le\frac a{1-a\varepsilon}\norm{D'x}.
\]
The same argument using $A$ gives $\norm{x-qx}\le a\norm{Dx}$. Both coefficients are nonnegative and $a/(1-a\varepsilon)\ge a$. Substitution into KS48 yields
\[
 \norm{q-q'}\le\frac{a\varepsilon}{1-a\varepsilon}.
\]
No assertion that $I-A'D'$ is orthogonal was needed. It supplies an exact comparison vector, while best approximation supplies the orthogonal bound.
\end{proof}
\source{KS49}
If a measured residual \(y\) satisfies \(\norm{D'x-y}\le\eta\), the distance to the new exact space is at most \(a(\norm y+\eta)/(1-a\varepsilon)\). The uncertainty radius remains a quantity to justify in the application. The theorem does not estimate it statistically.

For a surjective Hilbert-source operator, the right inverse need not be a separate unexplained input. The inverse of \(D:(\ker D)^\perp\to R\), followed by inclusion, constructs it. A known residual certificate bounds its norm. Surjectivity is still substantive, and cannot be assumed for a redundant residual map whose target contains unattainable directions.

\subsection{An equivalence between the moving kernels}
\begin{theorem}[Explicit kernel transport]\label{KS50}
Under the hypotheses of Theorem~\ref{KS49},
\[
 J:\ker D\simeq_L\ker D',\qquad Jx=x-A'D'x,
\]
with inverse \(y\mapsto y-ADy\). Moreover
\[
 \norm{Jx-x}\le\frac{a\varepsilon}{1-a\varepsilon}\norm x.
\]
\end{theorem}
\begin{proof}
Let $R_0=I-AD$ and $R_1=I-A'D'$. From $DA=D'A'=I$ we obtain $DR_0=0$ and $D'R_1=0$. Each $R_i$ fixes its corresponding kernel. Thus $R_1$ restricts to a map from $\ker D$ to $\ker D'$, while $R_0$ restricts in the other direction.

The inverse identities require two additional cancellations. First,
\[
 R_0A'=(I-AD)AZ=(A-ADA)Z=0.
\]
Second, using $D'A=I-T$ and $Z(I-T)=I$,
\[
 R_1A=A-AZD'A=A-AZ(I-T)=0.
\]
For $x\in\ker D$, these identities give
\[
 R_0R_1x=R_0(x-A'D'x)=R_0x=x.
\]
For $y\in\ker D'$, they give
\[
 R_1R_0y=R_1(y-ADy)=R_1y=y.
\]
Both restricted maps are bounded and linear as restrictions of bounded linear operators. They are therefore inverse continuous linear maps, establishing the asserted equivalence.

Finally, if $Dx=0$, then $D'x=(D'-D)x$. Hence
\[
 \norm{Jx-x}=\norm{A'D'x}
 \le\norm{A'}\varepsilon\norm x
 \le\frac{a\varepsilon}{1-a\varepsilon}\norm x.
\]
This derivation uses the selected right inverse $A$ throughout. It proves a controlled transport of exact vectors, not a claim that the selected map is the nearest-point projection onto the new kernel.
\end{proof}
\source{KS50}
The map \(J\) depends on the selected splitting. It is not generally the restriction of the orthogonal projection \(q'\). Equivalence of exact spaces and nearest-point repair answer different questions.

\subsection{An all-real worked model}
Let \(D_t(x,y)=x+ty\). Its kernel is the line \(\{(-tb,b):b\in\R\}\), and its orthogonal projector is
\begin{equation}\label{eq:lineprojector}
 q_t=\frac1{1+t^2}
 \begin{pmatrix}t^2&-t\\-t&1\end{pmatrix}.
\end{equation}
The common right inverse \(r\mapsto(r,0)\) supplies gain at most one for every \(t\), and \(\norm{D_s-D_t}=|s-t|\). The general theorem gives \(\norm{q_s-q_t}\le|s-t|\).

\begin{proposition}[Exact motion of the line]\label{KS51}
For every real \(t\),
\[
 \norm{(q_t-q_0)x}^2=\frac{t^2}{1+t^2}\norm x^2,
 \qquad \norm{q_t-q_0}=\frac{|t|}{\sqrt{1+t^2}}.
\]
In the small perturbation interface, kernel transport sends \((0,b)\) to \((-tb,b)\). Orthogonal projection instead sends it to \((-tb/(1+t^2),b/(1+t^2))\).
\end{proposition}
\begin{proof}
The vector \(v_t=(-t,1)\) spans \(\ker D_t\) and has squared norm \(1+t^2>0\). Projection onto this line is
\[
 x\longmapsto\frac{\langle x,v_t\rangle}{1+t^2}v_t,
\]
which gives the matrix in \eqref{eq:lineprojector}. In particular it is the orthogonal projector, not merely a map whose image is contained in the kernel.

Subtract \(q_0=\operatorname{diag}(0,1)\), obtaining the symmetric matrix
\[
 B_t=q_t-q_0=\frac1{1+t^2}
 \begin{pmatrix}t^2&-t\\-t&-t^2\end{pmatrix}.
\]
The two off-diagonal entries of its square are zero, since \(-t^3+t^3=0\). Both diagonal entries are \((t^4+t^2)/(1+t^2)^2=t^2/(1+t^2)\). Thus
\[
 B_t^{\mathsf T}B_t=B_t^2=\frac{t^2}{1+t^2}I.
\]
Pairing this identity with an arbitrary vector proves the all-vector norm formula. It bounds the operator norm above by \(|t|/\sqrt{1+t^2}\), and evaluating any unit vector gives equality. The argument includes \(t=0\), where both sides vanish.

For the common section \(A(r)=(r,0)\), the perturbation inverse leaves \(A\) unchanged because \(D_tA=I\). The constructed transport is therefore \(Jx=x-AD_tx\) on the old kernel. At \(x=(0,b)\), it gives \((-tb,b)\). Substitution into the projection formula gives instead \((-tb/(1+t^2),b/(1+t^2))\). The strict-radius theorem applies when \(|t|<1\), but these explicit line and projection calculations are valid for every real \(t\).
\end{proof}
\source{KS51}

\plate{12_moving_line}{The exact space can move continuously without remaining equal to the original kernel. Transport along a chosen splitting preserves a coordinate, while orthogonal projection minimizes distance. In the line example these two maps have different formulas even though both land in the same exact space.}{fig:movingline}

\subsection{What fails without the hypotheses}
Take \(D_0=\operatorname{diag}(1,0)\) and \(D_t=\operatorname{diag}(1,t)\). The operator difference has norm \(|t|\). For every nonzero \(t\), however, the new kernel is zero, while the old kernel is the vertical axis. Their projections differ in norm by one. Operator closeness alone does not imply kernel continuity.

More generally, suppose a unit vector \(x\) is old exact but orthogonal to the new kernel. Every new residual gain \(C'\) obeys
\[
 1=\norm{x-q'x}\le C'\norm{D'x}\le C'\varepsilon.
\]
Losing an exact direction under a small operator change forces the new gain to grow. At the strict radius boundary, \(D=I\), \(A=I\), and \(D'=0\) have \(a=\varepsilon=1\), yet \(D'\) is not surjective. Thus the strict condition in Theorem~\ref{KS49} cannot be replaced universally by a non-strict one.
\source{KS52}
\plate{18_rank_loss}{Two families can have the same small operator-change norm and very different exact-space motion. A moving line retains a uniform residual gain. The rank-changing diagonal family loses an exact direction for every nonzero parameter, so its projector gap remains one.}{fig:rankloss}

\subsection{Motion of the actual descent repair}
Compare pairs \((D,C)\) and \((D',C')\) on the same raw Hilbert space. Suppose each pair has closed local range and closed overlap range on its own locally exact subspace. Let \(G,G'\) be the resulting joint residual gains. If
\[
 \norm{D-D'}\le\varepsilon_D,\qquad
 \norm{C-C'}\le\varepsilon_C,
\]
then the maximum norm on the target product gives
\[
 \norm{(D,C)-(D',C')}\le\max(\varepsilon_D,\varepsilon_C).
\]
Applying Theorem~\ref{KS48} to the joint operators, and using the actual two-stage identity of Theorem~\ref{KS46}, gives
\begin{equation}\label{eq:descentmotion}
 \norm{P_{D,C}-P_{D',C'}}
 \le\max(G,G')\max(\varepsilon_D,\varepsilon_C).
\end{equation}
This is a comparison of the original repair maps, not a replacement definition that happens to have a similar kernel. Both the local equations and overlap conditions can change. The raw space and its inner product remain fixed.
\source{KS53}

\input{chapters/05a_certificates}
\input{chapters/05b_transport_cost}
\input{chapters/05c_iteration}
\clearpage
\chap{The residual floor and the closure of admissible changes}\label{sec:floor}
An algebraic obstruction asks whether a residual is exactly the defect of an allowed correction. A numerical floor asks how closely allowed defects can approximate it. These questions agree under a closedness hypothesis. Without that hypothesis they can disagree even in a complete Hilbert setting.

\subsection{Distance to the actual correction image}
Let \(R\) be a normed additive group and \(B\le R\) an additive subgroup. In the protected linear model, \(B=d(\ker p)\). For a residual \(\Omega\in R\), define
\[
 \nu_B(\Omega)=\inf_{b\in B}\norm{\Omega-b}.
\]
The subgroup contains zero, so this is a finite nonnegative real number. It is an infimum, not an assertion that a nearest admissible residual exists.

\begin{proposition}[The floor detects closure]\label{KS76}
For every \(b\in B\),
\[
 0\le\nu_B(\Omega)\le\norm{\Omega-b}.
\]
Moreover,
\[
 \nu_B(\Omega)=0\quad\Longleftrightarrow\quad\Omega\in\overline B.
\]
If \(B\) is closed, this is equivalent to \(\Omega\in B\).
\end{proposition}
\begin{proof}
The set of candidate distances is nonempty because \(0\in B\), bounded below by zero, and contains the finite value \(\norm\Omega\). Its infimum is therefore a finite nonnegative real, and the defining lower-bound property gives \(\nu_B(\Omega)\le\norm{\Omega-b}\) for every \(b\in B\).

Suppose \(\nu_B(\Omega)=0\) and fix \(\varepsilon>0\). There must be \(b\in B\) with \(\norm{\Omega-b}<\varepsilon\). Otherwise every candidate distance would be at least \(\varepsilon\), making \(\varepsilon\) a lower bound and forcing the infimum to be at least \(\varepsilon\), a contradiction. Thus every open ball about \(\Omega\) meets \(B\), which is the closure criterion.

Conversely, if \(\Omega\in\overline B\), every such ball supplies \(b\) with \(\norm{\Omega-b}<\varepsilon\). The infimum then lies between zero and \(\varepsilon\) for every positive \(\varepsilon\), so it is zero. If \(B\) is closed, membership in its closure is exactly membership in \(B\). No step asserts that the infimum has a minimizing element when \(B\) is not closed, and closedness alone is being used here only for the zero-floor membership conclusion.
\end{proof}
\source{KS76}
Translation by a member \(b_0\in B\) does not change the floor: replacing \(b\) by \(b-b_0\) gives the same set of distances. Thus the floor is well defined on the algebraic quotient \(R/B\). When \(B\) is a linear subspace, this is the usual quotient seminorm. A nonzero algebraic class can have seminorm zero when \(B\) is not closed. The retained source implements the floor on residual representatives and proves the closure criterion. This quotient interpretation is its elementary translation-invariance calculation, not a new claim that every obstruction quotient carries a separated norm.

\subsection{A precise three-way distinction in this fixed model}
For this specified subgroup and residual there are exactly three possibilities:
\begin{center}
\begin{tabularx}{\linewidth}{@{}p{.29\linewidth}p{.20\linewidth}X@{}}
\toprule Residual location & Algebraic class & Meaning of the floor\\\midrule
\(\Omega\in B\) & Zero & Exact cancellation exists, and the floor is zero.\\
\(\Omega\in\overline B\setminus B\) & Nonzero & Arbitrarily small corrected residuals exist, but no exact cancellation exists.\\
\(\Omega\notin\overline B\) & Nonzero & The floor is strictly positive and blocks every allowed correction.\\\bottomrule
\end{tabularx}
\end{center}
This is a classification for the stated normed-group correction problem, not a universal trichotomy for every possible Komorebi model. Testing which case holds can itself be difficult. The definitions do not supply a decision algorithm.

In finite-dimensional real or complex linear models, every linear subspace is closed. The middle row is therefore absent. It becomes possible in infinite dimensions, or for additive subgroups that are not closed real-linear subspaces. For example the rational subgroup of \(\R\) is dense, so the distance from any irrational real to that subgroup is zero even though its coset is nonzero.

Closedness also should not be confused with a numerical error bound back in the candidate space. The statement \(\Omega\in B\) supplies an admissible correction, but not a bound on its size. The closed-range estimates of Sections~\ref{sec:closed} and~\ref{sec:motion}, or their protected counterparts, are the additional analytic tools that control that size in their declared settings.

\subsection{A persistent floor blocks every correction sequence}
The correction image can vary with the resolution or the admissibility rule. Let \(B_n\le R\) and let \(\Omega_n\in R\) be given.

\begin{proposition}[Uniformly positive floors obstruct asymptotic repair]\label{KS77}
Suppose \(c>0\) and eventually
\[
 c\le\nu_{B_n}(\Omega_n).
\]
For any choices \(b_n\in B_n\), it is impossible that
\[
 \norm{\Omega_n-b_n}\longrightarrow0.
\]
\end{proposition}
\begin{proof}
Every allowed correction satisfies
\(\norm{\Omega_n-b_n}\ge\nu_{B_n}(\Omega_n)\). The corrected norms are therefore eventually at least \(c\). They cannot converge to zero.
\end{proof}
\source{KS77}
The quantifier covers all admissible sequences, not just the output of one repair routine. On the other hand, the hypothesis is an eventually \emph{uniform} positive lower bound. Having a nonzero obstruction class at each stage is not the same statement. The classes might have zero floors, or their positive floors might tend to zero.

A finite-dimensional illustration uses \(R=\R^2\), \(B=\R\times\{0\}\), and \(\Omega_n=(n,1)\). Every correction leaves the second coordinate equal to one, so the floor is one independently of \(n\). Replacing the residuals by \((n,1/(n+1))\) gives nonzero algebraic classes at every stage, but their floors tend to zero. The persistent-floor theorem applies to the first sequence and not the second.

\subsection{The bridge back to exactification}
Theorem~\ref{KS73} and Proposition~\ref{KS76} fit together closely. If corrections produce a norm-convergent, protected path with vanishing residual, the actual quotient class vanishes. If only the corrected residuals tend to zero, all that follows is membership in \(\overline{d(\ker p)}\).

In particular, when the class is nonzero but its floor is zero, \emph{no} protected path with vanishing residual can satisfy all the exactification hypotheses. Something must fail: the path may escape, its movements may not be summable, the candidate space may not be complete, or its residual and semantic maps may not have the required continuity. The next model isolates escape while keeping the spaces Hilbert and both maps continuous.

\input{chapters/05e_escape}
\clearpage
\chap{A bounded repair budget can recover existence}\label{sec:bounded-realization}
The preceding example does not merely lack a convergent iteration. It has no bounded family of increasingly accurate protected corrections. This suggests a converse question. If approximate corrections stay in one bounded region, does an exact answer exist?

For bounded linear equations in a Hilbert space, the answer is yes. This is an existence statement, not a claim that the given sequence converges in norm. A stronger conclusion becomes available when the candidates are the least norm solutions of nested finite constraint systems.

\begin{scopebox}
This chapter and the next three retain the WR01--WR12 identifiers. Every proof is printed. The present chapter remains written-only. The later regularized optimizer and parameter results WR06--WR07 retain their separately verified version 1.6.0 focused Lean counterpart. Fixed-parameter estimates within WR09 are also formalized, but its limit statement is not. Each result carries its own evidence label. Exact finite models have separate executable checks. The Hilbert and regularization background is classical \cite{clason-inverse}. No novelty is claimed for weak compactness or quadratic regularization.
\end{scopebox}

\subsection{The correction space is the protected kernel}
Let $E,Y,Z$ be real Hilbert spaces, let $d:E\to Y$ and $p:E\to Z$ be bounded linear maps, and fix a raw candidate $f\in E$. Write
\[
 H=\ker p,\qquad T=d|_H:H\to Y,\qquad r=df.
\]
The space $H$ is a closed Hilbert subspace with the norm inherited from $E$. An allowed correction is $a\in H$, and its repaired candidate is $f-a$. Thus
\[
 p(f-a)=pf,\qquad d(f-a)=r-Ta.
\]
The protected exact problem is precisely $Ta=r$. Restricting the domain before optimizing is important. A penalty on $\norm{pa}$ would only discourage a change in the protected data, not prohibit it.

The adjoint below always means $T^*:Y\to H$, taken with the inherited Hilbert metric on $H$. In ambient coordinates this is the projection of $d^*$ onto $H$. A coordinate basis of $\ker p$ need not be orthonormal. Its Gram matrix must then be included in any calculation of correction energy.

\subsection{The compactness needed for the existence argument}
\begin{lemma}[Bounded sequences have weakly convergent subsequences]\label{lem:weak-subsequence}
Every bounded sequence $(a_n)$ in a real Hilbert space has a subsequence converging weakly to some $a$. Moreover,
\[
 \norm a\le\liminf_k\norm{a_{n_k}}.
\]
\end{lemma}
\begin{proof}
Choose $M\ge0$ with $\norm{a_n}\le M$. If $M=0$, every term is zero and the conclusion is immediate. Otherwise let $H_0$ be the closed linear span of the countably many vectors $a_n$. Finite rational linear combinations of those vectors form a countable dense set, so $H_0$ is separable. Apply Gram--Schmidt to a countable dense sequence, omitting zero remainders. This gives either a finite orthonormal basis or a countable orthonormal sequence $(e_j)$ whose closed span is $H_0$. The ambient Hilbert space itself need not be separable.

Each scalar sequence $\langle a_n,e_j\rangle$ is bounded by $M$. For the first coordinate choose a convergent subsequence, for the second choose a subsequence of that subsequence, and continue. Taking the $k$th term from the $k$th nested subsequence, with increasing original indices, yields $a_{n_k}$ for which every fixed coordinate converges. Call the coordinate limits $c_j$. In the finite-dimensional case only finitely many successive selections are required.

For each finite $J$, Bessel's inequality and passage to the limit give
\[
 \sum_{j=1}^Jc_j^2
 =\lim_k\sum_{j=1}^J\langle a_{n_k},e_j\rangle^2\le M^2.
\]
The nondecreasing partial sums are therefore bounded. Orthogonality makes the partial vector sums $\sum_{j\le J}c_je_j$ Cauchy, so completeness defines $a\in H_0$ with $\norm a\le M$.

To verify weak convergence against $h\in H_0$, choose a finite basis combination $h_J$ close to $h$. Then
\[
 |\langle a_{n_k}-a,h\rangle|
 \le |\langle a_{n_k}-a,h_J\rangle|
      +2M\norm{h-h_J}.
\]
For a prescribed tolerance, choose $h_J$ first to control the second term and then $k$ to control the first by finitely many coordinate limits. For a general $h$ in the ambient space, write $h=P_{H_0}h+h_\perp$. The orthogonal component pairs to zero with every $a_{n_k}$ and with $a$, reducing to the preceding case. This proves weak convergence against every test vector.

If $a\ne0$, testing against $a/\norm a$ gives
\[
 \norm a=\lim_k\langle a_{n_k},a/\norm a\rangle
 \le\liminf_k\norm{a_{n_k}}.
\]
The zero limit case follows from nonnegativity. Thus the lower-semicontinuity conclusion concerns the actual norms on the selected subsequence, not only the original bound $M$.
\end{proof}

\begin{theorem}[Bounded approximate feasibility]\label{WR01}
For a bounded linear map $T:H\to Y$ between real Hilbert spaces and $r\in Y$, the following are equivalent:
\[
 \begin{split}
 &r\in\ran T,\\
 &\text{there is a bounded sequence }(a_n)\subset H
       \text{ with }\norm{Ta_n-r}\longrightarrow0.
 \end{split}
\]
If all $\norm{a_n}\le M$, an exact solution exists with norm at most $M$.
\end{theorem}
\begin{proof}
If $Ta=r$, the constant sequence $a_n=a$ is bounded and has zero residual. For the converse, let $\norm{a_n}\le M$ and $Ta_n\to r$ in norm. By the preceding lemma there is a subsequence $a_{n_k}$ converging weakly to $a\in H$ with $\norm a\le M$.

For every fixed $y\in Y$, the adjoint identity and weak convergence give
\[
 \langle Ta,y\rangle
 =\langle a,T^*y\rangle
 =\lim_k\langle a_{n_k},T^*y\rangle
 =\lim_k\langle Ta_{n_k},y\rangle.
\]
The last limit equals $\langle r,y\rangle$ because norm convergence implies convergence of each continuous pairing. Thus $\langle Ta-r,y\rangle=0$ for every $y$. Taking $y=Ta-r$ gives $\norm{Ta-r}^2=0$, hence $Ta=r$.

This proof does not claim that $Ta_{n_k}$ converges in norm merely because $a_{n_k}$ converges weakly. Its norm convergence to $r$ was a separate hypothesis. The boundedness of the candidate corrections is used to obtain a weakly convergent subsequence, and their accurate residuals are used to identify its image exactly. The bound $\norm a\le M$ comes from weak lower semicontinuity, so the exact answer stays inside the same closed budget ball.
\end{proof}
\written{WR01}

No closed range assumption appears. The theorem concerns a single datum and a bounded family of its corrections. Closed range is the stronger, uniform statement that every attainable datum admits a correction controlled by one residual gain.

An immediate contrapositive is useful. If $r\in\overline{\ran T}\setminus\ran T$, then every sequence with $Ta_n\to r$ satisfies $\norm{a_n}\to\infty$. Otherwise infinitely many terms would lie in one bounded ball, contradicting the theorem along that subsequence. This recovers the escape mechanism of Section~\ref{sec:escape} without using the diagonal formula.

\subsection{Existence does not make an arbitrary sequence converge strongly}
Take $H=Y=\ell^2$ and $Te_n=e_n/n$. The vectors $a_n=e_n$ have norm one and satisfy $Ta_n\to0$. An exact answer exists, namely zero. Yet $\norm{e_n-e_m}=\sqrt2$ for $n\ne m$, so no subsequence converges in norm.

The conclusion of WR01 is the existence of a weak limit along a subsequence that solves the equation. It is not summable movement, a strong convergence rate, or a selection rule for the original candidates. The minimum correction for the zero datum is zero, which this particular sequence never approaches in norm.

\subsection{Minimum energy finite repairs do converge strongly}
\fixedplate{31_nested_energy}{Minimum energy finite corrections behave differently for two data under the same diagonal operator. Bounded partial energy produces a strong Hilbert realization. Escaping partial energy accompanies a coherent product-space point outside the Hilbert space. The attainable energy lies below the explicit bound two.}{fig:nested-energy}

Suppose now that $A_1\supseteq A_2\supseteq\cdots$ are nonempty closed affine subsets of a real Hilbert correction space $H$. A typical $A_n$ consists of corrections satisfying the first $n$ bounded linear observation equations. The sets are affine because the equations are linear, and closed because the observations are continuous.

Let $a_n$ be the unique point of $A_n$ nearest to zero, and put $m_n=\norm{a_n}$. The Hilbert projection theorem supplies these points. It can also be proved directly: a sequence approaching the minimum norm is Cauchy by the parallelogram identity, its limit belongs to the closed convex set, and the same identity gives uniqueness.

\begin{theorem}[A uniform finite budget yields a strong realization]\label{WR02}
For the nested affine sets above, $m_n$ is nondecreasing, and for $m\ge n$,
\begin{equation}\label{eq:nested-energy}
 \norm{a_m-a_n}^2=m_m^2-m_n^2.
\end{equation}
Their intersection is nonempty if and only if $\sup_n m_n<\infty$. In that case $a_n$ converges in norm to the unique least norm point $a_*\in\bigcap_n A_n$, with
\[
 \norm{a_*-a_n}^2=\norm{a_*}^2-\norm{a_n}^2.
\]
If the intersection is empty, then $m_n\to\infty$.
\end{theorem}
\begin{proof}
The nearest-point theorem for nonempty closed convex subsets of a Hilbert space, proved in Appendix~\ref{app:proof-tools}, supplies $a_n$. Write $A_n=a_n+L_n$, where $L_n$ is its closed direction subspace. For $\ell\in L_n$ and every real $t$, the vector $a_n+t\ell$ belongs to $A_n$. Minimality implies
\[
 0\le\norm{a_n+t\ell}^2-\norm{a_n}^2
 =2t\langle a_n,\ell\rangle+t^2\norm\ell^2.
\]
Allowing both signs of arbitrarily small $t$ forces $\langle a_n,\ell\rangle=0$.

If $m\ge n$, nesting gives $a_m\in A_n$, so $a_m-a_n\in L_n$. Expanding its norm therefore gives
\[
 \norm{a_m}^2
 =\norm{a_n+(a_m-a_n)}^2
 =\norm{a_n}^2+\norm{a_m-a_n}^2.
\]
This is the exact energy identity. It also makes $m_n^2$ nondecreasing.

Suppose these energies are bounded, and let $L=\lim_n m_n^2$. Given $\epsilon>0$, choose $N$ with $L-m_N^2<\epsilon^2$. For $m\ge n\ge N$,
\[
 \norm{a_m-a_n}^2=m_m^2-m_n^2\le L-m_N^2<\epsilon^2.
\]
Thus $(a_n)$ is Cauchy and has a norm limit $a_*$. For each fixed $n$, all terms from that index onward belong to the closed set $A_n$, so $a_*\in A_n$. This places $a_*$ in their entire intersection.

If $a$ lies in the intersection, then $\norm a\ge\norm{a_n}$ for every $n$. Norm continuity gives $\norm a\ge\norm{a_*}$. Uniqueness follows from the nearest-point theorem applied to the intersection, which is closed and convex and now known to be nonempty. Alternatively, the midpoint of two different minimizers would have strictly smaller norm by the parallelogram identity. Taking $m\to\infty$ in the finite energy identity yields
\[
 \norm{a_*-a_n}^2=\norm{a_*}^2-\norm{a_n}^2.
\]
Conversely, any element of the intersection bounds every minimum norm from above. Therefore an empty intersection forces the nondecreasing sequence $m_n$ to be unbounded. A nondecreasing unbounded nonnegative sequence tends to infinity, giving the last conclusion.
\end{proof}
\written{WR02}

The finite answers need not extend one another coordinate by coordinate. Their shared feature is minimum energy in one common ambient Hilbert space. This is the extra geometry that forces their differences to be orthogonal to the earlier minimizers.

For the attainable diagonal datum $r_n=1/n^2$ under $Tu=(u_n/n)$, the first $N$ equations force $a_n=1/n$ for $n\le N$. Their minimum energy completion is zero beyond $N$, and
\[
 m_N^2=\sum_{n=1}^{N}\frac1{n^2}\le2.
\]
These corrections converge strongly to $(1/n)$. For the escaping datum $r_n=1/n$, the corresponding energy is $m_N^2=N$. Every finite equation is again solved, but the budget test now rules out an ambient solution.


\begin{insight}
Finite consistency is not enough by itself. For nested closed affine problems in one Hilbert space, a uniform bound on the minimum correction norms is exactly what upgrades finite feasibility to ambient feasibility. This is a datum-specific test, not a uniform inverse bound for the operator.
\end{insight}

\input{chapters/05g_budget_certificates}
\clearpage
\chap{Regularization selects a protected correction}\label{sec:regularized-repair}
A fixed budget asks how small the residual can become without exceeding an allowed correction size. Quadratic regularization reverses the emphasis. It asks for the best compromise between the residual and the size of the change. The resulting problem has a unique answer even when the exact protected equation has no solution.

That last sentence is both its usefulness and its limitation. The existence of a regularized answer is not a proof of exact realizability. The limiting behavior, rather than the mere existence of the minimizer, supplies the diagnosis.

\subsection{The normal equation is derived in the protected space}
For $\alpha>0$, define on $H=\ker p$
\[
 J_\alpha(a)=\norm{Ta-r}^2+\alpha\norm a^2.
\]
Both norms are the declared Hilbert norms. All minimizers below are corrections in $H$, so preservation is exact for every parameter, not just in the limit. This is ordinary quadratic Tikhonov regularization applied to the actual protected correction operator. The standard regularization background is discussed in \cite{clason-inverse}. The proof here does not require $T$ to be compact.

\begin{theorem}[A unique regularized correction]\label{WR06}
For every $\alpha>0$, there is a unique minimizer
\[
 a_\alpha=(T^*T+\alpha I)^{-1}T^*r.
\]
It satisfies $T^*(r-Ta_\alpha)=\alpha a_\alpha$. For every $a\in H$,
\begin{equation}\label{eq:regularization-energy}
 J_\alpha(a)-J_\alpha(a_\alpha)
 =\norm{T(a-a_\alpha)}^2+\alpha\norm{a-a_\alpha}^2.
\end{equation}
\end{theorem}
\begin{proof}
\emph{Step 1: the normal operator has an actual bounded inverse.} Let $B=T^*T+\alpha I:H\to H$. It is bounded and self-adjoint. For every $h\in H$,
\[
 \langle Bh,h\rangle=\norm{Th}^2+\alpha\norm h^2
 \ge\alpha\norm h^2.
\]
If $h\ne0$, Cauchy--Schwarz and division by $\norm h$ give $\norm{Bh}\ge\alpha\norm h$. The same inequality holds at zero. It implies injectivity of $B$.

To prove closedness of its range, suppose $Bh_n\to y$. Applying the lower bound to $h_n-h_m$ gives
\[
 \norm{h_n-h_m}\le\alpha^{-1}\norm{Bh_n-Bh_m}.
\]
Thus $h_n$ is Cauchy. Completeness of $H$ supplies a limit $h$, and continuity gives $Bh=y$. If $y$ is orthogonal to the range of $B$, then for every $h$,
\[
 0=\langle Bh,y\rangle=\langle h,By\rangle.
\]
Taking $h=By$ gives $By=0$, and injectivity gives $y=0$. A closed subspace with zero orthogonal complement is the whole Hilbert space, so $\ran B=H$.

The inverse therefore exists. Its linearity follows from uniqueness of preimages, and its bound is $\norm{B^{-1}y}\le\alpha^{-1}\norm y$. This construction does not assume closedness of $\ran T$. It uses the positive term $\alpha I$ in the normal operator.

\emph{Step 2: define the optimizer and prove its normal equation.} Set $a_\alpha=B^{-1}T^*r$. Then $Ba_\alpha=T^*r$, or equivalently
\[
 T^*(Ta_\alpha-r)+\alpha a_\alpha=0.
\]
Because the construction takes place in $H=\ker p$, it already satisfies the exact protection requirement.

\emph{Step 3: prove global minimality, rather than just stationarity.} Write any proposed correction as $a=a_\alpha+h$ and set $s=Ta_\alpha-r$. Expanding both squared norms gives
\[
\begin{aligned}
 J_\alpha(a_\alpha+h)-J_\alpha(a_\alpha)
 &=\norm{Th}^2+\alpha\norm h^2
     +2\langle s,Th\rangle+2\alpha\langle a_\alpha,h\rangle\\
 &=\norm{Th}^2+\alpha\norm h^2
     +2\langle T^*s+\alpha a_\alpha,h\rangle\\
 &=\norm{Th}^2+\alpha\norm h^2.
\end{aligned}
\]
This is the asserted energy identity. It is nonnegative and, since $\alpha>0$, vanishes only for $h=0$. Therefore $a_\alpha$ is the unique global minimizer. In the zero correction space, the same construction and identities reduce to its single element, zero.
\end{proof}
\written{WR06}

A declared regularization parameter is part of the problem. A calculation that changes $\alpha$ is solving a different optimization problem, even if the protected data and forward operator stay fixed. On the other hand, the preservation condition itself does not change with $\alpha$.

\subsection{Coordinate bases must retain the correction metric}
In finite dimensions, let $B$ have linearly independent columns spanning $\ker p$. Write $a=Bc$. The coefficient vector $c$ is a representation of the correction, not necessarily an isometric one. Set
\[
 S=dB,\qquad G=B^{\mathsf T}B.
\]
The correct coordinate equation is
\begin{equation}\label{eq:gram-regularization}
 (S^{\mathsf T}S+\alpha G)c=S^{\mathsf T}r.
\end{equation}
The matrix $G$ is positive definite because the columns are independent. Replacing $G$ by the identity changes the meaning of correction size unless the chosen basis is orthonormal.

For example, let $d(a_1,a_2)=a_1$, $p(a_1,a_2)=a_1-a_2$, $r=1$, and $\alpha=1$. A basis of $\ker p$ is the column $B=(1,1)^{\mathsf T}$. Here $S=1$ and $G=2$, so the correction is $(1/3,1/3)$. Penalizing the coordinate with $G=1$ instead gives $(1/2,1/2)$. The latter is not the optimizer for the ambient Euclidean norm. The difference is a metric change, not a discrepancy between equivalent calculations.

The executable example in this edition constructs the nullspace basis over the rationals, retains its Gram matrix, solves \eqref{eq:gram-regularization}, and returns both the correction and an exact normal-equation residual. It also handles a zero-dimensional protected kernel without pretending that a nonzero correction exists.

\subsection{Parameter monotonicity and the fixed-budget problem}
\begin{proposition}[Regularized answers lie on the budget frontier]\label{WR07}
If $0<\alpha<\gamma$, then
\[
 \norm{a_\alpha}\ge\norm{a_\gamma},\qquad
 \norm{r-Ta_\alpha}\le\norm{r-Ta_\gamma}.
\]
At $M_\alpha=\norm{a_\alpha}$,
\[
 \beta(M_\alpha)=\norm{r-Ta_\alpha}.
\]
When $v_\alpha=r-Ta_\alpha\ne0$, its normalization is an optimizing witness in WR04 at that budget.
\end{proposition}
\begin{proof}
Let $A_\alpha=\norm{a_\alpha}^2$ and $R_\alpha=\norm{r-Ta_\alpha}^2$. Comparing the two admissible minimizers in each other's objectives gives
\[
 R_\alpha+\alpha A_\alpha\le R_\gamma+\alpha A_\gamma,
 \qquad R_\gamma+\gamma A_\gamma\le R_\alpha+\gamma A_\alpha.
\]
Add these inequalities and cancel the residual terms. The result is
\[
 (\gamma-\alpha)(A_\alpha-A_\gamma)\ge0.
\]
Since $\gamma>\alpha$, it follows that $A_\alpha\ge A_\gamma$. Return to the first comparison to obtain
\[
 R_\alpha-R_\gamma\le\alpha(A_\gamma-A_\alpha)\le0.
\]
Taking square roots, which preserves order for nonnegative numbers, proves both monotonicity statements.

At $M_\alpha=\norm{a_\alpha}$, the vector $a_\alpha$ is budget-feasible. If another correction $a$ in that ball had strictly smaller residual norm, then
\[
 J_\alpha(a)
 =\norm{r-Ta}^2+\alpha\norm a^2
 <\norm{r-Ta_\alpha}^2+\alpha\norm{a_\alpha}^2,
\]
contradicting its defining minimality. Thus it attains the budget frontier.

For the witness statement, let $v=r-Ta_\alpha\ne0$. The normal equation gives $T^*v=\alpha a_\alpha$. Also $r=v+Ta_\alpha$, hence
\[
 \langle r,v\rangle
 =\norm v^2+\langle a_\alpha,T^*v\rangle
 =\norm v^2+\alpha\norm{a_\alpha}^2.
\]
For $z=v/\norm v$, the support penalty is
\[
 M_\alpha\norm{T^*z}
 =\frac{\alpha\norm{a_\alpha}^2}{\norm v}.
\]
Subtracting it from $\langle r,z\rangle$ yields $\norm v=\beta(M_\alpha)$, so this is an optimizing witness. This calculation also covers a zero correction $a_\alpha$ with nonzero residual. If $v=0$, normalization is unavailable, but the zero witness attains the zero budget residual.
\end{proof}
\written{WR07}

This links a penalized calculation to a hard-budget certificate. It does not say that every larger budget must produce a different answer. Once the least-squares floor is attained, excess budget can be irrelevant. If $T^*r=0$, every regularized correction is zero and the entire profile is constant.

\subsection{The limit separates projection from actual realizability}
Keep $C=\overline{\ran T}$, $b=P_Cr$, and $q=\norm{r-b}$. Since $r-b\in\ker T^*$, the regularized corrections depend on $b$, not on the component that is orthogonal to the closure of all attainable residuals.

\begin{theorem}[The regularization limit alternative]\label{WR08}
As $\alpha\downarrow0$,
\[
 Ta_\alpha\longrightarrow b,\qquad
 \norm{r-Ta_\alpha}\longrightarrow q,\qquad
 \alpha\norm{a_\alpha}^2\longrightarrow0.
\]
If $b\in\ran T$, let $a^\dagger$ be its unique minimum norm preimage. Then
\[
 a_\alpha\longrightarrow a^\dagger\quad\text{in norm}.
\]
If $b\notin\ran T$, then $\norm{a_\alpha}\to\infty$.
\end{theorem}
\begin{proof}
\emph{Step 1: identify the unavoidable component.} Since $r-b\perp\overline{\ran T}$, the adjoint identity gives $T^*(r-b)=0$. Thus $a_\alpha(r)=a_\alpha(b)$. For every correction the orthogonal decomposition is
\[
 \norm{r-Ta}^2=q^2+\norm{b-Ta}^2.
\]
In particular both the residual square and the minimum objective are at least $q^2$.

\emph{Step 2: take the small-parameter limit in the correct order.} Let $\epsilon>0$. Since $b$ lies in the closure of the range, choose a single $a\in H$ with $\norm{r-Ta}^2<q^2+\epsilon$. This $a$ is chosen before $\alpha$. Minimality gives
\[
 q^2\le\norm{r-Ta_\alpha}^2
 \le J_\alpha(a_\alpha)
 \le q^2+\epsilon+\alpha\norm a^2.
\]
For sufficiently small $\alpha$, the last term is below $q^2+2\epsilon$ (and if $a=0$ it is already below $q^2+\epsilon$). Hence both intervening quantities tend to $q^2$. Their difference is $\alpha\norm{a_\alpha}^2$, which tends to zero. The orthogonal decomposition now implies $\norm{Ta_\alpha-b}^2\to0$, and taking square roots gives the first two limits.

\emph{Step 3: when the projected datum is attained, obtain a uniform bound.} Suppose $b\in\ran T$. The preimage set is a nonempty closed affine subset, so projection of zero supplies its unique minimum-norm element $a^\dagger$. Comparing with it and subtracting $q^2$ gives
\[
 \norm{Ta_\alpha-b}^2+\alpha\norm{a_\alpha}^2
 \le\alpha\norm{a^\dagger}^2.
\]
Thus $\norm{a_\alpha}\le\norm{a^\dagger}$ for every $\alpha>0$. This bound does not follow merely from the previously proved vanishing weighted penalty.

\emph{Step 4: strengthen weak convergence to norm convergence.} Take an arbitrary sequence $\alpha_n\downarrow0$, or more generally positive $\alpha_n\to0$. The bound supplies a weakly convergent subsequence $a_{\alpha_{n_k}}\rightharpoonup a$. Since its images converge in norm to $b$, the adjoint pairing argument of WR01 yields $Ta=b$. Weak lower semicontinuity gives
\[
 \norm a\le\liminf_k\norm{a_{\alpha_{n_k}}}
 \le\norm{a^\dagger}.
\]
Minimality and uniqueness of $a^\dagger$ force $a=a^\dagger$. Moreover the same lower bound, together with the uniform upper bound, gives convergence of the norms to $\norm{a^\dagger}$. In a real Hilbert space,
\[
 \norm{a_{\alpha_{n_k}}-a^\dagger}^2
 =\norm{a_{\alpha_{n_k}}}^2+\norm{a^\dagger}^2
       -2\langle a_{\alpha_{n_k}},a^\dagger\rangle\longrightarrow0.
\]
Thus this subsequence converges strongly. If the full positive-parameter family failed to converge strongly, there would be an $\epsilon>0$ and parameters $\alpha_n\to0$ whose corrections all stayed at distance at least $\epsilon$ from $a^\dagger$. The argument just given would produce a strongly convergent subsequence of them, a contradiction. This proves the limit for the entire parameter family.

\emph{Step 5: in the unattained case the norms tend to infinity.} If $b\notin\ran T$ and the norms failed to tend to infinity, there would exist $M<\infty$ and positive $\alpha_n\to0$ with $\norm{a_{\alpha_n}}\le M$. Their images still tend to $b$ by Step 2. WR01 would then put $b$ in the actual range, a contradiction. This proves divergence of the norms, not just unboundedness somewhere along the path.
\end{proof}
\written{WR08}

The vanishing penalty $\alpha\norm{a_\alpha}^2$ is not a boundedness certificate for $a_\alpha$. Its coefficient is tending to zero. In the escaping case, the penalty and residual can both tend to zero while the actual corrections diverge.

\begin{center}
\begin{tabularx}{\linewidth}{@{}p{.17\linewidth}p{.26\linewidth}X@{}}
\toprule Residual floor & Is $b$ attained? & Limit of the protected corrections\\\midrule
$q=0$ & Yes & Strong convergence to the nearest exact protected correction.\\[3pt]
$q=0$ & No & Residual vanishes, correction norms diverge, no exact protected answer.\\[3pt]
$q>0$ & Yes & Strong convergence to the least norm least-squares correction, not an exact answer.\\[3pt]
$q>0$ & No & A positive residual floor and diverging correction norms occur together.\\\bottomrule
\end{tabularx}
\end{center}

The last case is not hypothetical. Let $Tu=(u_n/n)$ on $\ell^2$, define $Su=(Tu,0)$ into $\ell^2\oplus\R$, and take $r=((1/n),1)$. Then $b=((1/n),0)$ is in the closure but not the range, while $q=1$. The same escaping corrections as before coexist with an irreducible orthogonal residual. Positive floor and unbounded correction cost are independent features.

\subsection{Noise is controlled at a fixed parameter}
\fixedplate{34_noise_filter}{For one singular value $\sigma$, the response to residual-data noise is $\sigma/(\sigma^2+\alpha)$. Its maximum occurs at $\sigma=\sqrt\alpha$ and equals $1/(2\sqrt\alpha)$. The plot evaluates this formula and the universal horizontal bounds at three explicit positive parameters.}{fig:noise-filter}

\begin{theorem}[Data sensitivity and a convergent parameter choice]\label{WR09}
For two residual data $r,\widetilde r$ with $\norm{r-\widetilde r}\le\delta$ and the same $T,\alpha$,
\[
 \norm{a_\alpha(r)-a_\alpha(\widetilde r)}
 \le\frac{\delta}{2\sqrt\alpha},\qquad
 \norm{T(a_\alpha(r)-a_\alpha(\widetilde r))}\le\delta.
\]
The coefficient $1/(2\sqrt\alpha)$ is sharp over bounded linear Hilbert operators. If $b=P_Cr\in\ran T$, $\alpha(\delta)\to0$, and $\delta/\sqrt{\alpha(\delta)}\to0$, then the noisy corrections converge strongly to $a^\dagger$ as $\delta\to0$.
\end{theorem}
\begin{proof}
Both corrections are constructed at the same positive parameter and in the same Hilbert correction space. Let $h=a_\alpha(\widetilde r)-a_\alpha(r)$ and $e=\widetilde r-r$. Subtract the two normal equations to obtain
\[
 (T^*T+\alpha I)h=T^*e.
\]
Pair with $h$ and use the adjoint identity:
\[
 \norm{Th}^2+\alpha\norm h^2
 =\langle e,Th\rangle\le\delta\norm{Th}.
\]
Put $t=\norm{Th}\ge0$. If $t>0$, dropping the nonnegative penalty and dividing by $t$ gives $t\le\delta$. If $t=0$, that inequality is already true. Completing the scalar square gives
\[
 \alpha\norm h^2\le\delta t-t^2
 =\frac{\delta^2}{4}-(t-\delta/2)^2
 \le\frac{\delta^2}{4}.
\]
Because $\alpha>0$, taking square roots proves the correction bound. This also explains why the coefficient is $1/(2\sqrt\alpha)$ rather than the coarser $\norm T/\alpha$ obtained by bounding the inverse alone.

For sharpness at a fixed parameter, take $H=Y=\R$ and $T=\sqrt\alpha\,I$. The solution map is multiplication by
\[
 \frac{\sqrt\alpha}{\alpha+\alpha}=\frac1{2\sqrt\alpha}.
\]
Any nonzero data perturbation attains the first bound. This is worst-case sharpness across operators at that parameter, not a claim of equality for every fixed operator and every parameter.

For the limiting statement, insert the exact-data correction at the selected parameter:
\[
 \norm{a_{\alpha(\delta)}(\widetilde r)-a^\dagger}
 \le\frac{\delta}{2\sqrt{\alpha(\delta)}}
       +\norm{a_{\alpha(\delta)}(r)-a^\dagger}.
\]
The first term tends to zero by the noise/parameter hypothesis. The second tends to zero by the attained-projection case of WR08, because the selected parameter tends to zero. Both statements require the same fixed operator and Hilbert metrics. The conclusion is a least-squares correction when the unavoidable residual component is nonzero.
\end{proof}
\written{WR09}

The theorem does not manufacture an exact solution when $r$ is outside the range. For the positive-floor attainable case it converges to a least-squares correction. For the unattainable projected datum it cannot produce a bounded exact-data limit. The range hypothesis must therefore precede the parameter-choice conclusion.


\begin{insight}
Regularization gives a unique, exactly protected answer at every positive parameter. Whether those answers converge to an exact realization, converge only to a best approximation, or escape is determined by the projected datum's actual range membership. A smooth regularized path does not erase that distinction.
\end{insight}

\chap{A complete regularization diagnostic for the escaping model}\label{sec:diagonal-regularization}
The abstract alternatives become explicit for the diagonal model. They also answer a natural objection to the hard prefix corrections: perhaps a smoother correction could fit increasingly well without escaping. It can fit more efficiently, but it cannot avoid the obstruction.

\subsection{The entire regularized family is explicit}
Use $Tu=(u_n/n)$ and $r=w=(1/n)$ as in Section~\ref{sec:escape}. Choose $\alpha=m^{-2}$ for an integer $m\ge1$. The normal equation in each coordinate gives
\[
 a^{(m)}_n=\frac{m^2}{m^2+n^2},\qquad
 v^{(m)}_n=r_n-(Ta^{(m)})_n=\frac{n}{m^2+n^2}.
\]
The repaired ambient candidate is $(-a^{(m)},1)$. Its protected coordinate is one for every $m$.

\begin{proposition}[Energy, residual, and certified tails]\label{WR10}
The displayed $a^{(m)}$ belongs to $\ell^2$ and is the unique regularized correction at $\alpha=m^{-2}$. Set
\[
 A_m=\norm{a^{(m)}}^2=\sum_{n\ge1}\frac{m^4}{(m^2+n^2)^2},\qquad
 B_m=\norm{v^{(m)}}^2=\sum_{n\ge1}\frac{n^2}{(m^2+n^2)^2}.
\]
Then
\begin{equation}\label{eq:soft-bounds}
 \frac m4\le A_m\le\frac{4m}{3},\qquad
 \frac1{4(m+1)}\le B_m\le\frac2m.
\end{equation}
For the sums $A_{m,K},B_{m,K}$ through any integer $K\ge m$,
\begin{equation}\label{eq:soft-tails}
 0\le A_m-A_{m,K}\le\frac{m^4}{3K^3},\qquad
 0\le B_m-B_{m,K}\le\frac1K.
\end{equation}
\end{proposition}
\begin{proof}
For $n>m$, one has $(a_n^{(m)})^2\le m^4/n^4$, so the correction sequence is square-summable. The diagonal operator is bounded with norm one, and $w=(1/n)\in\ell^2$, hence the residual sequence is also in $\ell^2$. The operator is self-adjoint, and the coordinate equality
\[
 \frac1n\frac{n}{m^2+n^2}
 =m^{-2}\frac{m^2}{m^2+n^2}
\]
proves $T^*v^{(m)}=\alpha a^{(m)}$ as an equality of Hilbert vectors. WR06 therefore identifies this actual infinite vector as the unique minimizer, rather than just a coordinatewise candidate.

For $1\le n\le m$, $m^2/(m^2+n^2)\ge1/2$. There are $m$ such coordinates, so $A_m\ge m/4$. For the upper bound, the first $m$ squares are at most one, and the decreasing function $x^{-4}$ gives
\[
 \sum_{n>m}\frac1{n^4}\le\int_m^\infty x^{-4}\,dx=\frac1{3m^3}.
\]
Thus $A_m\le m+m^4/(3m^3)=4m/3$.

For $n\le m$, the residual square is at most $n^2/m^4$. The exact identity
\[
 \sum_{n=1}^m n^2=\frac{m(m+1)(2m+1)}6\le m^3
\]
holds for $m\ge1$, so this part contributes at most $1/m$. For $n>m$, each residual square is at most $1/n^2$, whose tail is at most $1/m$. This proves $B_m\le2/m$.

For the lower bound, $n>m$ implies $m^2+n^2\le2n^2$, hence
\[
 \frac{n^2}{(m^2+n^2)^2}\ge\frac1{4n^2}.
\]
The other direction of the decreasing-function integral comparison gives
\[
 \sum_{n>m}\frac1{n^2}\ge\int_{m+1}^\infty x^{-2}\,dx=\frac1{m+1}.
\]
This proves $B_m\ge1/(4(m+1))$.

Finally, for $K\ge m$, both omitted tails consist of nonnegative terms. Bound the first by $m^4\sum_{n>K}n^{-4}$ and the second by $\sum_{n>K}n^{-2}$. The same upper integral comparisons give the two certified remainders in \eqref{eq:soft-tails}. No finite truncation is substituted for an equality of infinite energies.
\end{proof}
\written{WR10}

Thus the correction norm grows on the order of $\sqrt m$, the residual norm decays on the order of $m^{-1/2}$, and the penalty $\alpha A_m=A_m/m^2$ tends to zero. The smooth correction avoids the hard discontinuity at a prefix boundary, but it still approaches the forbidden constant sequence coordinate by coordinate.

\fixedplate{35_regularized_escape}{At each fixed coordinate, the soft correction approaches one as $m$ tends to infinity. Its transition moves outward as $m$ increases. The chart compares three exact coordinate formulas $m^2/(m^2+n^2)$, not finite approximations to an unspecified optimization procedure. Every displayed sequence has finite energy, but the family's energies are unbounded.}{fig:regularized-escape}

\subsection{The sharp asymptotic correction cost}
The preceding bounds establish divergence and decay. The leading constant can also be computed, which permits a meaningful comparison with hard prefix repair.

\begin{theorem}[A sharp budget asymptotic for the diagonal datum]\label{WR11}
As $m\to\infty$,
\[
 \frac{A_m}{m}\longrightarrow\frac\pi4,\qquad
 mB_m\longrightarrow\frac\pi4.
\]
For the exact budget profile of this infinite problem,
\begin{equation}\label{eq:sharp-budget}
 \lim_{M\to\infty}M\beta_{T,w}(M)=\frac\pi4.
\end{equation}
The hard prefix correction at budget $M=\sqrt N$ instead has
\[
 \sqrt N\,\norm{w-T\mathbf1_N}\longrightarrow1.
\]
\end{theorem}
\begin{proof}
\emph{Step 1: control the infinite Riemann sums uniformly.} Define
\[
 f(x)=\frac1{(1+x^2)^2},\qquad
 g(x)=\frac{x^2}{(1+x^2)^2}.
\]
Then $A_m/m=m^{-1}\sum_{n\ge1}f(n/m)$ and $mB_m=m^{-1}\sum_{n\ge1}g(n/m)$. Fix $L\ge2$ and put $N=\lfloor mL\rfloor$. On $[0,L]$, both functions are uniformly continuous. If $\omega_f(h)$ is a modulus of continuity there, comparison on the $N$ intervals of length $1/m$ gives
\[
 \left|\frac1m\sum_{n=1}^N f(n/m)-\int_0^{N/m}f(x)\,dx\right|
 \le L\omega_f(1/m).
\]
The missing interval $[N/m,L]$ has length below $1/m$ and a bounded integrand, so it contributes a term tending to zero. The same argument applies to $g$. This proves convergence of each truncated sum to its integral on $[0,L]$.

The tails require a bound independent of sufficiently fine discretization. Since $N\ge mL/2$, the inequalities $f(x)\le x^{-4}$ and $g(x)\le x^{-2}$ give
\[
 \frac1m\sum_{n>N}f(n/m)
 \le m^3\sum_{n>N}n^{-4}\le\frac{m^3}{3N^3}
 \le\frac8{3L^3},
\]
\[
 \frac1m\sum_{n>N}g(n/m)
 \le m\sum_{n>N}n^{-2}\le\frac mN\le\frac2L.
\]
The integral tails are at most $1/(3L^3)$ and $1/L$. Given an error tolerance, choose $L$ first to control both discrete and integral tails, then choose $m$ to control the compact-interval sums. This justifies passing from finite-interval Riemann convergence to the full half-line sums. No uncontrolled interchange of two infinite limits is involved.

\emph{Step 2: evaluate the integrals.} The substitution $x=\tan\theta$ gives $dx=\sec^2\theta\,d\theta$ and $1+x^2=\sec^2\theta$, so
\[
 \int_0^\infty f(x)\,dx=\int_0^{\pi/2}\cos^2\theta\,d\theta=\frac\pi4,
 \qquad
 \int_0^\infty g(x)\,dx=\int_0^{\pi/2}\sin^2\theta\,d\theta=\frac\pi4.
\]
Step 1 now proves the two stated asymptotics.

\emph{Step 3: move from regularization parameters to arbitrary budgets.} Set $M_m=\sqrt{A_m}$. Every positive coordinate $m^2/(m^2+n^2)$ increases strictly with $m$, so $A_m$, and hence $M_m$, increases strictly. Since $A_m/m\to c=\pi/4>0$, these budgets tend to infinity and
\[
 \left(\frac{M_{m+1}}{M_m}\right)^2
 =\frac{A_{m+1}/(m+1)}{A_m/m}\frac{m+1}{m}\longrightarrow1.
\]
WR07 identifies $\beta(M_m)=\sqrt{B_m}$. Thus
\[
 M_m\beta(M_m)=\sqrt{(A_m/m)(mB_m)}\longrightarrow c.
\]
For every sufficiently large $M$, choose $m$ with $M_m\le M\le M_{m+1}$. Monotonicity of $\beta$ gives
\[
 M_m\beta(M_{m+1})\le M\beta(M)\le M_{m+1}\beta(M_m).
\]
The lower bound is $(M_m/M_{m+1})[M_{m+1}\beta(M_{m+1})]$ and the upper is $(M_{m+1}/M_m)[M_m\beta(M_m)]$. Both tend to $c$, proving the full budget limit, not merely convergence on a selected subsequence of budgets.

\emph{Step 4: compare the hard-prefix method.} That correction has norm $\sqrt N$ and residual square $\sum_{n>N}n^{-2}$. Integral comparison yields
\[
 \frac{N}{N+1}\le N\sum_{n>N}\frac1{n^2}\le1.
\]
Taking the limit and then a square root gives its leading product one. The two constants are consequently compared at large actual correction budgets, rather than at arbitrarily equated parameter labels.
\end{proof}
\written{WR11}

At large matched correction budgets, the optimal soft repair therefore has about $\pi/4$ times the hard-prefix residual. This improves the efficiency of approximation, not the existence result. Both residuals tend to zero only as the budget tends to infinity. The limiting constant is a calculation for this declared model, not a universal regularization efficiency factor.

\fixedplate{36_budget_product}{The plotted intervals enclose the actual infinite quantities $\sqrt{A_mB_m}$, using finite sums plus the proved tails. The corresponding hard-prefix intervals enclose $\sqrt{N\sum_{n>N}n^{-2}}$ with $N=m$. These have different natural budget sequences. Their limits give the matched-budget asymptotic constants $\pi/4$ and one.}{fig:budget-product}

\subsection{The same operator has an attainable datum}
Change only the residual to $b_n=1/n^2$. Its minimum norm preimage is $a^\dagger_n=1/n\in\ell^2$. At $\alpha=m^{-2}$, the regularized correction is
\[
 c^{(m)}_n=\frac{m^2}{n(m^2+n^2)},\qquad
 a^\dagger_n-c^{(m)}_n=\frac{n}{m^2+n^2}.
\]
The bias square is exactly the previously computed $B_m$. This links the two cases without conflating their existence hypotheses: the escaping residual of one problem is the regularization bias of an attainable problem.

\begin{corollary}[An explicit noisy attainable repair]\label{WR12}
For the attainable datum $b=(1/n^2)$, suppose $\norm{\widetilde b-b}\le\delta$ with $0<\delta\le1$. Choose
\[
 m=\lceil\delta^{-2/3}\rceil,\qquad\alpha=m^{-2}.
\]
Then the noisy regularized correction satisfies
\[
 \norm{a_\alpha(\widetilde b)-a^\dagger}
 \le\sqrt{\frac2m}+\frac{\delta m}{2}
 \le(\sqrt2+1)\delta^{1/3}.
\]
\end{corollary}
\begin{proof}
The exact-data bias has coordinates $n/(m^2+n^2)$, so its squared norm is precisely $B_m$. By WR10,
\[
 \norm{a_{m^{-2}}(b)-a^\dagger}\le\sqrt{2/m}.
\]
WR09 at the same parameter gives
\[
 \norm{a_{m^{-2}}(\widetilde b)-a_{m^{-2}}(b)}
 \le\frac{\delta}{2\sqrt{m^{-2}}}=\frac{\delta m}{2}.
\]
The triangle inequality proves the first bound.

Because $0<\delta\le1$, the number $u=\delta^{-2/3}$ is at least one. The ceiling satisfies $u\le m<u+1\le2u$. Therefore
\[
 \sqrt{2/m}\le\sqrt2\,\delta^{1/3},
 \qquad \frac{\delta m}{2}\le\delta u=\delta^{1/3}.
\]
Adding these estimates proves the displayed rate with every ceiling factor accounted for. The positive-noise condition is needed to define this particular finite integer choice. The statement does not prescribe an infinite $m$ at $\delta=0$, nor does it use a probabilistic assumption on the perturbation.
\end{proof}
\written{WR12}

This is a deterministic guarantee for the specified datum and Hilbert noise bound. It is not a statistical confidence statement. The identical parameter rule in the escaping problem cannot produce convergence to an exact protected Hilbert solution, because no such solution exists.

\subsection{What the calculation permits a machine to certify}
The finite program can check a normal equation, exact preservation, and a quadratic optimality identity. It can return a dual witness ruling out a smaller residual at a given budget. In the infinite diagonal example it can return partial sums together with rational tail enclosures. Those are distinct certificates with distinct conclusions.

None of them alone establishes unrestricted semantic realizability. The bounded-budget theorem needs a genuinely uniform bound. The exactification theorem needs control of candidate movement. The regularization limit needs actual attainment of the projected datum. The repeated appearance of a small residual does not replace any of those hypotheses.

The useful next diagnostic is therefore not just ``Did the residual decrease?'' It is ``What happened to the correction norm, what quantity was preserved, and what existence theorem applies to this datum?''

\chap{Certifying an inexact protected correction}\label{sec:inexact-repair}
An existence theorem for a regularized optimizer does not guarantee that a numerical procedure has reached it. The procedure may stop early, use a poorly scaled basis, or report the residual of the wrong equation. These are mathematical distinctions, not implementation details that can be deferred until after a certificate has been interpreted.

The relevant question is now more specific. Given a proposed correction, what quantities can be checked without first computing the exact optimizer? The answer is a normal-equation residual and, independently, a primal and dual objective gap. Both give quantitative statements about the declared regularized problem. Neither promises that the original unregularized realization problem is solvable.

\begin{scopebox}
The current formal extension constructs the positive-parameter optimizer rather than assuming it exists. It proves the energy identity, global uniqueness, and the estimates in this chapter, with a separate bridge to the actual protected kernel. The weak compactness, infinite-limit, and diagonal asymptotic arguments retain their separately stated verification scopes. The exact Python certificate verifier is tested, not universally verified as a refinement of Lean.
\end{scopebox}

\subsection{There are three different residuals}
Keep $H=\ker p$, $T=d|_H$, and a residual datum $r$. Fix $\alpha>0$. Write $a_\alpha$ for WR06's unique minimizer and let $a\in H$ be any proposed correction. Three quantities now have different meanings:
\[
 v=r-Ta,\qquad
 g=T^*(Ta-r)+\alpha a,\qquad
 e=a-a_\alpha.
\]
The vector $v$ is the fit residual. The vector $g$ is the normal-equation residual, equal to one half of the derivative of the quadratic objective after identifying the dual with $H$. The vector $e$ is the unknown correction error. The exact optimizer generally has $v\ne0$ but always has $g=0$.

The adjoint in $g$ is the restricted adjoint on $H$, not the unrestricted ambient adjoint. Equivalently,
\[
 T^*v=P_Hd^*v.
\]
Thus the correction certificate must measure the part of the normal equation that belongs to the permitted correction space. A numerical program that forgets this projection can incorrectly reject an exact protected optimizer.

\subsection{A posteriori correction and fit bounds}
\begin{theorem}[Normal residual as a solver certificate]\label{WR13}
For every proposed $a\in H$,
\begin{align}
 \norm{a-a_\alpha}&\le\frac{\norm g}{\alpha},\label{eq:solver-error}\\
 \norm{T(a-a_\alpha)}&\le\frac{\norm g}{2\sqrt\alpha},\label{eq:solver-fit}\\
 0\le J_\alpha(a)-J_\alpha(a_\alpha)&\le\frac{\norm g^2}{\alpha}.
 \label{eq:solver-objective}
\end{align}
\end{theorem}
\begin{proof}
Let $e=a-a_\alpha$ and $B=T^*T+\alpha I$. The normal equation of the constructed optimizer gives
\[
 Be=(T^*T+\alpha I)a-T^*r=g.
\]
Pairing with $e$ yields the exact identity
\begin{equation}\label{eq:solver-energy}
 \norm{Te}^2+\alpha\norm e^2=\langle g,e\rangle
 \le\norm g\norm e.
\end{equation}
If $e=0$, all three conclusions hold with zero error. Otherwise drop the first nonnegative term and divide by $\alpha\norm e>0$ to obtain $\norm e\le\norm g/\alpha$.

To obtain the sharper fitted-data bound, retain the quadratic dependence on $\norm e$ instead of substituting the first estimate. Multiplying \eqref{eq:solver-energy} by $\alpha$ and rearranging gives
\[
\begin{aligned}
 \alpha\norm{Te}^2
 &\le\alpha\norm g\norm e-\alpha^2\norm e^2\\
 &=\frac{\norm g^2}{4}
     -\left(\alpha\norm e-\frac{\norm g}{2}\right)^2
 \le\frac{\norm g^2}{4}.
\end{aligned}
\]
Division by $\alpha$ and square roots give the factor $1/(2\sqrt\alpha)$.

By WR06 the objective gap is exactly $\norm{Te}^2+\alpha\norm e^2$, hence by \eqref{eq:solver-energy} it is also $\langle g,e\rangle$. It is nonnegative and satisfies
\[
 J_\alpha(a)-J_\alpha(a_\alpha)
 \le\norm g\norm e\le\frac{\norm g^2}{\alpha}.
\]
Every certificate quantity on the right is computed from the proposal $a$ and the declared problem. The unknown optimizer was used to prove the guarantee, but it is not required to evaluate the certificate.
\end{proof}
\written{WR13}

\paragraph{Sharp scalar examples.} For the first sharpness statement take $T=0$ on $\R$ and a nonzero error. Then $g=\alpha e$. For the second, take $T=\sqrt\alpha\,I$ on $\R$, so $g=2\alpha e$ and $\norm{Te}=\norm g/(2\sqrt\alpha)$. These examples establish worst-case sharpness, not simultaneous equality for every operator.

For a fixed raw vector $f$, the corresponding repaired candidates are $f-a$ and $f-a_\alpha$. Their distance is exactly $\norm e$. Both preserve $pf$ because both corrections belong to $H$. The theorem therefore controls a genuine protected candidate error, not just an intermediate coordinate vector.

\fixedplate{37_solver_amplification}{A fixed normal-equation tolerance has different consequences at different regularization parameters. The curves are universal certified upper bounds for a declared normal-residual norm, not observed errors of a particular experiment. The correction estimate grows as $\alpha^{-1}$ while the fitted-data estimate grows as $\alpha^{-1/2}$.}{fig:solver-amplification}

\subsection{Why a vanishing normal residual is not enough}
Let $T=0$ on $\R$, let $r=0$, and take $\alpha_n=n^{-2}$. The exact optimizer is zero for every $n$. The fixed proposal $a_n=1$ instead has
\[
 v_n=0,\qquad g_n=\alpha_n\longrightarrow0,\qquad
 J_{\alpha_n}(a_n)=\alpha_n\longrightarrow0,
 \qquad\abs{a_n-a_{\alpha_n}}=1.
\]
Thus the fit residual, the normal residual, and even the absolute objective value can all tend to zero without convergence to the selected optimizer. The certificate correctly reports $\abs{g_n}/\alpha_n=1$.

This is not an infinite-dimensional obstruction. It occurs in one dimension because the regularizing penalty itself disappears. The question is not whether the reported residual is small in isolation, but whether it is small relative to the parameter that converts it into an error bound.

\subsection{A dual certificate does not require solving the normal equation}
For an arbitrary $z\in Y$, define
\[
 \mathcal D_\alpha(z)=2\langle r,z\rangle-\norm z^2
                    -\frac{\norm{T^*z}^2}{\alpha}.
\]
Unlike the fixed-budget witness in WR04, this dual witness has no unit-norm constraint. Its quadratic penalties supply the control.

\begin{theorem}[Exact regularized primal and dual gap]\label{WR14}
For every $a\in H$ and $z\in Y$,
\begin{equation}\label{eq:regularized-duality-gap}
 J_\alpha(a)-\mathcal D_\alpha(z)
 =\norm{Ta-r+z}^2+
   \frac{\norm{T^*z-\alpha a}^2}{\alpha}.
\end{equation}
Consequently $\mathcal D_\alpha(z)\le J_\alpha(a)$ for every pair. At
$z_\alpha=r-Ta_\alpha$, equality holds, so the dual maximum is attained and equals the primal minimum. Moreover,
\[
 \alpha\norm{a-a_\alpha}^2
 \le J_\alpha(a)-\mathcal D_\alpha(z).
\]
For the easily computed choice $z=r-Ta$,
\begin{equation}\label{eq:residual-dual-gap}
 J_\alpha(a)-\mathcal D_\alpha(r-Ta)=\frac{\norm g^2}{\alpha}.
\end{equation}
\end{theorem}
\begin{proof}
For arbitrary $a,z$, expand the first square as
\[
 \norm{Ta-r+z}^2
 =\norm{Ta-r}^2+\norm z^2+2\langle Ta,z\rangle-2\langle r,z\rangle,
\]
and the scaled second square as
\[
 \frac{\norm{T^*z-\alpha a}^2}{\alpha}
 =\frac{\norm{T^*z}^2}{\alpha}+\alpha\norm a^2
       -2\langle T^*z,a\rangle.
\]
The adjoint identity cancels $2\langle Ta,z\rangle$ against the last term. The remaining expression is
\[
 \norm{Ta-r}^2+\alpha\norm a^2
 -\left(2\langle r,z\rangle-\norm z^2
                    -\alpha^{-1}\norm{T^*z}^2\right),
\]
which is the asserted gap. Its nonnegativity gives weak duality for every pair, without any constraint on the witness norm.

Set $a=a_\alpha$ and $z_\alpha=r-Ta_\alpha$. The first square is zero by this definition. The normal equation gives $T^*z_\alpha=\alpha a_\alpha$, so the second is zero as well. Therefore
\[
 \mathcal D_\alpha(z_\alpha)=J_\alpha(a_\alpha).
\]
For every other $z$, weak duality with $a_\alpha$ gives $\mathcal D_\alpha(z)\le J_\alpha(a_\alpha)$. This proves attainment of the dual maximum and equality with the primal minimum directly.

For a general proposed pair $(a,z)$, decompose its gap as
\[
 J_\alpha(a)-\mathcal D_\alpha(z)
 =[J_\alpha(a)-J_\alpha(a_\alpha)]
  +[J_\alpha(a_\alpha)-\mathcal D_\alpha(z)].
\]
The second bracket is nonnegative by weak duality, while WR06 bounds the first below by $\alpha\norm{a-a_\alpha}^2$. This proves the correction-error certificate.

Finally take $z=r-Ta$. The first square in the gap is zero, while $T^*z-\alpha a=-g$. The second becomes $\norm g^2/\alpha$, proving the residual-witness formula. A small duality gap controls optimizer error only after division by the specified positive parameter.
\end{proof}
\written{WR14}

A proposed correction and a proposed dual witness can come from different procedures. The verifier only needs the mathematical problem, the two vectors, and the displayed identities. A poor witness remains valid as a lower bound, but can give a loose gap. The residual witness avoids a second optimization routine and recovers WR13's bound exactly.

\subsection{A protected example with every certificate calculated}
Take
\[
 d(x,y)=x,\qquad p(x,y)=x-y,\qquad f=(1,0),\qquad r=df=1,\qquad\alpha=1.
\]
Every allowed correction has the form $(t,t)$. The objective is
\[
 J(t,t)=(t-1)^2+2t^2,
\]
so $a_\alpha=(1/3,1/3)$ and $J(a_\alpha)=2/3$. Suppose a procedure returns only $a=(1/4,1/4)$. Then
\[
 v=\frac34,\qquad T^*v=(3/8,3/8),\qquad
 g=(-1/8,-1/8),\qquad\norm g^2=\frac1{32}.
\]
The actual correction-error square is $1/72$, while the certified upper bound is $1/32$. The actual fitted-error square is $1/144$, while the sharper bound in WR13 is $1/128$.

The proposed objective is $11/16$. Its actual excess over the optimum is $1/48$. The residual witness gives
\[
 \mathcal D_1(3/4)=\frac{21}{32},\qquad
 J(a)-\mathcal D_1(3/4)=\frac1{32}.
\]
The duality gap is an upper bound for suboptimality, not necessarily the exact suboptimality. Both repaired candidates preserve $p(f)=1$.

\fixedplate{38_objective_certificate}{The objective on the actual protected line is $3t^2-2t+1$. The proposed correction at $t=1/4$, exact minimizer at $t=1/3$, and residual-witness lower bound are distinct. The vertical gap to the dual level is the checkable certificate. The smaller gap to the minimum is the true suboptimality.}{fig:objective-certificate}

The supplied program \path{examples/inexact_repair_certificate.py} verifies this calculation without calling the optimizer. It also accepts a user-supplied dual witness. Every reported scalar is rational, and all norm error bounds are stored as squared quantities so that no rounding of a square root is required.

\chap{Coordinates, noise, and a justified stopping rule}\label{sec:stopping-rules}
A certificate can be exact and still be interpreted incorrectly. Two common errors are to measure a coordinate covector with the wrong norm, or to keep a solver tolerance fixed while sending the regularization parameter to zero. Both change what a small reported number means.

\subsection{The residual of the coordinate equation is a covector}
In finite dimensions, let the columns of $B$ form a basis of $H=\ker p$. Write $a=Bc$, $S=dB$, and $G=B^{\mathsf T}B$. As before, $G$ is positive definite. For an inexact coefficient vector $c$, the residual of the coordinate normal equation is
\[
 q=(S^{\mathsf T}S+\alpha G)c-S^{\mathsf T}r.
\]
The vector $q$ stores pairings against basis vectors. It is not the coordinate vector of the Hilbert normal residual unless the basis is orthonormal.

\begin{proposition}[The invariant norm of a coordinate certificate]\label{WR15}
The actual normal residual is
\begin{equation}\label{eq:covector-norm}
 g=BG^{-1}q,\qquad \norm g^2=q^{\mathsf T}G^{-1}q.
\end{equation}
For any invertible change of basis $B'=BC$, the same proposed correction has
\[
 c'=C^{-1}c,\quad G'=C^{\mathsf T}GC,\quad q'=C^{\mathsf T}q,
\]
and
\[
 (q')^{\mathsf T}(G')^{-1}q'=q^{\mathsf T}G^{-1}q.
\]
\end{proposition}
\begin{proof}
Because $B$ has independent columns, for every nonzero coefficient vector $c$,
\[
 c^{\mathsf T}Gc=c^{\mathsf T}B^{\mathsf T}Bc=\norm{Bc}^2>0.
\]
Thus $G$ is positive definite and invertible. The matrix $P=BG^{-1}B^{\mathsf T}$ is symmetric and obeys
\[
 P^2=BG^{-1}(B^{\mathsf T}B)G^{-1}B^{\mathsf T}=P,
 \qquad PB=B.
\]
Its range is precisely the column space $H$, so it is the orthogonal projection onto the entire permitted correction space.

The restricted Hilbert normal residual is
\[
 g=P d^{\mathsf T}(dBc-r)+\alpha Bc.
\]
Substitute $S=dB$, $G=B^{\mathsf T}B$ and the expression for $P$:
\[
 g=BG^{-1}\bigl(S^{\mathsf T}(Sc-r)+\alpha Gc\bigr)=BG^{-1}q.
\]
Equivalently $q=B^{\mathsf T}g$ records the pairings of $g$ with the basis columns. Its squared Hilbert norm is
\[
 g^{\mathsf T}g
 =q^{\mathsf T}G^{-1}B^{\mathsf T}BG^{-1}q
 =q^{\mathsf T}G^{-1}q,
\]
using symmetry of $G^{-1}$. This proves both coordinate formulas.

For the basis change $B'=BC$, the same ambient vector is represented by $c'=C^{-1}c$. Also $S'=SC$ and $G'=C^{\mathsf T}GC$. Substitute these formulas into the new coordinate normal equation:
\[
\begin{aligned}
 q'&=((S')^{\mathsf T}S'+\alpha G')c'-(S')^{\mathsf T}r\\
 &=C^{\mathsf T}(S^{\mathsf T}S+\alpha G)CC^{-1}c
      -C^{\mathsf T}S^{\mathsf T}r=C^{\mathsf T}q.
\end{aligned}
\]
The inverse metric is $(G')^{-1}=C^{-1}G^{-1}(C^{\mathsf T})^{-1}$. Hence
\[
 (q')^{\mathsf T}(G')^{-1}q'
 =q^{\mathsf T}C C^{-1}G^{-1}(C^{\mathsf T})^{-1}C^{\mathsf T}q
 =q^{\mathsf T}G^{-1}q.
\]
For a zero-dimensional kernel these are empty matrices and zero vectors, and the resulting scalar norm is zero. The spanning assumption is indispensable: a basis of a proper subspace would compute a projection to a different correction problem.
\end{proof}
\written{WR15}

In the previous example, $B=(1,1)^{\mathsf T}$, $G=2$, and the proposal $c=1/4$ gives $q=-1/4$. Thus $q^{\mathsf T}G^{-1}q=1/32$, as required. Replacing $B$ by $sB$ changes $q$ to $sq$ and $G$ to $s^2G$. The Euclidean norm of the coordinate covector can be made arbitrarily small or large by changing the nonzero scalar $s$. The actual correction, normal residual, and certified error do not change.

\fixedplate{39_covector_metric}{Rescaling the basis of the same protected line rescales the raw coordinate covector. Its Euclidean norm is therefore not a representation-invariant stopping criterion. The Gram-corrected norm stays equal to $1/\sqrt{32}$ for every displayed scale. All systems represent the same proposed correction.}{fig:covector-metric}

The finite verifier checks that a supplied basis spans the \emph{whole} protected kernel. A basis of a smaller subspace might give a correct certificate for that smaller optimization problem while missing an allowed direction that decreases the objective. The rank, kernel-membership, and independence checks therefore belong to the certificate contract.

\subsection{Separate data error, numerical error, and regularization bias}
Let $r$ be the intended datum and $\widetilde r$ the measured datum. Suppose
\[
 \norm{\widetilde r-r}\le\delta,
 \qquad \norm{T^*(Ta-\widetilde r)+\alpha a}\le\varepsilon.
\]
The maps $T,p$ and their Hilbert metrics remain fixed. By WR13 and the fixed-parameter part of WR09,
\begin{equation}\label{eq:combined-solver-data}
 \norm{a-a_\alpha(r)}
 \le \frac{\varepsilon}{\alpha}+\frac{\delta}{2\sqrt\alpha}.
\end{equation}
The first term is the solver's contribution. The second is the worst-case propagation of data error. They should not be collapsed into one unqualified residual tolerance.

\begin{theorem}[A consistent inexact stopping rule]\label{WR16}
Assume the projected intended datum belongs to the actual range of $T$, with minimum-norm preimage $a^\dagger$. Along a sequence of positive parameters and measured problems, suppose
\[
 \alpha_n\to0,\qquad
 \frac{\delta_n}{\sqrt{\alpha_n}}\to0,\qquad
 \frac{\varepsilon_n}{\alpha_n}\to0.
\]
If the proposed corrections $a_n\in H$ satisfy the displayed data and normal-residual bounds, then $a_n\to a^\dagger$ in norm. At every index,
\[
 \norm{a_n-a^\dagger}
 \le\frac{\varepsilon_n}{\alpha_n}
      +\frac{\delta_n}{2\sqrt{\alpha_n}}
      +\norm{a_{\alpha_n}(r)-a^\dagger}.
\]
\end{theorem}
\begin{proof}
At index $n$, distinguish two exact minimizers: let $\widetilde a_n=a_{\alpha_n}(\widetilde r_n)$ solve the measured problem and let $a_n^0=a_{\alpha_n}(r)$ solve the intended-data problem. All three corrections, including the proposal $a_n$, belong to the same space $H$.

The supplied normal-residual bound is for the measured datum. Therefore WR13 applies first to that problem and gives
\[
 \norm{a_n-\widetilde a_n}\le\frac{\varepsilon_n}{\alpha_n}.
\]
The pointwise data estimate WR09, at the same parameter, then gives
\[
 \norm{\widetilde a_n-a_n^0}\le\frac{\delta_n}{2\sqrt{\alpha_n}}.
\]
Insert these two intermediate vectors between $a_n$ and $a^\dagger$ in the triangle inequality to obtain the displayed three-term estimate.

The first term tends to zero by the solver-tolerance hypothesis and the second by the data/parameter hypothesis. Since the projected intended datum is attained and $\alpha_n\to0$, the last term tends to zero by WR08. For a direct epsilon argument, make each term smaller than one third of a prescribed tolerance beyond its own threshold and take the maximum of the three thresholds. This gives norm convergence of the proposals.

Protection is not merely asymptotic: $a_n\in\ker p$ implies $p(f-a_n)=pf$ at every index. The proof compares no changing operators or changing kernels. Its pointwise certificate is formally verified, while the appeal to WR08 is the separately identified written convergence step.
\end{proof}
\written{WR16}

The current Lean addition proves the pointwise estimate \eqref{eq:combined-solver-data}. The limit conclusion here still uses the written WR08 limit theorem. This distinction is recorded in the catalogue rather than hidden by the fact that the last step is a triangle inequality.

\subsection{Retaining the explicit diagonal rate with an inexact solver}
For WR12's attainable datum and $0<\delta\le1$, take $m=\lceil\delta^{-2/3}\rceil$ and $\alpha=m^{-2}$ as before. If an inexact solver satisfies
\[
 \varepsilon\le\kappa\alpha\delta^{1/3}
 \quad(\kappa\ge0),
\]
then the same argument gives
\[
 \norm{a-a^\dagger}\le(\sqrt2+1+\kappa)\delta^{1/3}.
\]
Thus the inner solve can retain the proven outer accuracy if its normal residual decreases on the scale $\alpha\delta^{1/3}$. Since $\alpha$ is on the scale $\delta^{4/3}$, the corresponding normal-residual scale is $\delta^{5/3}$, with the stated ceiling factors retained in the exact rule. This is a sufficient rule for this problem, not a universal optimal stopping law.

\fixedplate{40_stopping_schedule}{For the attainable diagonal example, the three displayed terms are the proved bias bound, propagated data-noise bound, and a solver-error budget chosen to preserve the same $\delta^{1/3}$ rate. The ceiling in $m=\lceil\delta^{-2/3}\rceil$ is evaluated explicitly. The plot shows bounds, not measurements from a trained system.}{fig:stopping-schedule}

For the escaping datum, the same rule does not create convergence to a nonexistent protected Hilbert answer. It only makes the numerical proposals track the regularized path accurately. Tracking an escaping path and converging to a realization are different conclusions.

\subsection{The executable verification contract}
The program takes $d,p,f,\alpha$ and a proposed ambient correction, with optional measured residuals, a full kernel basis, and an independent dual witness. It rejects a correction that changes the protected data and checks the basis's dimensions, rank, and kernel membership. It does not silently project a rejected proposal into another problem.

The output records $f-a$, protection before and after, measured fit, actual model residual, the Gram-corrected normal certificate, objective values, and squared error bounds. Verification recomputes the record without calling an optimizer. Independent determinant-based tests compare the bounds with actual errors in small ambient constrained systems. A certificate establishes consistency with its included problem, not authentication that this is the externally intended task.

\begin{insight}
A certificate needs a declared target. Fit measures agreement with data. The normal residual bounds optimizer error at a specified parameter. The duality gap bounds objective error. A realization theorem is still needed to conclude convergence to an intended exact object.
\end{insight}

\input{chapters/06_cases}
\chap{Mathematical meaning and the remaining boundaries}\label{sec:meaning}
\subsection{A relative theory can be complete without being absolute}
The strict space \(\Gamma(F)\) is canonical relative to its diagram: its universal property determines it up to the appropriate unique isomorphism. The additive obstruction is canonical relative to its admissible corrections. A marked realization problem is invariant under the specified comparison equivalences. A repair can be canonical relative to a protected observable, or orthogonal relative to a fixed inner product.

These are meaningful and different kinds of canonicity. They need not culminate in one object independent of every choice of views, coefficients, topology, semantics, and geometry. Some changes of view intentionally forget information. Others strengthen the question. Treating every such change as an equivalence would remove the very distinctions the calculus is designed to preserve.

The appropriate central statement is therefore not that every possible mathematical presentation lives in one fixed universal space. It is that a declared system of representations can be organized into coherent realizations, with explicit extension, comparison, and repair problems. A proposed bridge between two such systems has to show what structure it preserves and what information it loses.

\subsection{Internal realization and intended objects}
Suppose a separate category \(\mathcal S\) contains the intended semantic objects of an application, and a realization construction gives a comparison functor
\[
 \mathcal R:\mathcal S\longrightarrow\mathcal K,
\]
where \(\mathcal K\) is the relevant category or groupoid of coherent data. Defining \(\mathcal K\) does not establish that \(\mathcal R\) is faithful, full, or essentially surjective. These are new mathematical questions about the actual semantic comparison.

Faithfulness asks whether distinct semantic arrows can become indistinguishable. Fullness asks whether every compatible comparison comes from a semantic one. Essential surjectivity asks whether every intended coherent datum is realized by an actual semantic object. Depending on the application, one may seek only a restricted image or a classified obstruction to belonging to it.

This distinction prevents a formal globalization of encoded data from being mistaken for an unrestricted arithmetic or physical existence theorem. The abstract construction does not manufacture missing semantic objects.

\subsection{Every finite test can pass while the infinite problem fails}
\begin{example}[Finite prefixes without a global point]\label{KS54}
Let \(X_n=\{m\in\N:m\ge n\}\), with bonding maps given by inclusion. Every level is nonempty, and every finite prefix has a coherent point. The inverse limit is nevertheless empty.
\end{example}
\begin{proof}
For a fixed terminal level \(N\), the integer \(N\) belongs to every fiber \(X_n\) with \(n\le N\). Assign that same integer at each such level. Every restriction is an inclusion, so it preserves the assigned integer and all finite coherence equations hold. This proves finite solvability with the quantifiers in the order ``for every \(N\), there exists a prefix.''

Suppose instead that \(x=(x_n)\) is a global point. Equality along each inclusion says that the underlying natural numbers of \(x_{n+1}\) and \(x_n\) are equal. Induction therefore makes every underlying coordinate equal to the number \(m\) underlying \(x_0\). Membership of \(x_{m+1}\) in \(X_{m+1}\) now gives \(m\ge m+1\), a contradiction. Thus no global point exists.

There is no conflict with a gluing theorem for prefixes. The prefix chosen above for length \(N+1\) restricts to the value \(N+1\), not to our chosen value \(N\) on the shorter prefix. Nor could another choice of one solution for every length give a matching family: its value at level zero would force the same forbidden global integer. Separate existence at every length is strictly weaker than existence of a compatible family of prefix solutions.
\end{proof}
\source{KS54}
This counterexample is formalized in \code{FinitePrefixBoundary.lean} using the original resolution and partial-section types. It explains why inverse-limit existence needs more than separate finite solvability. Finiteness, compactness, surjectivity, or another coherence principle must be proved in the relevant application. The example does not contradict the strict gluing theorem: a compatible matching family of all prefixes was never supplied.

\subsection{A positive infinite-existence criterion}
The failed example isolates a missing compactness or lifting hypothesis. Finite fibers supply a concrete positive criterion without requiring every bonding map to be surjective.

\begin{proposition}[Finite fibers and persistent images]\label{KS58}
Let \((X_n,\rho_n)\) be a sequential inverse system of nonempty finite sets. Its inverse limit is nonempty. More precisely, if \(\rho_{n,m}:X_m\to X_n\) is composite restriction, then
\[
 \im(\operatorname{ev}_n)=Y_n,
 \qquad Y_n=\bigcap_{m\ge n}\im\rho_{n,m}.
\]
Thus a prescribed value at level \(n\) extends precisely when it survives every later image.
\end{proposition}
\begin{proof}
Write \(I_{n,m}=\rho_{n,m}(X_m)\) for \(m\ge n\). Each \(I_{n,m}\) is nonempty. The composition law gives
\[
 I_{n,m+1}=\rho_{n,m}\bigl(\rho_m(X_{m+1})\bigr)\subseteq I_{n,m}.
\]
These are subsets of the finite set \(X_n\). Every strict decrease removes an element permanently, so only finitely many strict decreases are possible. There is an index \(s_n\ge n\) such that \(I_{n,m}=I_{n,s_n}\) for all \(m\ge s_n\). Consequently \(Y_n=I_{n,s_n}\) is nonempty.

We next prove that restriction induces a \emph{surjection} \(Y_{n+1}\to Y_n\), even though \(X_{n+1}\to X_n\) need not be surjective. For membership, let \(z\in Y_{n+1}\). Given \(m\ge n+1\), write \(z=\rho_{n+1,m}(w)\) for some \(w\in X_m\). Then \(\rho_n z=\rho_{n,m}(w)\). Together with the automatic condition at \(m=n\), this shows \(\rho_nz\in Y_n\).

For surjectivity, fix \(y\in Y_n\) and choose
\[
 m\ge\max\{s_n,s_{n+1},n+1\}.
\]
Since \(y\in I_{n,m}\), there is \(w\in X_m\) with \(\rho_{n,m}(w)=y\). Put \(z=\rho_{n+1,m}(w)\). Stabilization at level \(n+1\) gives \(z\in Y_{n+1}\), and composition gives \(\rho_nz=y\).

Choose \(x_0\in Y_0\) and recursively choose \(x_{n+1}\in Y_{n+1}\) above \(x_n\). The preceding surjectivity justifies each step. Adjacent coherence implies every composite coherence equation, so this gives a point. This is an existence construction using countably many choices, not an asserted effective procedure for locating the stabilization indices.

Finally, fix a level \(n\). Every global point satisfies \(x_n=\rho_{n,m}(x_m)\) for each \(m\ge n\), hence \(x_n\in Y_n\). Conversely, start with any \(y\in Y_n\), lift recursively above level \(n\) through the surjections just proved, and define the lower coordinates by \(x_i=\rho_{i,n}(y)\). The composite restriction laws verify coherence across and below level \(n\). The resulting point evaluates to \(y\). These two inclusions prove the asserted equality of evaluation images, not merely nonemptiness of the inverse limit.
\end{proof}
\source{KS58}
This is the standard finite-set inverse-limit principle, here with its sequential persistent-image proof~\cite{stacks-limits}. Its formalization in \code{FiniteRealizationExistence.lean} constructs an explicit equivalence with the pinned library's section type, then applies the finite cofiltered-system and eventual-range theorems. The printed sequential proof and this formal route establish the same stated criterion. Both use ordinary classical choice. Surjective bonding maps on arbitrary nonempty sets give a second sufficient condition by direct recursive lifting. Neither criterion says that arbitrary data in a non-surjective finite tower extend.

\subsection{A finite computation that exposes the preservation contract}
Consider a raw vector \(x=(2,0,5)\in\R^3\) with defect \(d(x)=x_1-x_2\). Unconstrained orthogonal repair gives \((1,1,5)\), with squared error two. If the first and third coordinates must be preserved, the unique exact repair is \((2,2,5)\), with squared error four.

Both outputs are exact. Only one satisfies the second preservation problem. Neither theorem contradicts the other because the admissible changes differ. An implementation that reports only zero residual would not distinguish the two tasks.

For a finite linear system over a field, the same question admits independently checkable positive and negative certificates. Let \(A\in\K^{m\times n}\) encode compatibility, \(P\in\K^{p\times n}\) the protected values, and \(x\in\K^n\) the raw candidate. Let the columns of \(N\) be a basis of \(\ker P\). The correction equation is
\[
 ANz=Ax,\qquad y=x-Nz.
\]
It is equivalent to the stacked system \(Ay=0\), \(Py=Px\).

\begin{proposition}[A finite certificate alternative]\label{KS57}
Exactly one of the following occurs:
\begin{enumerate}
\item There is \(y\) with \(Ay=0\) and \(Py=Px\). Every answer is \(y+H\), where \(H=\ker A\cap\ker P\).
\item There are row vectors \(\lambda\in\K^{1\times m}\) and \(\mu\in\K^{1\times p}\) with
\[
 \lambda A=\mu P,\qquad\lambda Ax=1.
\]
\end{enumerate}
The negative witness annihilates every admissible correction defect, since \(\lambda AN=\mu PN=0\). The ambiguity dimension in the positive case is
\[
 \dim H=n-\operatorname{rank}\begin{pmatrix}A\\P\end{pmatrix}.
\]
\end{proposition}
\begin{proof}
Set \(T=\begin{pmatrix}A\\P\end{pmatrix}\) and \(b=\begin{pmatrix}0\\Px\end{pmatrix}\). A vector \(y\) is an exact protected answer if and only if \(Ty=b\). Perform Gaussian row operations on the augmented matrix, keeping their product \(U\). Each elementary row operation is invertible over the given field. Thus \(Ty=b\) and \(UTy=Ub\) have exactly the same solutions.

If no row of the reduced augmented system has zero left side and nonzero right side, assign arbitrary values to its free variables and solve for its pivot variables. This supplies a solution \(y_0\). For any vector \(y\),
\[
 Ty=b\quad\Longleftrightarrow\quad T(y-y_0)=0.
\]
Translation by \(y_0\) is therefore a bijection from \(\ker T\) to the answer set, with inverse \(y\mapsto y-y_0\). Rank--nullity applied to the actual linear map \(T:\mathbb F^n\to\mathbb F^{m+p}\) gives the ambiguity dimension. Dependent prescription rows do not alter this argument.

If the reduced system is inconsistent, choose a row with zero left side and right side \(\alpha\ne0\). Let \(w\) be the corresponding row of \(U\), and split it according to the two blocks as \(w=(u,v)\). The retained row operations give the literal identities
\[
 uA+vP=0,\qquad vPx=\alpha.
\]
Division is permitted because \(\alpha\) is a nonzero field element. Define \(\lambda=-\alpha^{-1}u\) and \(\mu=\alpha^{-1}v\). Multiplying the displayed equations by \(\alpha^{-1}\) gives
\[
 \lambda A=\mu P,\qquad \lambda Ax=\mu Px=1.
\]
For every admissible correction \(a\in\ker P\), these same row equations imply \(\lambda Aa=\mu Pa=0\). The witness therefore annihilates the entire correction image while taking value one on the required defect.

The alternatives cannot occur together. If \(y\) were an exact protected answer in the presence of such a witness, then
\[
 1=\lambda Ax=\mu Px=\mu Py=\lambda Ay=0.
\]
A field has \(1\ne0\), so this is impossible. Consistency or inconsistency of the reduced system proves exhaustiveness, and the contradiction proves exclusivity. Empty row or column blocks cause no extra division: normalization is used only in the inconsistent case, where the nonzero scalar has already been obtained. The proof does not transfer to integral coefficients by silently allowing the same division.
\end{proof}
\source{KS57}
The formal proof gives these literal row-vector equations over an arbitrary field. It uses algebraic dual separation of the stacked image rather than verifying the row-reduction program. The affine answer equivalence and rank-nullity formula are proved separately. Thus a normalized obstruction is an actual checkable witness, while universal correctness of the supplied Python elimination implementation remains a different question.

\subsection{One equation, three preservation problems}
For the following three cases work over \(\Q\), with \(A=(1,-1,0)\) and \(x=(2,0,5)\). With \(P(x)=(x_1,x_3)\), the unique answer is \((2,2,5)\). With \(P(x)=x_3\), the entire answer set is \(\{(t,t,5):t\in\Q\}\). With \(P(x)=(x_1,x_2)\), no answer exists. In this last case
\[
 \lambda=(1/2),\qquad\mu=(1/2,-1/2)
\]
is a normalized negative certificate over \(\Q\) or \(\R\): \(\lambda A=\mu P\) and \(\lambda Ax=1\).

The package implements this finite interface over \(\Q\) in \code{examples/protected\_repair.py}. It returns a particular exact answer and an ambiguity basis, or a normalized dual witness. Its verifier checks the certificate equations directly. A Gaussian particular answer is \emph{not} advertised as the nearest answer of Theorem~\ref{KS56}. That is a separate geometric selection.

Exact field elimination does not solve the integer parity problem by changing its coefficient domain. Over \(\Q\), \(2m=1\) has the solution \(1/2\). Over \(\Z\), it has none. The rational interface rejects floating-point inputs rather than silently treating an approximate decimal as exact arithmetic. Its finite tests establish the tested implementation behavior, not an arbitrary-input Lean refinement theorem.

\subsection{Why this is relevant to machine reasoning}
The framework gives a machine explicit mathematical objects to manipulate: views, restrictions, partial realizations, typed defects, admissible corrections, and comparison witnesses. A successful step can carry both an output and a certificate of what remained fixed. A failed step can return an obstruction in a specified quotient rather than an unexplained failure flag.

The split coordinate theorem is especially direct. In its stated class, the same candidate is represented by protected information and residual data, and repair zeros only the residual coordinate. The descent theorem explains which local certificates can be assembled. The analytic results explain when a small measured defect has a quantitatively reliable meaning.

These observations motivate a machine oriented mathematical language. They do not establish an autonomous reasoning system, an empirical intelligence gain, or a general efficient algorithm. Classical choice, infinite products, quotient membership, and bounded inverse constructions can be mathematically valid without yielding executable finite procedures. Implementation claims need their own specifications and tests.

\subsection{What this monograph closes, and what it leaves open}
The book supplies a coherent account of strict realization spaces, additive and nonabelian extension problems, split tower repair, retained multiscale detail, a genuine composition-site stack, presentation transport, graph holonomy, closed-range stability, and moving exact spaces. The additional convergence and obstruction development separates small measurements from small movements, coherent coordinate families from their semantic image, and zero residual from zero transformation cost. The new bounded-budget development supplies a positive existence criterion and keeps range membership separate from regularized solvability. Its fixed-budget dual witness can certify a limit on every admissible correction, while its diagonal example computes the actual leading cost of approximation. Each layer is attached to the data that make it meaningful.

The following broader targets are not conclusions of this edition: arbitrary higher categorical coefficients and infinity-stack descent, unrestricted changes of view doctrine, norm reconstruction for every bounded infinite residual hierarchy, nonlinear or unknown-connection repair, perturbation across changing Hilbert geometries without new comparison data, and semantic essential surjectivity for unrestricted arithmetic or physical applications.

The accepted integrated infrastructure has \EditionFormalModules\ production modules, a single root, compiler records, declaration axiom reports, and rejection controls. All five earlier completion items KS54--KS58 have formal counterparts. The separate version 1.6.0 regularization extension establishes the focused counterparts of WR06, WR07, WR13, and WR14 and the pointwise estimates in WR09 and WR16. The remaining ten WR groups and the limiting portions of the two partial groups retain written proofs. This edition changes neither source tree and claims no new compiler run. The longer exposition and analytic background appendix improve the printed argument without enlarging the verified source scope. Further extensions require their own hypotheses and evidence.

\begin{insight}
\textbf{The organizing principle.} Coherent realization, obstruction to extension, freedom of completion, and stability of repair belong to one mathematical conversation. Their unification is useful precisely when it does not identify them with one another.
\end{insight}

\appendix
\input{chapters/09_workshop}
\input{chapters/10_hypothesis_atlas}
\input{chapters/07_concordance}
\clearpage
\chap{Extension proofs and their verification scope}\label{app:written}
The WR catalogue is deliberately separate from the eighty KS exposition items. It began with twelve complete written theorem groups. Version 1.6.0 added WR13--WR16 and supplied a focused formal counterpart for selected parts of the catalogue. The present edition expands all sixteen printed proofs while retaining those evidence boundaries. It does not promote a written limit theorem merely because a related pointwise estimate has now been compiled.

\begin{center}
\begin{tabularx}{\linewidth}{@{}p{.19\linewidth}X@{}}
\toprule Coverage & Exact interpretation\\\midrule
Complete focused & WR06, WR07, WR13 and WR14. The principal universal statements have typed counterparts in the three new modules and their actual protected bridge. Explanatory examples retain their stated separate scope.\\[5pt]
Partial focused & WR09 and WR16. The fixed-parameter and pointwise solver/data inequalities are formal. Their parameter-limit consequences still use written WR08.\\[5pt]
Written only & WR01--WR05, WR08, WR10--WR12 and WR15. Full proofs are printed. No new typed counterpart is assigned.\\\bottomrule
\end{tabularx}
\end{center}

The retained version 1.6.0 focused replay recompiled seven targets from source: three new modules, three inherited dependencies, and the root. It uses the pinned official Lean compiler and imported dependency caches. The canonical \path{formal/} tree remains unchanged, with its original accepted cumulative evidence. The source resolver binds WR references to the focused imported declaration types. This is an editorial correspondence, not automatic natural-language proof verification.

\subsection{Existence and budget geometry}
\textbf{WR01. Bounded approximate feasibility.} The real Hilbert operator and residual datum are fixed. A uniform correction-norm bound, not residual decay alone, gives an exact preimage by weak subsequence compactness. This remains a written proof.

\textbf{WR02. Nested affine minimum corrections.} Nonempty closed affine constraint sets are nested in one Hilbert space. The exact energy identity makes bounded minimum-norm corrections converge strongly. No formal compactness or affine-limit module is added here.

\textbf{WR03. Attained budget profile.} A finite nonnegative budget has an attained residual minimum even when the full operator range is not closed. The proof retains the weak topology argument rather than substituting finite-dimensional compactness.

\textbf{WR04. Sharp budget support.} The general support formula uses the actual restricted adjoint. The focused module proves the support inequality and supplies equality at a regularized optimizer's own budget. It does not prove the entire general budget-attainment route, so WR04 remains written-only as a group.

\textbf{WR05. Robust budget comparison.} The operators act on the same correction space, with operator and data errors bounded by $\eta$ and $\delta$. The bound $\delta+\eta M$ does not apply automatically when the protected kernel changes.

\subsection{Regularization and its limit}
\textbf{WR06. Constructed regularized minimizer.} Positive $\alpha$ makes $T^*T+\alpha I$ coercive. The formal proof constructs an actual continuous linear equivalence, defines the correction using its inverse, and proves the energy identity and uniqueness. The manuscript's self-adjoint closed-range proof and the formal coercive-form proof establish the same stated result. The bridge uses the existing \path{ProtectedHilbertRepair.restricted} operator.

\textbf{WR07. Parameter and budget comparison.} The formal statements compare two positive parameters, prove the two norm monotonicities, and show that no correction of smaller or equal norm has a lower residual than the regularized answer. The normalized residual has unit norm and attains the support expression whenever it is nonzero. The Lean formulation gives the optimizer and universal competitor inequality directly rather than introducing a second numerical profile definition.

\textbf{WR08. The regularization limit.} Projection onto the closure of the range is separated from actual range membership. Strong convergence of corrections requires the projected datum to be attained. Otherwise the correction norms diverge. These infinite-limit arguments remain written-only.

\textbf{WR09. Data sensitivity and parameter choice.} The focused proof supplies the coefficients $1/(2\sqrt\alpha)$ for correction noise and one for fitted-data noise. The scalar sharpness calculations and the parameter-choice convergence consequence retain their printed proofs. Formalizing a fixed-parameter inequality does not formalize its convergence application automatically.

\textbf{WR10--WR12. The explicit diagonal models.} The exact regularized family, omitted-tail bounds, matched-budget limit $\pi/4$, and attainable noisy datum remain written with exact finite enclosure tests. No general infinite-sum formalization or new diagonal Lean instance is claimed.

\subsection{Inexact solution and coordinate interpretation}
\textbf{WR13. A posteriori solver bounds.} The correction error, fitted error, and objective excess are bounded in terms of the actual normal residual and positive parameter. The focused theorem quantifies over every proposal in the Hilbert source. The protected bridge transfers its candidate-distance bound to the actual kernel subtype. Scalar sharpness examples are separated from the principal theorem and checked independently.

\textbf{WR14. A regularized primal and dual certificate.} An exact two-square identity proves weak duality, a dual maximizer constructed from the primal correction, global objective bracketing, and a correction-distance bound. The witness is arbitrary. Setting it to the proposal's fit residual gives the exact normal-residual gap.

\textbf{WR15. The coordinate covector metric.} A full independent kernel basis gives $g=BG^{-1}q$ and $\norm g^2=q^{\mathsf T}G^{-1}q$. The invariance proof and its invertibility hypotheses are printed and rationally tested. They have not been compiled as a new matrix theorem.

\textbf{WR16. A justified stopping rule.} The focused theorem proves the pointwise sum of the solver and data bounds. The printed convergence consequence uses written WR08 and the additional assumptions $\alpha\to0$, $\delta/\sqrt\alpha\to0$, and $\epsilon/\alpha\to0$. The displayed diagonal schedule also uses written WR12.

\subsection{Independent exact implementation and restart}
The rerun regularization suite contributes \EditionRegularizationChecks\ assertions and the rerun inexact-certificate suite \EditionInexactChecks. The independent oracle solves ambient constrained systems by determinant formulas rather than using the certificate program's kernel coordinates. The verifier checks a full kernel basis, retains its Gram metric, preserves the observable, and distinguishes measured fit from the actual residual. All outputs are exact rationals or squared bounds. It recomputes records and rejects altered claims without calling an exact optimizer.


\clearpage
\chap{Analytic facts used in the proofs}\label{app:proof-tools}
This appendix supplies the analytic ingredients used at several points in the book. They are classical facts, not additional novelty claims. Writing out their proofs makes the dependence on completeness, convexity, the inner product, and the declared operator norm visible. The statements are restricted to the real Hilbert and real Banach settings used in the main arguments. Their role in the source correspondence is background support, not a new KS or WR entry.

The projection and inverse constructions have standard formal counterparts in the mathematical library cited in \cite{mathlib-projection,mathlib-banach}. The regularization context is also discussed in \cite{clason-inverse}. The arguments below are ordinary written proofs. Retaining an accepted Lean source tree does not by itself give a new compiler record for this appendix.

\subsection{Projection onto a closed convex set}
\begin{lemma}[Existence, uniqueness, and the variational inequality]\label{lem:projection-tool}
Let \(H\) be a real Hilbert space, let \(C\subseteq H\) be nonempty, closed, and convex, and fix \(x\in H\). There is a unique \(c\in C\) minimizing \(\norm{x-c}\). It is characterized by
\begin{equation}\label{eq:projection-variational}
 \langle x-c,z-c\rangle\le0\qquad(z\in C).
\end{equation}
If \(C=c_0+L\) is a closed affine subspace with direction space \(L\), the condition is \(x-c\perp L\).
\end{lemma}
\begin{proof}
\textbf{Produce a Cauchy minimizing sequence.} Put \(d=\inf_{z\in C}\norm{x-z}\). Nonemptiness makes this a finite nonnegative real number. Choose \(c_n\in C\) so that \(\norm{x-c_n}\to d\). Convexity places \((c_n+c_m)/2\) in \(C\). The parallelogram identity gives
\[
 \norm{c_n-c_m}^2
 =2\norm{x-c_n}^2+2\norm{x-c_m}^2
   -4\left\|x-\frac{c_n+c_m}{2}\right\|^2
 \le2\norm{x-c_n}^2+2\norm{x-c_m}^2-4d^2.
\]
The right side tends to zero as both indices tend to infinity. Thus \((c_n)\) is Cauchy. Completeness supplies a limit \(c\in H\), closedness puts it in \(C\), and continuity of the norm gives \(\norm{x-c}=d\). A minimizing sequence in an arbitrary closed set need not have been Cauchy. The midpoint was the essential use of convexity.

\textbf{Prove uniqueness.} If \(c,c'\in C\) both attain \(d\), the same midpoint inequality gives
\[
 \norm{c-c'}^2\le2d^2+2d^2-4d^2=0.
\]
Therefore \(c=c'\).

\textbf{Derive the variational inequality.} For \(z\in C\) and \(0<t\le1\), convexity makes \(c+t(z-c)\) feasible. Minimality and the inner-product expansion yield
\[
 0\le\norm{x-c-t(z-c)}^2-\norm{x-c}^2
   =-2t\langle x-c,z-c\rangle+t^2\norm{z-c}^2.
\]
Divide by the positive number \(t\) and let \(t\downarrow0\). This gives \eqref{eq:projection-variational}. Conversely, if that inequality holds, then for every \(z\in C\),
\[
 \norm{x-z}^2
 =\norm{x-c}^2-2\langle x-c,z-c\rangle+\norm{z-c}^2
 \ge\norm{x-c}^2.
\]
It therefore characterizes the unique minimizer.

For an affine subspace, \(c+h\) and \(c-h\) are feasible for every \(h\in L\). Applying the inequality with both signs gives \(\langle x-c,h\rangle=0\). Orthogonality conversely implies the same variational inequality. In this case the preceding expansion becomes the Pythagorean equality, not merely an inequality.
\end{proof}

For a closed linear subspace \(L\), denote its projection by \(P_L\). The lemma gives a unique decomposition
\[
 x=P_Lx+(x-P_Lx),\qquad P_Lx\in L,\quad x-P_Lx\in L^\perp.
\]
Adding two such decompositions, or multiplying one by a scalar, proves linearity by uniqueness. The Pythagorean identity gives \(\norm{P_Lx}\le\norm{x}\), so \(P_L\) is bounded. It fixes \(L\), annihilates \(L^\perp\), and satisfies \(P_L^2=P_L\). For \(x,y\in H\), orthogonality gives
\[
 \langle P_Lx,y\rangle=\langle P_Lx,P_Ly\rangle
 =\langle x,P_Ly\rangle.
\]
Thus \(P_L\) is self-adjoint. These properties are conclusions of the construction, not extra data assumed for a proposed projection.

\subsection{Riesz representation and the adjoint}
\begin{lemma}[Representing a bounded functional]\label{lem:riesz-tool}
For every bounded real linear functional \(\ell:H\to\R\) on a real Hilbert space, there is a unique \(u\in H\) with
\[
 \ell(h)=\langle h,u\rangle\quad(h\in H),\qquad \norm{\ell}=\norm{u}.
\]
\end{lemma}
\begin{proof}
If \(\ell=0\), take \(u=0\). Otherwise choose \(x\) with \(\ell(x)\ne0\). Boundedness makes \(L=\ker\ell\) closed. Use Lemma~\ref{lem:projection-tool} to form \(v=x-P_Lx\). Then \(v\perp L\) and \(\ell(v)=\ell(x)\ne0\), in particular \(v\ne0\).

For any \(h\in H\), the vector \(h-\ell(h)v/\ell(v)\) belongs to \(L\). Pairing it with \(v\) therefore gives
\[
 \langle h,v\rangle=\frac{\ell(h)}{\ell(v)}\norm{v}^2.
\]
Consequently \(u=\ell(v)v/\norm{v}^2\) represents \(\ell\). If \(u'\) also represents it, pairing the difference of the two formulas with \(h=u-u'\) gives \(\norm{u-u'}^2=0\), hence uniqueness.

Cauchy--Schwarz gives \(|\ell(h)|\le\norm{u}\norm{h}\), so \(\norm{\ell}\le\norm{u}\). If \(u\ne0\), the unit vector \(h=u/\norm{u}\) gives the reverse inequality. The zero case was already included.
\end{proof}

Let \(T:H\to R\) be a bounded linear map between real Hilbert spaces. For fixed \(z\in R\), the functional \(h\mapsto\langle Th,z\rangle\) is bounded by \(\norm{T}\norm{z}\norm{h}\). The lemma therefore defines a unique vector \(T^*z\in H\) with
\[
 \langle Th,z\rangle=\langle h,T^*z\rangle.
\]
Uniqueness of representatives proves that \(T^*\) is linear. The same bound proves \(\norm{T^*z}\le\norm{T}\norm{z}\). Conversely,
\[
 \norm{Th}=\sup_{\norm{z}\le1}|\langle Th,z\rangle|
 \le\norm{T^*}\norm{h},
\]
where the first equality follows by testing \(z=Th/\norm{Th}\) when \(Th\ne0\). Hence \(\norm{T^*}=\norm{T}\).

The defining pairing also gives the identity
\begin{equation}\label{eq:adjoint-range-orthogonal}
 (\ran T)^\perp=\ker T^*.
\end{equation}
Indeed, \(z\) is orthogonal to every \(Th\) exactly when \(\langle h,T^*z\rangle=0\) for every \(h\), which is equivalent to \(T^*z=0\). Orthogonal complements do not change on taking closure. The projection decomposition then gives
\[
 \overline{\ran T}=(\ker T^*)^\perp.
\]
This is why a self-adjoint injective operator has dense range. It does not yet say its range is closed. In the regularized normal-operator proof, the coercive lower bound supplies that separate fact.

\subsection{Why a bounded bijection has a bounded inverse}
The closed-range arguments use the inverse theorem only after proving that the relevant source and target subspaces are complete and that the restricted map is bijective. The following derivation records the missing analytic step.

\begin{lemma}[Baire's complete-space principle]\label{lem:baire-tool}
A nonempty complete metric space cannot be the union of countably many closed sets with empty interior.
\end{lemma}
\begin{proof}
Suppose \(F_1,F_2,\ldots\) are closed with empty interior. Choose a nonempty open ball whose closed ball \(B_1\), of positive radius at most \(1/2\), avoids \(F_1\). Such a choice is possible because the complement of \(F_1\) is open and meets every nonempty open set. Recursively choose a closed ball \(B_{n+1}\) of positive radius at most \(2^{-(n+1)}\), contained in the open ball with the same center and radius as \(B_n\), and disjoint from \(F_{n+1}\). The intersection of that open ball with the complement of \(F_{n+1}\) is nonempty and open, so it contains a sufficiently small closed ball. This also covers the case of isolated points.

The centers form a Cauchy sequence: all centers after stage \(n\) lie in \(B_n\), whose diameter is at most \(2^{1-n}\). Completeness gives a limit. For each fixed \(n\), a tail of centers lies in the closed set \(B_n\), so the limit lies there. It therefore belongs to no \(F_n\), contradicting that these sets cover the space.
\end{proof}

\begin{lemma}[Continuous inverse for Banach spaces]\label{lem:inverse-tool}
Let \(S:X\to Y\) be a bounded linear surjection between real Banach spaces. There is a constant \(C<\infty\) such that every \(y\in Y\) has a preimage \(x\) with \(\norm{x}\le C\norm{y}\). If \(S\) is also injective, its inverse is bounded linear.
\end{lemma}
\begin{proof}
If \(Y=\{0\}\), take the zero preimage. Otherwise write \(B_X\) for the closed unit ball. Surjectivity gives
\[
 Y=\bigcup_{n\ge1}\overline{S(nB_X)}.
\]
By Lemma~\ref{lem:baire-tool}, at least one of the closed sets on the right has nonempty interior. Thus there are \(n\ge1\), \(y_0\in Y\), and \(\rho>0\) such that
\[
 B(y_0,\rho)\subseteq\overline{S(nB_X)}.
\]
Here \(B(y_0,\rho)\) is an open ball. If \(\norm{h}<\rho\), both \(y_0+h\) and \(y_0\) belong to this closure. Choose approximating image sequences for both. Subtracting their preimages shows \(h\in\overline{S(2nB_X)}\). Scaling by \(2n\) gives a constant \(c=\rho/(2n)>0\) with
\begin{equation}\label{eq:inverse-closure-ball}
 B(0,c)\subseteq\overline{S(B_X)}.
\end{equation}
This only gives approximate preimages. We next turn them into an actual preimage with a summable budget.

For a nonzero residual \(r\in Y\), set \(\lambda=2\norm{r}/c\). The vector \(r/\lambda\) has norm \(c/2\), so \eqref{eq:inverse-closure-ball} supplies \(z\in B_X\) with
\[
 \norm{Sz-r/\lambda}\le c/4.
\]
Then \(u=\lambda z\) satisfies
\[
 \norm{u}\le 2\norm{r}/c,\qquad
 \norm{r-Su}\le\norm{r}/2.
\]
If the residual is zero, choose \(u=0\). Start from \(r_0=y\), repeat this construction, and define \(r_{j+1}=r_j-Su_j\). Induction yields
\[
 \norm{r_j}\le2^{-j}\norm{y},\qquad
 \sum_{j\ge0}\norm{u_j}\le\frac{4}{c}\norm{y}.
\]
The series of preimages converges in \(X\) by completeness, to a vector \(x\) with the displayed norm bound. Continuity of \(S\), together with the finite telescoping equation
\[
 S\left(\sum_{j=0}^{N-1}u_j\right)=y-r_N,
\]
proves \(Sx=y\). Thus we may take \(C=4/c\).

If \(S\) is bijective, the preimage is unique. Uniqueness makes the inverse additive and homogeneous, and the bound just obtained makes it continuous. In the merely surjective case the construction does not assert a linear choice of preimage. In particular it does not assert that every Banach-space surjection has a bounded linear right inverse.
\end{proof}

For example, in Theorem~\ref{KS43} the source of the inverse is \((\ker D)^\perp\) and its target is \(\ran D\). Both are closed subspaces of the stated Banach spaces. The inverse lemma applies only after the proof has established those closedness and bijectivity facts. It would not apply with an unclosed range equipped with its inherited norm.

\subsection{Neumann inversion with both inverse identities}
\begin{lemma}[A strict perturbation radius]\label{lem:neumann-tool}
Let \(W\) be a real Banach space and \(T:W\to W\) bounded linear, with \(\norm{T}\le\theta<1\) and \(\theta\ge0\). Then
\[
 Z=\sum_{j=0}^{\infty}T^j
\]
exists in operator norm and satisfies
\[
 (I-T)Z=Z(I-T)=I,\qquad \norm{Z}\le\frac1{1-\theta}.
\]
\end{lemma}
\begin{proof}
First recall why the space of bounded operators on \(W\) is complete. An operator-norm Cauchy sequence \(A_n\) makes \(A_nx\) Cauchy for each \(x\), and hence defines \(Ax=\lim_nA_nx\). Pointwise limits preserve linearity. The sequence of operator norms is bounded, so the pointwise limit is a bounded operator. For a fixed sufficiently large \(n\), pass \(m\to\infty\) in \(\norm{(A_n-A_m)x}\le\varepsilon\norm{x}\) to obtain \(\norm{(A_n-A)x}\le\varepsilon\norm{x}\). Taking the supremum over the unit ball proves operator-norm convergence. Completeness of \(W\) is the step that produced each pointwise limit.

Submultiplicativity gives \(\norm{T^j}\le\theta^j\), with the convention \(T^0=I\). The finite partial sums are therefore operator-norm Cauchy, since for \(M>N\),
\[
 \left\|\sum_{j=N+1}^{M}T^j\right\|
 \le\sum_{j=N+1}^{M}\theta^j\le\frac{\theta^{N+1}}{1-\theta}.
\]
Their limit \(Z\) has norm at most \(\sum_{j\ge0}\theta^j=1/(1-\theta)\). The operator identities for finite sums are
\[
 (I-T)\sum_{j=0}^{N}T^j=I-T^{N+1},\qquad
 \left(\sum_{j=0}^{N}T^j\right)(I-T)=I-T^{N+1}.
\]
Multiplication by a fixed bounded operator is continuous in operator norm, and \(\norm{T^{N+1}}\le\theta^{N+1}\to0\). Passing to the limit proves both inverse identities. When \(W\) is the zero space the same statements hold, with its unique identity map. When \(\theta=1\), the summable tail and decay used here disappear, so the proof supplies no endpoint conclusion.
\end{proof}

\subsection{The dependency order}
The analytic arguments can now be read without a circular invocation of their conclusions. Projection gives the orthogonal decomposition and Riesz representation. The latter gives adjoints and the relation between range and adjoint kernel. Baire's principle independently gives the Banach inverse theorem. Those ingredients yield closed-range estimates and the positive regularized normal inverse. Neumann's series gives the separate small-perturbation inverse and its explicit radius.

The weak-subsequence proof in Section~\ref{sec:bounded-realization} uses projection onto the closed separable span generated by the sequence, followed by a coordinate diagonal argument. It does not assume that the whole ambient Hilbert space is separable. The budget-attainment and regularization-limit proofs use that sequential result. Strong convergence is proved later from the minimum-norm condition and the norm identities, not inferred from weak convergence alone.

\clearpage
\chap{Background and source documents}\label{app:references}
The references identify the standard background and frozen project sources. The principal constructions have printed proofs, with the expanded arguments retaining their original KS and WR identifiers. Appendix~\ref{app:proof-tools} also supplies the analytic tools used repeatedly in the main text. The five additions KS54 through KS58 have their earlier integrated formal counterparts. The regularization groups retain the separately stated version 1.6.0 focused scope. Live mathlib documentation is a navigation aid, not a replacement for the pinned source revision.

\begin{thebibliography}{99}
\bibitem{stacks-limits} The Stacks Project Authors. \emph{Limits and colimits over preordered sets}. Section 4.21, Tag 002Z, including the finite-set inverse-limit theorem in Lemma 4.21.7. \url{https://stacks.math.columbia.edu/tag/002Z}. Accessed 10 October 2026.
\bibitem{stacks-stack} The Stacks Project Authors. \emph{Stacks}. Section 8.4, Tag 0268. In particular the effective descent and equivalence-of-descent-categories formulations. \url{https://stacks.math.columbia.edu/tag/0268}. Accessed 10 October 2026.
\bibitem{mathlib-stack} The mathlib Community. \emph{CategoryTheory.Sites.Descent.IsStack}. Standard pseudofunctor stack interface. \url{https://leanprover-community.github.io/mathlib4_docs/Mathlib/CategoryTheory/Sites/Descent/IsStack.html}.
\bibitem{mathlib-projection} The mathlib Community. \emph{Analysis.InnerProductSpace.Projection.Basic}. Orthogonal projection onto complete subspaces. \url{https://leanprover-community.github.io/mathlib4_docs/Mathlib/Analysis/InnerProductSpace/Projection/Basic.html}.
\bibitem{mathlib-banach} The mathlib Community. \emph{Analysis.Normed.Operator.Banach}. Banach open mapping and continuous inverse constructions. \url{https://leanprover-community.github.io/mathlib4_docs/Mathlib/Analysis/Normed/Operator/Banach.html}.
\bibitem{foundation} Komorebi P. \emph{Komorebi Space Theory}, version 0.20.0. Recovered source archive and replay evidence, 28 September 2026. Included unchanged under \code{evidence/archives/}. Its source documents include \emph{Ambient Sections and Extension Obstructions}, \emph{Split Tower Coordinates}, \emph{Infinite Reconstruction}, \emph{Nonabelian Coherence}, \emph{Effective Groupoid Descent}, \emph{A Stack on the Composition Site}, and \emph{Nonisometric Transport and Amplified Reconstruction}.
\bibitem{moving} Komorebi P. \emph{Komorebi Space: Controlled Motion of Exact Realization Spaces}, focused version 0.32.0. Historical source archive and replay evidence, 30 September 2026. Included unchanged under \code{evidence/archives/}.
\bibitem{hilbert-extra} Komorebi P. \emph{HilbertHolonomyRepair.lean}. Separately recovered source implementing closed-range amplified holonomy reconstruction. Included under \code{evidence/selected/}. Its original separate provenance is retained. A byte-identical canonical copy was subsequently compiled with its dependencies in the accepted version 1.2.0 integrated replay.
\bibitem{clason-inverse} Christian Clason. \emph{Regularization of Inverse Problems}. Graduate lecture notes, arXiv:2001.00617v2, 2021. Sections 2, 4, and 6 provide standard background on weak convergence, parameter choice, and Tikhonov regularization. \url{https://arxiv.org/abs/2001.00617v2}. Consulted 12 October 2026.
\end{thebibliography}


\end{document}
